\chapter{Finite-time extinction}
\label{chap:finite-time-extinction}
\label{energy}



Our purpose in this chapter is to prove the following finite-time
extinction theorem for certain Ricci flows with surgery which, as we
shall show below, when combined with the theorem on the existence of
Ricci flows with surgery defined for all $t\in [0,\infty)$
(Theorem~\ref{MAIN}), immediately yields Theorem~\ref{Theorem1},
thus completing the proof of the Poincar\'e Conjecture and the
$3$-dimensional space-form conjecture.






\section{The result}




\begin{theorem}[Finite-time extinction for Ricci flow with surgery]
\label{thm:finite-time-extinction-main}
\label{extinct}
\uses{MAIN}
Let $(M,g(0))$ be a compact, connected normalized Riemannian $3$-manifold.
Suppose that the fundamental group of $M$ is a free product of finite groups
and infinite cyclic groups. Then $M$ contains no $\Ar P^2$ with trivial normal
bundle.  Let $({\mathcal M},G)$ be the Ricci flow with surgery defined for all
$t\in [0,\infty)$ with $(M,g(0))$ as initial conditions given by
Theorem~\ref{MAIN}. This Ricci flow with surgery becomes extinct after a finite
time in the sense that the time-slices $M_T$ of ${\mathcal M}$ are empty for
all $T$ sufficiently large.
\end{theorem}


Let us quickly show how this theorem implies our main result
Theorem~\ref{Theorem1}.

\begin{proof}(of Theorem~\ref{Theorem1} assuming
Theorem~\ref{thm:finite-time-extinction-main}).
\uses{thm:finite-time-extinction-main,MAIN,Tgood0good}
Fix a normalized metric $g(0)$ on $M$, and let
$({\mathcal M},G)$ be the Ricci flow with surgery defined for all $t\in
[0,\infty)$ produced by Theorem~\ref{MAIN} with initial conditions $(M,g(0))$.
According to Theorem~\ref{thm:finite-time-extinction-main} there is $T>0$ for which the time-slice
$M_T$ is empty. By Corollary~\ref{Tgood0good}, if there is $T$ for which $M_T$
is empty, then for any $T'<T$ the manifold $M_{T'}$ is a disjoint union of
connected sums of $3$-dimensional spherical space forms and $2$-sphere bundles
over $S^1$. Thus, the manifold $M=M_0$ is a connected sum of $3$-dimensional
space-forms and $2$-sphere bundles over $S^1$. This proves
Theorem~\ref{Theorem1}. In particular, if $M$ is simply connected, then $M$ is
diffeomorphic to $S^3$, which is the statement of the Poincar\'e Conjecture.
Similarly, if $\pi_1(M)$ is finite then $M$ is diffeomorphic to a connected sum
of a $3$-dimensional spherical space-form and $3$-spheres, and hence $M$ is
diffeomorphic to a $3$-dimensional spherical space-form.
\end{proof}


The rest of this chapter is devoted to the proof of
Theorem~\ref{thm:finite-time-extinction-main} which will then complete the proof of
Theorem~\ref{Theorem1}.

\subsection{History of this approach}

The basic idea for proving finite-time extinction is to use a min-max function
based on the area (or the closely related energy) of $2$-spheres or $2$-disks
in the manifold. The critical points of the energy functional are harmonic maps
and they play a central role in the proof. For a basic reference on harmonic
maps see \cite{SacksUhlenbeck}, \cite{SchoenYau}, and \cite{Jost}. Let us
sketch the argument. For a compact Riemannian manifold $(M,g)$ every non-zero
element $\boldsymbol{\beta}\in \pi_2(M)$ has associated with it an area, denoted
$W_2(\boldsymbol{\beta},g)$, which is the infimum over all maps $S^2\to M$ in the free
homotopy class of $\boldsymbol{\beta}$ of the energy of the map. We find it convenient to
set $W_2(g)$ equal to the minimum over all non-zero homotopy classes $\boldsymbol{\beta}$
of $W_2(\boldsymbol{\beta},g)$. In the case of a Ricci flow $g(t)$ there is an estimate
(from above) for the forward difference quotient of $W_2(g(t))$ with respect to
$t$. This estimate shows that after a finite time $W_2(g(t))$ must go negative.
This  is absurd since $W_2(g(t))$ is  always non-negative. This means that the
Ricci flow cannot exist for all forward time. In fact, using the
distance-decreasing property for surgery in Proposition~\ref{surgerydist} we
see that, even in a Ricci flow with surgery, the same forward difference
quotient estimate holds for as long as  $\pi_2$ continues to be non-trivial,
i.e., is not killed by the surgery. The forward difference quotient estimate
means that eventually all of $\pi_2$ is killed in a Ricci flow with surgery and
we arrive at a time $T$ for which every component of the $T$ time-slice, $M_T$,
has trivial $\pi_2$. This result holds for all Ricci flows with surgery
satisfying the conclusion of Theorem~\ref{MAIN}.

Now we fix $T_0$ so that every component of $M_{T_0}$ has trivial
$\pi_2$. It follows easily from the description of surgery that the
same statement holds for all $T\ge T_0$. We wish to show that, under
the group-theoretic hypothesis of Theorem~\ref{thm:finite-time-extinction-main}, at some
later time $T'>T_0$ the time-slice $M_{T'}$ is empty. The argument
here is similar in spirit. There are two approaches. The first
approach is due to Perelman \cite{Perelman3}. Here, one represents a
non-trivial element in $\pi_3(M_{T_0},x_0)$ by  a non-trivial
element in $\pi_2(\Lambda M,*)$, where $\Lambda M$ is the free loop
space on $M$ and $*$ is the trivial loop at $x_0$. For any compact
family $\Gamma$ of homotopically trivial loops in $M$ we consider
the areas of minimal spanning disks for each of the loops in the
family and set $W(\Gamma)$ equal to the maximal area of these
minimal spanning disks. For a given element in $\gamma\in
\pi_2(\Lambda M)$ we set $W(\gamma)$ equal to
 the infimum over all representative $2$-sphere families $\Gamma$
for $\gamma$  of $W(\Gamma)$. Under Ricci flow, the forward
difference quotient of this invariant satisfies an inequality and
the distance-decreasing property of surgery
(Proposition~\ref{surgerydist}) says that the inequality remains
valid for Ricci flow with surgery. The inequality implies that the
value $W(\gamma)$ goes negative in finite time, which is impossible.

The other approach, by Colding-Minicozzi \cite{ColdingMinicozzi}, is
to represent a non-trivial element in $\pi_3(M_T)$ as a non-trivial
element in $\pi_1(\operatorname{Maps}(S^2,M))$, and associate to such an
element the infimum over all representative families of the maximal
energy of the $2$-spheres in the family. Again, one shows that under
Ricci flow the forward difference quotient of this minimax satisfies
an inequality that implies that it goes negative in finite time. As
before, the distance-decreasing property of surgery
(Proposition~\ref{surgerydist}) implies that this inequality is
valid for Ricci flows with surgery. This tells us that the manifold
must completely disappear in finite time.




Our first reaction was that, of the two approaches, the one
considered by Colding-Minicozzi was preferable since it seemed more
natural and it had the advantage of avoiding the boundary issues
that occupy most of Perelman's analysis in \cite{Perelman3}.
 In the Colding-Minicozzi approach one must construct paths of $2$-spheres with the
 property that when the energy of the $2$-sphere is close to the
 maximum value along the path, then the $2$-sphere in question
 represents a point in the space $\operatorname{Maps}(S^2,M)$ that is
 close to a (usually) non-minimal critical point for the energy
 functional on this space. Such paths are needed in order to
 establish the forward difference quotient result alluded to above.
In Perelman's approach, one deals only with area-minimizing disks so
that one avoids having to deal with non-minimal critical points at
the expense of dealing with the technical issues related to the
boundary. Since the latter are one-dimensional in nature, they are
much easier to handle. In the end we decided to follow Perelman's
approach, and that is the one we present here.  In \cite{Perelman3}
there were two points that we felt required quite a bit of argument
beyond what Perelman presented. In \S 2.2 on page 4 of
\cite{Perelman3}, Perelman asserts that there is a local, pointwise
curvature estimate that can be obtained by adapting arguments in the
literature; see Lemmas~\ref{lem:curve-shrinking-shi-estimate} and~\ref{lem:curve-shrinking-shi-estimate-first-inequality} for the precise
statement. To implement this adaption required further non-trivial
arguments. We present these arguments in Section~\ref{18.6}. In \S
2.5 on page 5 of \cite{Perelman3} Perelman asserts that an
elementary argument establishes a lower bound on the length of a
boundary curve of a minimal annulus; see
Proposition~\ref{prop:annulus-boundary-length-estimate} for a precise statement. While the
statement seems intuitively clear, we found the argument, while
elementary, was quite intricate. We present this argument in
Section~\ref{18.5}.

The first use of these types of ideas to show that geometric objects
must disappear in finite time under Ricci flow is due to Hamilton
\cite{HamiltonNSRF3M}. He was considering a situation where a
time-slice $(M,g(t_0))$ of a $3$-dimensional Ricci flow had
submanifolds on which the metric was close to (a truncated version)
of a hyperbolic metric of finite volume. He wished to show that
eventually the boundary tori of the truncation were incompressible
in the $3$-manifold. If not, then there would be an immersed minimal
disk in $M$ whose boundary was a non-trivial loop on the torus. He
represented this relative homotopy class by a minimal energy disk in
$(M,g(t_0))$ and proved the same sort of forward difference quotient
estimate for the area of the minimal disk in the relative homotopy
class.  The same contradiction -- the forward difference quotient
implies that after a finite time the area would go negative if the
disk continued to exist --- implies that after a finite amount of
time this compressing disk must disappear. Using this he showed that
for sufficiently large time all the boundary tori of almost
hyperbolic submanifolds in $(M,g(t))$ were incompressible.

In the next section we deal with $\pi_2$ and, using $W_2$, we show
that given a Ricci flow with surgery as in Theorem~\ref{MAIN} there
is $T_1<\infty$ such that for all $T\ge T_1$ every connected
component of $M_T$ has trivial $\pi_2$. Then in the section after
that, by analyzing $W_3$, we show that, under the group-theoretic
hypothesis of Theorem~\ref{thm:finite-time-extinction-main}, there is a $T_2<\infty$ such
that $M_T=\emptyset$ for all $T\ge T_2$. In both these arguments we
need the same type of results -- a forward difference inequality for
the energy function;  the statement that away from surgery times
this function is continuous; and lastly, the statement that the
value of the energy function at a surgery time is at most the liminf
of its values at a sequence of times approaching the surgery time
from below.

\subsection{Existence of the Ricci flow with surgery}

Let $(M,g(0))$ be as in the statement of Theorem~\ref{thm:finite-time-extinction-main}, so that $M$ is
a compact, connected $3$-manifold whose fundamental group is a free product of
finite groups and infinite cyclic groups. By scaling $g(0)$ by a sufficiently
large constant, we can assume that $g(0)$ is normalized. Let us show that such
a manifold cannot contain an embedded $\RR P^2$ with trivial normal bundle.
First note that since $\RR P^2$ has Euler characteristic one, it is not the
boundary of a compact $3$-manifold. Hence, an $\RR P^2$ embedded with trivial
normal bundle does not separate the connected component of $M$ containing it.
Also, any non-trivial loop in $\Ar P^2$ has non-trivial normal bundle in $M$ so
that inclusion of $\RR P^2$ into $M$ induces an injection on fundamental
groups. Under the fundamental group hypotheses, $M$ decomposes as a connected
sum of $3$-manifolds with finite fundamental groups and $2$-sphere bundles over
$S^1$, see \cite{Hempel}.
 Given
an $\RR P^2$ with trivial normal bundle embedded in a connected sum,
it can be replaced by one contained in one of the connected factors.
[Proof: Let $\Sigma=\Sigma_1\cup\cdots\cup\Sigma_n$ be the spheres
giving the connected sum decomposition of $M$.  Deform the $\Ar P^2$
until it is transverse to $\Sigma$ and let $\gamma$ be a circle of
intersection of $\Ar P^2$ with one of the $\Sigma_i$ that is
innermost on $\Sigma_i$ in the sense that $\gamma$ bounds a disk $D$
in $\Sigma_i$ disjoint from all other components of intersection of
$\Sigma_i$ and $\Ar P^2$. The loop $\gamma$ also bounds a disk $D'$
in $\Ar P^2$. Replace $D'$ by $D$ and push $D$ slightly off to the
correct side of $\Sigma_i$. This will produce a new embedded $\Ar
P^2$ with trivial normal bundle in $M$ and at least one fewer
component of intersection with $\Sigma$. Continue inductively until
all components of intersection with $\Sigma$ are removed.]

Now suppose that we have an $\Ar P^2$ with trivial normal bundle
embedded disjointly from $\Sigma$, and hence embedded in one of the
prime factors of $M$.
 Since it
does not separate this factor, by the Mayer-Vietoris sequence (see p. 149 of
\cite{Hatcher}) the first homology of the factor in question maps onto $\Zee$
and hence the factor in question has infinite fundamental group. But this group
also contains the cyclic subgroup of order two, namely the image of $\pi_1(\RR
P^2)$ under the map induced by the inclusion. Thus, the fundamental group of
this prime factor is not finite and is not infinite cyclic. This is a
contradiction. (We have chosen to give a topological argument for this result.
There is also an argument using the theory of groups acting on trees which is
more elementary in the sense that it uses no $3$-manifold topology. Since it is
a more complicated, and to us, a less illuminating argument, we decided to
present the topological argument.)

Thus, by Theorem~\ref{MAIN},  for any compact $3$-manifold $M$ whose
fundamental group is a free product of finite groups and infinite cyclic groups
and for any normalized metric $g(0)$ on $M$ there is a Ricci flow with surgery
$({\mathcal M},G)$ defined for all time $t\in [0,\infty)$ satisfying the
conclusion of Theorem~\ref{MAIN} with $(M,g(0))$ as the initial conditions.

\begin{definition}[Path of components]
\label{def:path-of-components}
Let $I$ be an interval (which is allowed to be open or closed at
each end and finite or infinite at each end). By a {\em path of
components} of a Ricci flow with surgery $({\mathcal M},G)$ defined
for all $t\in I$ we mean a connected, open subset ${\mathcal
X}\subset {\bf t}^{-1}(I)$ with the property that for every $t\in I$
the intersection ${\mathcal X}(t)$ of ${\mathcal X}$ with each
time-slice $M_t$ is a connected component of $M_t$.
\end{definition}

Let ${\mathcal X}$ be a path of components in a Ricci flow with surgery
$({\mathcal M},G)$, a path defined for all $t\in I$. Let $I'$ be a subinterval
of $I$ with the property that no point of $I'$ except possibly its initial
point is a surgery time.  Then the intersection of ${\mathcal X}$ with ${\bf
t}^{-1}(I')$ is  the Ricci flow on the time interval $I'$ applied to ${\mathcal
X}(t)$ for any $t\in I'$.  Thus, for such intervals $I'$ the intersection,
${\mathcal X}(I')$, of ${\mathcal X}$ with ${\bf t}^{-1}(I')$ is determined by
the time-slice ${\mathcal X}(t)$ for any $t\in I'$. That is no longer
necessarily the case if some point of $I'$ besides its initial point is a
surgery time. Let $t\in I$ be a surgery time, distinct from the initial point
of $I$ (if there is one), and let $I'\subset I$ be an interval of the form
$[t',t)$ for some $t'<t$ sufficiently close to $t$ so that there are no surgery
times in $[t',t)$. Then, as we have just seen, ${\mathcal X}(I')$ is a Ricci
flow on the connected manifold ${\mathcal X}(t')$. There are several possible
outcomes of the result of surgery at time $t$ on this manifold. One possibility
is that the surgery leaves this connected component unchanged (affecting only
other connected components). In this case, there is no choice for ${\mathcal
X}(t)$: it is the continuation to time $t$ of the Ricci flow on ${\mathcal
X}(t')$. Another possibility is that ${\mathcal X}(t')$ is completely removed
by the surgery at time $t$. In this case the manifold ${\mathcal X}$ cannot be
continued to time $t$, contradicting the fact that the path of components
${\mathcal X}$ exists for all $t\in I$. The last possibility is that at time
$t$ surgery is done on ${\mathcal X}(t')$ using one or more $2$-spheres
contained in ${\mathcal X}(t')$. In this case the result of surgery on
${\mathcal X}(t')$ results in one or several connected components and
${\mathcal X}(t)$ can be any one of these.


\section{Disappearance of components with non-trivial $\pi_2$}

Let $({\mathcal M},G)$ be a Ricci flow with surgery satisfying the
conclusions of Theorem~\ref{MAIN}. We make no assumptions about the
fundamental group of the initial manifold $M_0$. In this section we
shall show that at some finite time $T_1$ every connected component
of $M_{T_1}$ has trivial $\pi_2$ and that this condition persists
for all times $T\ge T_1$. There are two steps in this argument.
First, we show that there is a finite time $T_0$ such that after
time $T_0$ every $2$-sphere surgery is performed along a
homotopically trivial $2$-sphere. (Using Kneser's theorem on
finiteness of topologically non-trivial families of $2$-spheres, one
can actually show by the same argument that after some finite time
all $2$-sphere surgeries are done along $2$-spheres that bound
$3$-balls. But in fact, Kneser's theorem will follow from what we do
here.)

After time $T_0$ the number of components with non-trivial $\pi_2$
is a weakly monotone decreasing function of time. The reason is the
following. Consider a path of components ${\mathcal X}$ defined for
$t\in [T_0,t']$ with the property that each time-slice ${\mathcal
X}(t)$ has non-trivial $\pi_2$. Using the fact that after time $T_0$
all the $2$-sphere surgeries are along homotopically trivial
$2$-spheres,  one shows easily that ${\mathcal X}$ is determined by
its initial time-slice ${\mathcal X}(T_0)$. Also, it is easy to see
that if there is a component of $M_{t'}$ with non-trivial $\pi_2$,
then it is the final time-slice of some path of components defined
for $t\in [T_0,t']$ with every time-slice of this path having
non-trivial $\pi_2$. This then produces an injection from the set of
connected components of $M_{t}$ with non-trivial $\pi_2$ into the
set of connected components of $M_{T_0}$ with non-trivial $\pi_2$.



The second step in the argument is to fix a path ${\mathcal X}(t)$,
$T_0\le t\le t'$, of connected components with non-trivial $\pi_2$
and to consider the function $W_2=W_2^{\mathcal X}$ that assigns to
each $t\in [T_0,t']$ the minimal area of a homotopically non-trivial
$2$-sphere mapping into ${\mathcal X}(t)$. We show that this
function is continuous except at the surgery times. Furthermore, we
show that if $t$ is a surgery time, then $W_2(t)\le \mathrm{liminf}_{t'\rightarrow t^-}W_2(t)$. Lastly, we show that at any
point $t\ge T_0$ we have
$$\frac{dW_2}{dt}(t)\le -4\pi-\frac{1}{2}R_{\min}(t)W_2(t),$$
in the sense of forward difference quotients. It follows easily from
the bound $R_{\min}(t)\ge -6/(4t+1)$ that there is $T_1({\mathcal
X})$  such that $W_2$ with these three properties cannot be
non-negative for all $t\in [T_0,T_1({\mathcal X})]$ and hence
$t'<T_1$. Since there are only finitely many components with
non-trivial $\pi_2$ at time $T_0$ it follows that there is
$T_1<\infty$ such that every component of $M_T$ has trivial $\pi_2$
for every $T\ge T_1$.


\subsection{A group-theory lemma}

To bound the number of homotopically non-trivial $2$-spheres in a
compact $3$-manifold we need the following group theory lemma.

\begin{lemma}[Rank bound for free product decompositions]
\label{lem:free-product-rank-bound}
\label{grouptheory}
Suppose that $G$ is a finitely generated group,
say generated by $k$ elements. Let $G=G_1*\cdots*G_\ell$ be a free
product decomposition of $G$ with non-trivial free factors, i.e.,
with $G_i\not=\{1\}$ for each $i=1,\ldots,\ell$. Then $\ell\le k$.
\end{lemma}


\begin{proof}
This is a consequence of Grushko's theorem \cite{Stallings}, which
says that given a map of a finitely generated free group $F$ onto
the  free product $G$, one can decompose the free group as a free
product of free groups $F=F_1*\cdots*F_\ell$ with $F_i$ mapping onto
$G_i$.
\end{proof}

\subsection{Homotopically non-trivial families of $2$-spheres}



\begin{definition}[Homotopically essential family of 2-spheres]
\label{def:homotopically-essential-family}
Let $X$ be a compact $3$-manifold (possibly disconnected). An
embedded $2$-sphere in $X$ is said to be {\em homotopically
essential} if the inclusion of the $2$-sphere into $X$ is not
homotopic to a point map of the $2$-sphere to $X$. More generally,
let $F=\{\Sigma_1,\ldots,\Sigma_n\}$ be a family of disjointly
embedded $2$-spheres in $X$. We say that the family is {\em
homotopically essential} if
\begin{enumerate}
\item[(i)] each $2$-sphere in the family is homotopically essential, and
\item[(ii)] for any $1\le i<j\le n$, the inclusion of $\Sigma_i$ into $X$ is not
homotopic in $X$ to the inclusion of $\Sigma_j$ into $X$.
\end{enumerate}
\end{definition}

Notice that if $F=\{\Sigma_1,\ldots,\Sigma_n\}$ is a homotopically
essential family of disjointly  embedded $2$-spheres in $X$, then
any subset $F$ is also homotopically essential.

\begin{lemma}[Finiteness of homotopically essential families of 2-spheres]
\label{lem:homotopically-essential-family-finite}
Let $X$ be a compact $3$-manifold (possibly disconnected). Then
there is a finite upper bound to the number of spheres in any
homotopically essential family of disjointly embedded $2$-spheres.
\end{lemma}

\begin{proof}
\uses{lem:free-product-rank-bound}
Clearly, without loss of generality we can assume that $X$ is
connected. If $F$ is a homotopically essential family of $2$-spheres
in $X$, then by van Kampen's theorem, see p. 40 of \cite{Hatcher},
there is an induced graph of groups decomposition of $\pi_1(X)$ with
all the edge groups being trivial. Since the family is homotopically
essential, it follows that the group associated with each vertex of
order $1$ and each vertex of order $2$ is non-trivial group. The
rank of the first homology of the graph underlying the graph of
groups, denoted $k$, is bounded above by the rank of $H_1(X)$.
Furthermore, by the theory of graphs of groups there is a free
product decomposition of $\pi_1(X)$ with the free factors being the
vertex groups and then $k$ infinite cyclic factors. Denote by $V_i$
the number of vertices of order $i$ and by $E$ the number of edges
of the graph. The number $E$ is the number of $2$-spheres in the
family $F$.
 An elementary combinatorial argument shows that
$$2V_1+V_2\ge E+3(1-k).$$
Thus, we have a free product decomposition of $\pi_1(X)$ with at least
$E+3(1-k)$ non-trivial free factors. Since $k$ is bounded by the rank of
$H_1(X)$,
 applying
Lemma~\ref{lem:free-product-rank-bound} and using the fact that the fundamental group of a
compact manifold is finitely presented establishes the result.
\end{proof}

\subsection{Two-sphere surgeries are trivial after finite
time}



\begin{definition}[Homotopically essential surgery]
\label{def:homotopically-essential-surgery}
Let $({\mathcal M},G)$ be a Ricci flow with surgery. We say that a
surgery along a $2$-sphere $S_0(t)$ at time $t$ in $({\mathcal
M},G)$ is a {\em homotopically essential surgery} if, for every
$t'<t$ sufficiently close to $t$,  flowing  $S_0(t)$ backwards from
time $t$ to time $t'$ results in a homotopically essential
$2$-sphere $S_0(t')$ in $M_{t'}$.
\end{definition}


\begin{proposition}[Finitely many homotopically essential surgeries]
\label{prop:finitely-many-homotopically-essential-surgeries}
\label{homess}
Let $({\mathcal M},G)$ be a Ricci flow with surgery satisfying
Assumptions (1) -- (7) in Chapter~\ref{sect:surgery}. Then there can
be only finitely many homotopically essential surgeries along
$2$-spheres in $({\mathcal M},G)$.
\end{proposition}


\begin{proof}
 Associate to each compact $3$-manifold $X$
 the invariant $s(X)$ which is the maximal number of spheres in
 any homotopically essential family of embedded $2$-spheres in $X$.
 The main step in establishing the corollary is the following:

\begin{claim}[Monotonicity of the essential-sphere invariant $s$]
\label{claim:invariant-s-monotone}
Let $({\mathcal M},G)$ be a Ricci flow with surgery and for each $t$ set
$s(t)=s(M_t)$. If $t'<t$ then $s(t')\ge s(t)$. If  we do surgery at time $t$
along at least one homotopically essential $2$-sphere, then $s(t)<s(t')$ for
any $t'<t$.
\end{claim}

\begin{proof}
\uses{surgerytoptype}
Clearly, for any $t_0$ we have $s(t)=s(t_0)$ for $t\ge t_0$
sufficiently close to $t_0$. Also, if $t$ is not a surgery time,
then $s(t)=s(t')$ for all $t'<t$ and sufficiently close to $t$.
According to Proposition~\ref{surgerytoptype}, if $t$ is a surgery
time then for $t'<t$ but sufficiently close to it, the manifold
$M_t$ is obtained from $M_{t'}$  by doing surgery on a finite number
of $2$-spheres and removing certain components of the result. We
divide the operations into three types: (i) surgery along
homotopically trivial $2$-spheres in $M_{t'}$, (ii) surgery along
homotopically non-trivial $2$-spheres in $M_{t'}$, (iii) removal of
components. Clearly, the first operation does not change the
invariant $s$ since it simply creates a manifold that is the
disjoint union of a manifold homotopy equivalent to the original
manifold with a collection of homotopy $3$-spheres. Removal of
components will not increase the invariant. The last operation to
consider is surgery along a homotopically non-trivial $2$-sphere.
Let $F_t$ be a homotopically essential family of disjointly embedded
$2$-spheres in $M_t$. This family of $2$-spheres in $M_t$ can be
deformed to miss the $3$-disks (the surgery caps) in $M_t$ that we
sewed in doing the surgery at time $t$ along a homotopically
non-trivial $2$-sphere. After deforming the spheres in the family
$F_t$ away from the surgery caps, they produce a disjoint family
$F'_{t'}$ of $2$-spheres in the manifold $M_{t'}$, for $t'<t$ but
$t'$ sufficiently close to $t$. Each $2$-sphere in $F'_{t'}$ is
disjoint from the homotopically essential $2$-sphere $S_0$ along
which we do surgery at time $t$. Let $F_{t'}$ be the family
$F'_{t'}\cup \{S_0\}$. We claim that $F_{t'}$ is a homotopically
essential family in $M_{t'}$.

First, suppose that one of the spheres $\Sigma$ in $F_{t'}$ is
homotopically trivial in $M_{t'}$. Of course, we are in the case
when the surgery $2$-sphere is homotopically essential, so $\Sigma$
is not $S_0$ and hence is the image of one of the $2$-spheres in
$F_t$. Since $\Sigma$ is homotopically trivial, it is the boundary
of a homotopy $3$-ball $B$ in $M_{t'}$. If $B$ is disjoint from the
surgery $2$-sphere $S_0$, then it exists in $M_t$ and hence $\Sigma$
is homotopically trivial in $M_t$, which is not possible from the
assumption about the family $F_t$. If $B$ meets the surgery
$2$-sphere $S_0$, then since the spheres in the family $F_{t'}$ are
disjoint, it follows that $B$ contains the surgery $2$-sphere $S_0$.
This is not possible since in this case $S_0$ would be homotopically
trivial in $M_{t'}$, contrary to assumption.

We also claim that no distinct members of $F_{t'}$ are homotopic.
For suppose that two of the members $\Sigma$ and $\Sigma'$ are
homotopic. It cannot be the case that one of $\Sigma$ or $\Sigma'$
is the surgery $2$-sphere $S_0$ since, in that case, the other one
would be homotopically trivial after surgery, i.e., in $M_t$. The
$2$-spheres $\Sigma$ and $\Sigma'$ are the boundary components of a
submanifold  $A$ in $M_{t'}$ homotopy equivalent to $S^2\times I$.
If $A$ is disjoint from the surgery $2$-sphere $S_0$, then $A$
exists in $M_t$ and $\Sigma$ and $\Sigma'$ are homotopic in $M_{t}$,
contrary to assumption. Otherwise, the surgery sphere $S_0$ must be
contained in $A$. Every $2$-sphere in $A$ is either homotopically
trivial in $A$ or is homotopic in $A$ to either boundary component.
If $S_0$ is homotopically trivial in $A$, then it would be
homotopically trivial in $M_{t'}$ and this contradicts our
assumption. If $S_0$ is homotopic in $A$ to each of $\Sigma$ and
$\Sigma'$, then each of $\Sigma$ and $\Sigma'$ is homotopically
trivial in $M_{t'}$, contrary to assumption. This shows that the
family $F_{t'}$ is homotopically essential. It follows immediately
that doing surgery on a homotopically non-trivial $2$-sphere
strictly decreases the invariant $s$.
\end{proof}


Proposition~\ref{prop:finitely-many-homotopically-essential-surgeries} is immediate from this claim and the
previous lemma.
\end{proof}


\subsection{For all $T$ sufficiently large $\pi_2(M_T)=0$}


We have just established that given any Ricci flow with surgery $({\mathcal
M},G)$ satisfying the conclusion of Theorem~\ref{MAIN} there is $T_0<\infty$,
depending  on $({\mathcal M},G)$, such that all surgeries after time $T_0$
either are along homotopically trivial $2$-spheres or remove entire components
of the manifold. Suppose that $M_{T_0}$ has a component ${\mathcal X}(T_0)$
with non-trivial $\pi_2$, and suppose that we have a path of components
${\mathcal X}(t)$ defined for $t\in [T_0,T)$ with the property that each
time-slice has non-trivial $\pi_2$. If $T$ is not a surgery time, then there is
a unique extension of ${\mathcal X}$ to a path of components with non-trivial
$\pi_2$ defined until the first surgery time after $T$. Suppose that $T$ is a
surgery time and let us consider the effect of surgery at time $T$ on
${\mathcal X}(t)$ for $t<T$ but close to it. Since no surgery after time $T_0$
is done on a homotopically essential $2$-sphere there are three possibilities:
(i) ${\mathcal X}(t)$ is untouched by the surgery, (ii) surgery is performed on
one or more homotopically trivial $2$-spheres in ${\mathcal X}(t)$, or (iii)
the component ${\mathcal X}(t)$ is completely removed by the surgery. In the
second case, the result of the surgery on ${\mathcal X}(t)$ is a disjoint union
of components one of which is homotopy equivalent to ${\mathcal X}(t)$, and
hence has non-trivial $\pi_2$, and all others are homotopy $3$-spheres. This
implies that there is a unique extension of the path of components preserving
the condition that every time-slice has non-trivial $\pi_2$, unless the
component ${\mathcal X}(t)$ is removed by surgery at time $T$, in which case
there is no extension of the path of components to time $T$. Thus, there is a
unique maximal such path of components starting at ${\mathcal X}(T_0)$ with the
property that every time-slice has non-trivial $\pi_2$. There are two
possibilities for the interval of definition of this maximal path of components
with non-trivial $\pi_2$. It can be $[T_0,\infty)$ or it is of the form
$[T_0,T)$, where the surgery at time $T$ removes the component ${\mathcal
X}(t)$ for $t<T$ sufficiently close to it.







\begin{proposition}[Eventual triviality of $\pi_2$]
\label{prop:eventual-trivial-pi2}
\label{pi2}
Let $({\mathcal M},G)$ be a Ricci flow with surgery satisfying the conclusion
of Theorem~\ref{MAIN}. Then there is some time $T_1<\infty$ such that every
component of $M_T$ for any $T\ge T_1$ has trivial $\pi_2$. For every $T\ge
T_1$, each component of $M_T$ either has finite fundamental group, and hence
has a homotopy $3$-sphere as universal covering, or has  contractible universal
covering.
\end{proposition}

If $M$ is a connected $3$-manifold with $\pi_2(M)=0$, then the
universal covering, $\widetilde M$, of $M$ is a $2$-connected
$3$-manifold. The covering $\widetilde M$ is compact if and only if
$\pi_1(M)$ is finite. In this case $\widetilde M$ is a homotopy
$3$-sphere. If $\widetilde M$ is non-compact then $H_3(\widetilde
M)=0$, so that all its homology groups and hence, by the Hurewicz
theorem, all its homotopy groups vanish. It follows from the
Whitehead theorem  that $\widetilde M$ is contractible in this case.
This proves the last assertion in the proposition modulo the first
assertion.

The proof of the first assertion of this proposition occupies the
rest of this subsection. By the above discussion we see that the
proposition holds unless there is a path of components ${\mathcal
X}$ defined for all $t\in [T_0,\infty)$ with the property that every
time-slice has non-trivial $\pi_2$. We must rule out this
possibility. To achieve this we introduce the area functional.

\begin{lemma}[Minimal energy gap for homotopically non-trivial 2-spheres]
\label{lem:minimal-sphere-energy-gap}
Let $X$ be a compact Riemannian manifold with $\pi_2(X)\not=0$. Then
there is a positive number $e_0=e_0(X)$ with the following two
properties:
\begin{enumerate} \item Any map $f\colon S^2\to X$ with area less
than $e_0$ is homotopic to a point map. \item There is a minimal
$2$-sphere $f\colon S^2\to X$, which is a branched immersion, with
the property that the area of $f(S^2)=e_0$ and with the property
that $f$ is not homotopic to a point map.
\end{enumerate}
\end{lemma}


\begin{proof}
The first statement is Theorem 3.3 in \cite{SacksUhlenbeck}. As for
the second, following Sacks-Uhlenbeck, for any $\alpha>1$ we
consider the perturbed energy $E_\alpha$ given by
$$E_\alpha(s)=\int_{S^2}\left(1+|ds|^2\right)^\alpha da.$$
 According to
\cite{SacksUhlenbeck} this energy functional is Palais-Smale on the
space of $H^{1,2\alpha}$ maps and has an absolute minimum among
homotopically non-trivial maps, realized by a map $s_\alpha\colon
S^2\to X$. We consider a decreasing sequence of $\alpha$ tending to
$1$ and the minimizers $s_\alpha$ among homotopically non-trivial
maps. According to \cite{SacksUhlenbeck}, after passing to a
subsequence, there is a weak limit which is a strong limit on the
complement of a finite set of points in $S^2$. This limit extends to
a harmonic map of $S^2\to M$, and its energy is less than or equal
to the limit of the $\alpha$-energies of $s_\alpha$. If the result
is homotopically non-trivial then it realizes a minimum value of the
usual energy among all homotopically non-trivial maps, for were
there a homotopically non-trivial map of smaller energy, it would
have smaller $E_\alpha$ energy than $s_\alpha$ for all $\alpha$
sufficiently close to $1$. Of course if the limit is a strong limit,
then the map is homotopically non-trivial, and the proof is
complete.

We must examine the case when the limit is truly a weak limit. Let
$s_n$ be a sequence as above with a weak limit $s$. If the limit is
truly a weak limit, then there is bubbling. Let $x\in S^2$ be a
point where the limit $s$  is not a strong limit. Then according to
\cite{SacksUhlenbeck} pre-composing with a sequence of conformal
dilations $\rho_n$ centered at this point leads to a sequence of
maps $s_n'$ converging uniformly on compact subsets of $\Ar ^2$ to a
non-constant harmonic map $s'$ that extends over the one-point
compactification $S^2$. The energy of this limiting map $s'$ is at
most the limit of the $\alpha$-energies of the $s_\alpha$. If $s'$
is homotopically non-trivial, then, arguing as before, we see that
it realizes the minimum energy among all homotopically non-trivial
maps, and once again we have completed the proof. We rule out the
possibility that $s'$ is homotopically trivial. Let $\alpha$ be the
area, or equivalently the energy, of $s'$. Let $D\subset \Ar^2$ be a
disk centered at the origin which contains three-quarters of the
energy of $s'$ (or equivalently three-quarters of the area of $s'$),
and let $D'$ be the complementary disk to $D$ in $S^2$.  For all $n$
sufficiently large the area of $s_n'|D$ minus the area of $s_n'|D'$
is at least $\alpha/3$. The restrictions of $s_n'$ on $\partial D$
are converging smoothly to $s'|\partial D'$. Let $D_n\subset S^2$ be
$\rho_n^{-1}(D)$. Then the area of $s_n|D_n$ equals the area of
$s_n'|D$ and hence is at least the area of $s'|D'$ plus $\alpha/4$
for all $n$ sufficiently large. Also, as $n$ tends to infinity the
image $s_n(D_n)$ converges smoothly, after reparameterization, to
$s'(\partial D)$. Thus, for all $n$ large, we can connect
$s_n(\partial D_n)$ to $s'(\partial D')$ by an annulus $A_n$
contained in a small neighborhood of $s'(\partial D')$ and whose
area tends to $0$ as $n$ goes to infinity.  For all $n$ sufficiently
large, the resulting $2$-sphere $\Sigma_n$ made out of
$s_n|(S^2\setminus D_n)\cup A_n\cup S'(D')$ is homotopic to $s(S^2)$
since $s'$ is homotopically trivial. Also, for all $n$ sufficiently
large, the area of $\Sigma_n$ is less than the area of $s_n$ minus
$\alpha/5$. Reparameterizing this $2$-sphere by a conformal map
leads to a homotopically non-trivial map of energy less than the
area of $s_n$ minus $\alpha/5$. Since as $n$ tends to infinity, the
limsup of the areas of the $s_n$ converge to at most $e_0$, for all
$n$ sufficiently large we have constructed a homotopically
non-trivial map of energy less than $e_0$, which contradicts the
fact that the minimal $\alpha$ energy for a homotopically
non-trivial map tends to $e_0$ as $\alpha$ tends to $1$.


Of course, any minimal energy map of $S^2$ into $M$ is conformal because there
is no non-trivial holomorphic quadratic differential on $S^2$. It follows that
such a map is a branched immersion.
\end{proof}


Now suppose that ${\mathcal X}$ is a path of components defined for
all $t\in [T_0,\infty)$ with $\pi_2({\mathcal X}(t))\not= 0$ for all
$t\in [T_0,\infty)$.
 For
each $t\ge T_0$ we define $W_2(t)$ to be $e_0({\mathcal X}(t))$, where $e_0$ is
the invariant  given in the previous lemma. Our assumption on ${\mathcal X}$
means that $W_2(t)$ is defined and positive for all $t\in [T_0,\infty)$.

\begin{lemma}[Forward difference quotient for $W_2$]
\label{lem:w2-forward-difference-quotient}
\label{W_2}
\uses{lem:minimal-sphere-energy-gap,lem:forward-difference-quotient}
$$\frac{d}{dt}W_2(t)\le -4\pi-\frac{1}{2}R_{\min}(t)W_2(t)$$
in the sense of forward difference quotients. If $t$ is not a
surgery time, then $W_2(t)$ is continuous at $t$, and if $t$ is a
surgery time, then $$W_2(t)\le \liminf_{t'\rightarrow
t^-}W_2(t').$$
\end{lemma}

Let us show how this lemma implies Proposition~\ref{prop:eventual-trivial-pi2}. Because
the curvature is pinched toward positive, we have
$$R_{\min}(t)\ge (-6)/(1+4t).$$
Let $w_2(t)$ be the function satisfying the differential equation
$$\frac{dw_2}{dt}=-4\pi+\frac{3w_2}{1+4t}$$
and $w_2(T_0)=W_2(T_0)$. Then by Lemma~\ref{lem:forward-difference-quotient} and
Lemma~\ref{lem:w2-forward-difference-quotient} we have $W_2(t)\le w_2(t)$ for all $t\ge T_0$. On
the other hand, we can integrate to find
$$w_2(t)=w_2(T_0)\frac{(4t+1)^{3/4}}{(4T_0+1)^{3/4}}+4\pi
(4T_0+1)^{1/4}(4t+1)^{3/4}-4\pi (4t+1).$$ Thus, for $t$ sufficiently
large, $w_2(t)<0$. This is a contradiction since $W_2(t)$ is always
positive, and $W_2(t)\le w_2(t)$.

This shows that to complete the proof of Proposition~\ref{prop:eventual-trivial-pi2} we
need only establish Lemma~\ref{lem:w2-forward-difference-quotient}.


\begin{proof}
\uses{claim:minimal-sphere-area-variation}
(of Lemma~\ref{lem:w2-forward-difference-quotient})
Let $f\colon S^2\to (X(t_0),g(t_0))$ be a minimal $2$-sphere.

\begin{claim}[Variation of area of a minimal 2-sphere]
\label{claim:minimal-sphere-area-variation}
\label{varareaform}
$$\frac{d\Area_{g(t)}(f(S^2))}{dt}(t_0)\le -4\pi -\frac{1}{2}
R_{\min}(g(t_0))\Area_{g(t_0)}f(S^2).$$
\end{claim}



\begin{proof}
 Recall that, for any immersed surface $f\colon S^2\to (M,g(t_0))$, we
have ([Ha])
 \begin{eqnarray}\label{ha1}
 \frac{d}{dt}\Area_{g(t)}(f(S^2))\bigl|_{t=t_0}\bigr. & = &
 \int_{S^2}\frac{1}{2}\operatorname{Tr}|_{S^2}\Bigl(\frac{\partial g}{\partial t}\Bigr)\Bigl|_{t=t_0}\Bigr.da \\
 & = &
 -\int_{S^2}(R-\Ric({\bf n},{\bf n}))da \nonumber
 \end{eqnarray}
 where $R$ denotes the scalar curvature of $M$, $\Ric$ is the
Ricci curvature of $M$, and ${\bf n}$ is the unit normal vector field of $f(S^2)$
in $M$. Now suppose that $f(S^2)$ is minimal. We can rewrite this as
 \begin{eqnarray}
 \frac{d}{dt}\Area_{g(t)}(f(S^2))\big |_{t=0}
 &= &
 -\int_{S^2}K_{S^2}da-\frac{1}{2}\int_{S^2}(|A|^2+R)da,  \label{ha2}
 \end{eqnarray}
where $K_{S^2}$ is the Gaussian curvature of $S^2$ and $A$ is the
second fundamental form of $f(S^2)$ in $M$. (Of course, since
$f(S^2)$ is minimal, the determinant of its second fundamental form
is $-|A|^2/2$.) Even if $f$ is only a branched minimal surface,
(\ref{ha2}) still holds when  the integral on the right is replaced
by the integral over the immersed part of $f(S^2)$. Then by the
Gauss-Bonnet theorem we have
 \begin{eqnarray}
\frac{d}{dt}\Area_{g(t)}(f(S^2))\bigl|_{t=t_0}\bigr.  & \leq &
-4\pi-\frac{1}{2}\Area_{g(t_0)}(S^2)\min_{x\in M}\{R_{g}(x,t_0)\}.  \label{ha3}
 \end{eqnarray}
\end{proof}




Since $f(S^2)$ is a homotopically non-trivial sphere in ${\mathcal X}(t)$ for
all $t$ sufficiently close to $t_0$ we see that $W_2(t)\le \Area_{g(t)}f(S^2)$. Since $\Area_{g(t)}f(S^2))$ is a smooth function of
$t$, the forward difference quotient statement in Lemma~\ref{lem:w2-forward-difference-quotient} follows
immediately from Claim~\ref{claim:minimal-sphere-area-variation}.

We turn now to continuity at non-surgery times. Fix $t'\ge T_0$
distinct from all surgery times. We show that  the function
$e_0(t')$ is continuous at $t'$. If $f\colon S^2\to {\mathcal
X}(t')$ is the minimal area, homotopically non-trivial sphere, then
the area of $f(S^2)$ with respect to a nearby metric $g(t)$ is close
to the area of $f(S^2)$ in the metric $g(t')$. Of course, the area
of $f(S^2)$ in the metric $g(t)$ is greater than or equal to
$W_2(t)$. This proves that $W_2(t)$ is  upper semi-continuous at
$t'$. Let us show that it is lower semi-continuous at  $t'$.

\begin{claim}[Area growth bound for a fixed 2-sphere map]
\label{claim:sphere-area-growth-bound}
Let $(M,g(t)),\ t_0\le t\le t_1$, be a Ricci flow on a compact
manifold. Suppose that $|\Ric_{g(t)}|\le D$ for all $t\in
[t_0,t_1]$ Let $f\colon S^2\to (M,g(t_0))$ be a $C^1$-map. Then
$$\Area_{g(t_1)}f(S^2)\le \Area_{g(t_0)}f(S^2)e^{4D(t_1-t_0)}.$$
\end{claim}

\begin{proof}
The rate of change of the area of $f(S^2)$ at time $t$ is
$$\int_{f(S^2)}\frac{\partial g}{\partial t}(t)da=-2\int_{f(S^2)}\operatorname{Tr}|_{TS^2}(\Ric_{g(t)})da\le
4D\Area_{g(t)}f(S^2) .$$ Integrating from $t_0$ to $t_1$ gives
the result.
\end{proof}

Now suppose that we have a family of times $t_n$ converging to a
time $t'$ that is not a surgery time. Let $f_n\colon S^2\to
{\mathcal X}(t_n)$ be the minimal area non-homotopically trivial
$2$-sphere in ${\mathcal X}(t_n)$, so that the area of $f_n(S^2)$ in
${\mathcal X}(t_n)$ is $e_0(t_n)$. Since $t'$ is not a surgery time,
for all $n$ sufficiently large we can view the maps $f_n$ as
homotopically non-trivial maps of $S^2$ into ${\mathcal X}(t')$. By
the above claim, for any $\delta>0$ for all $n$ sufficiently large,
the area of $f_n(S^2)$ with respect to the metric $g(t')$ is at most
the area of $f_n(S^2)$ plus $\delta$. This shows that for any
$\delta>0$ we have $W_2(t')\le W_2(t_n)+\delta$ for all $n$
sufficiently large, and hence $W_2(t')\le \mathrm{liminf}_{n\rightarrow\infty}W_2(t_n)$. This is the lower
semi-continuity.



The last thing to check is the behavior of $W_2$ near a surgery time
$t$. According to the description of the surgery process given in
Section~\ref{sect:process}, we write ${\mathcal X}(t)$ as the union
of a compact subset $C(t)$ and a finite number of surgery caps. For
every $t'<t$ sufficiently close to $t$ we have an embedding
$n_{t'}\colon C(t)\cong C(t')\subset {\mathcal X}(t')$ given by
flowing $C(t)$ backward under the flow to time $t'$. As
$t'\rightarrow t$ the maps $\eta_{t'}$ converge in the
$C^\infty$-topology to isometries, in the sense that the
$n_{t'}^*(g(t'))|_{C(t')}$ converge smoothly to $g(t)|_{C(t)}$.
Furthermore, since the $2$-spheres along which we do surgery are
homotopically trivial they separate $M_{t'}$. Thus, the maps
$n_{t'}^{-1}\colon C(t')\to C(t)$ extend to maps $\psi_{t'}\colon
{\mathcal X}(t')\to {\mathcal X}(t)$. The image under $\psi_{t'}$ of
${\mathcal X}(t')\setminus C(t')$ is contained in the union of the
surgery caps.  Clearly, since all the $2$-spheres on which we do
surgery at time $t$ are homotopically trivial, the maps $\psi_{t'}$
are homotopy equivalences. If follows from
Proposition~\ref{surgerydist} that for any $\eta>0$ for all $t'<t$
sufficiently close to $t$, the map $\psi_{t'}\colon {\mathcal
X}(t')\to {\mathcal X}(t)$ is a
 homotopy equivalence that is a $(1+\eta)$-Lipschitz map. Thus,
 given $\eta>0$ for all $t'<t$ sufficiently close to $t$,
 for any minimal $2$-sphere $f\colon S^2\to ({\mathcal X}(t'),g(t'))$ the area of
 $\psi_{t'}\circ f\colon S^2\to ({\mathcal X}(t),g(t))$ is at most $(1+\eta)^2$
 times the area of $f(S^2)$. Thus, given $\eta>0$ for all $t'<t$ sufficiently close to $t$
we see that $W_2(t)\le (1+\eta)^2W_2(t')$.
 Since this is true for every $\eta>0$, it follows that
 $$W_2(t)\le \liminf_{t'\rightarrow t^-}W_2(t').$$

This establishes all three statements in Proposition~\ref{prop:eventual-trivial-pi2} and
completes the proof of the proposition.
\end{proof}


As an immediate corollary of Proposition~\ref{prop:eventual-trivial-pi2}, we obtain the
sphere theorem for closed $3$-manifolds.

\begin{corollary}[Sphere theorem for closed 3-manifolds without embedded $\mathbb{RP}^2$]
\label{cor:sphere-theorem-no-rp2}
\label{spherethm}
Suppose that $M$ is a closed, connected $3$-manifold containing no
embedded $\Ar P^2$ with trivial normal bundle, and suppose that
$\pi_2(M)\not= 0$. Then either $M$ can be written as a connected sum
$M_1\# M_2$ where neither of the $M_i$ is homotopy equivalent to
$S^3$ or $M_1$ has a prime factor that is a $2$-sphere bundle over
$S^1$. In either case, $M$ contains an embedded $2$-sphere which is
homotopically non-trivial.
\end{corollary}

\begin{proof}
\uses{prop:eventual-trivial-pi2}
Let $M$ be as in the statement of the corollary. Let $g$ be a
normalized metric on $M$, and let $({\mathcal M},G)$ be the Ricci
flow with surgery defined for all time with $(M,g)$ as initial
conditions. According to Proposition~\ref{prop:eventual-trivial-pi2} there is $T<\infty$
such that every component of $M_T$ has trivial $\pi_2$. Thus, by the
analysis above, we see that there must be surgeries that kill
elements in $\pi_2$: either the removal of a component with
non-trivial $\pi_2$ or surgery along a homotopically non-trivial
$2$-sphere. We consider the first such surgery in $M$.
 The only
components with non-trivial $\pi_2$ that can be removed by surgery
are $S^2$-bundles over $S^1$ and $\Ar P^3\#\Ar P^3$. Since each of
these has homotopically non-trivially embedded $2$-spheres, if the
first surgery killing an element in $\pi_2$ is removal of such a
component, then, because all the earlier $2$-sphere surgeries are
along homotopically trivial $2$-spheres, the homotopically
non-trivial embedded $2$-sphere in this component deforms back to an
embedded, homotopically non-trivial $2$-sphere in $M$.  The other
possibility is that the first time an element in $\pi_2(M)$ is
killed it is by surgery along a homotopically non-trivial
$2$-sphere. Once again, using the fact that all previous surgeries
are along homotopically trivial $2$-spheres, deform this $2$-sphere
back to $M$ producing a homotopically non-trivial $2$-sphere in $M$.
\end{proof}

\begin{remark}[Components with non-trivial $\pi_2$ that disappear]
\label{rem:pi2-disappearing-components}
\uses{thm:finite-time-extinction-main}
Notice that it follows from the list of disappearing components that
the only ones with non-trivial $\pi_2$ are those based on the
geometry $S^2\times \Ar$; that is to say, $2$-sphere bundles over
$S^1$ and $\Ar P^3\#\Ar P^3$. Thus, once we have reached the level
$T_0$ after which all $2$-sphere surgeries are performed on
homotopically trivial $2$-spheres the only components that can have
non-trivial $\pi_2$ are components of these types. Thus, for example
if the original manifold has no $\Ar P^3$ prime factors and no
non-separating $2$-spheres, then when we reach time $T_0$ we have
done a connected sum decomposition into components each of which has
trivial $\pi_2$. Each of these components is either covered by a
contractible $3$-manifold or by a homotopy $3$-sphere, depending on
whether its fundamental group has infinite or finite order.
\end{remark}





\section{Extinction}

Now we assume that the Ricci flow with surgery $({\mathcal M},G)$ satisfies the
conclusion of Theorem~\ref{MAIN} and also has initial condition $M$ that is a
connected $3$-manifold whose fundamental group satisfies the hypothesis of
Theorem~\ref{thm:finite-time-extinction-main}. The argument showing that  components with non-trivial
$\pi_3$ disappear after a finite time is, in spirit, very similar to the
arguments above, though the technical details are more intricate in this case.

\subsection{Forward difference quotient for $\pi_3$}

Let $M$ be a compact, connected $3$-manifold. Fix a base point
$x_0\in M$. Denote by $\Lambda M$ the free loop space of $M$. By
this we mean the space of $C^1$-maps of $S^1$ to $M$ with the
$C^1$-topology. The components of $\Lambda M$ are the conjugacy
classes of elements in $\pi_1(M,x_0)$. The connected component of
the identity of $\Lambda M$ consists of all homotopically trivial
loops in $M$. Let $*$ be the trivial loop at $x_0$.

\begin{claim}[$\pi_2$ of the loop space identified with $\pi_3$]
\label{claim:loop-space-pi2-pi3-iso}
Suppose that $\pi_2(M,x_0)=0$. Then $\pi_2(\Lambda M,*)\cong
\pi_3(M,x_0)$ and $\pi_2(\Lambda M,*)$ is identified with the free
homotopy classes of maps of $S^2$ to the component of $\Lambda M$
consisting of homotopically trivial loops.
\end{claim}


\begin{proof}
An element in $\pi_2(\Lambda M,*)$ is represented by a map
$S^2\times S^1\to M$ that sends $\{\mathrm{pt}\}\times S^1$ to $x_0$.
Hence, this map factors through the quotient of $S^2\times S^1$
obtained by collapsing $\{\mathrm{pt}\}\times S^1$ to a point. The
resulting quotient space is homotopy equivalent to $S^2\vee S^3$,
and a map of this space into $M$ sending the wedge point to $x_0$
is, up to homotopy, the same as an element of $\pi_2(M,x_0)\oplus
\pi_3(M,x_0)$. But we are assuming that $\pi_2(M,x_0)=0$. The first
statement follows. For the second, notice that since $\pi_2(M,x_0)$
is trivial, $\pi_3(M,x_0)$ is identified with $H_3$ of the universal
covering $\widetilde M$ of $M$.  Hence, for any map of $S^2$ into
the component of $\Lambda M$ containing the trivial loops, the
resulting map $S^2\times S^1\to M$ lifts to  $\widetilde M$. The
corresponding element in $\pi_3(M,x_0)$ is the image of the
fundamental class of $S^2\times S^1$ in $H_3(\widetilde
M)=\pi_3(M)$.
\end{proof}


\begin{definition}[Spanning-disk energy functional $W(\xi)$]
\label{def:w-xi-spanning-disk-functional}
\label{Wdefn}
\uses{prop:eventual-trivial-pi2,lem:w2-forward-difference-quotient}
Fix a homotopically trivial loop $\gamma\in \Lambda M$. We set
$A(\gamma)$ equal to the infimum of the areas of any spanning disks
for $\gamma$, where by definition a spanning disk is a Lipschitz map
$D^2\to M$ whose boundary is, up to reparameterization, $\gamma$.
Notice that $A(\gamma)$ is a continuous function of $\gamma$ in
$\Lambda M$. Also, notice that $A(\gamma)$ is invariant under
reparameterization of the curve $\gamma$. Now suppose that
$\Gamma\colon S^2\to \Lambda M$ is given with the image consisting
of homotopically trivial loops.
 We define $W(\Gamma)$ to be equal to the maximum over all $c\in S^2$ of
$A(\Gamma(c))$. More generally, given a homotopy class $\xi\in \pi_2(\Lambda
M,*)$ we define $W(\xi)$ to be equal to the infimum over all (not necessarily
based) maps $\Gamma\colon S^2\to\Lambda M$ into the component of $\Lambda M$
consisting of homotopically trivial loops representing $\xi$ of $W(\Gamma)$.
\end{definition}




Now let us formulate the analogue of Proposition~\ref{prop:eventual-trivial-pi2} for
$\pi_3$. Suppose that ${\mathcal X}$ is a path of components of the
Ricci flow with surgery $({\mathcal M},G)$ defined for $t\in
[t_0,t_1]$. Suppose that $\pi_2({\mathcal X}(t_0),x_0)=0$ and that
$\pi_3({\mathcal X}(t_0),x_0)\not=0$. Then,  the same two conditions
hold for ${\mathcal X}(t)$ for each $t\in [t_0,t_1]$. The reason is
that at a surgery time $t$, since all the $2$-spheres in ${\mathcal
X}(t')$ ($t'<t$ but sufficiently close to $t$) along which we are
doing surgery are homotopically trivial, the result of surgery is a
disjoint union of connected components: one connected component is
homotopy equivalent to ${\mathcal X}(t')$ and all other connected
components are homotopy $3$-spheres. This means that either
${\mathcal X}(t)$ is homotopy equivalent to ${\mathcal X}(t')$ for
$t'<t$ or ${\mathcal X}(t)$ is a homotopy $3$-sphere. In either case
both homotopy group statements hold for ${\mathcal X}(t)$. Even more
is true: The distance-decreasing map ${\mathcal X}(t')\to {\mathcal
X}(t)$ given by Proposition~\ref{surgerydist} is either a homotopy
equivalence or a degree one map of ${\mathcal X}(t')\to {\mathcal
X}(t)$. In either case, it induces an injection of $\pi_3({\mathcal
X}(t'))\to \pi_3({\mathcal X}(t))$. In this way a non-zero element
in $\xi(t_0)\in\pi_3({\mathcal X}(t_0))$ produces a family of
non-zero elements $\xi(t)\in \pi_3({\mathcal X}(t))$ with the
property that under Ricci flow these elements agree and at a surgery
time $t$ the degree one map constructed in
Proposition~\ref{surgerydist} sends $\xi(t')$ to $\xi(t)$ for all
$t'<t$ sufficiently close to it. Since $\pi_2({\mathcal X}(t))$ is
trivial for all $t$, we identify $\xi(t)$ with a homotopy class of
maps of $S^2$ to $\Lambda {\mathcal X}(t)$. We now define a function
$W_\xi(t)$ by associating to each $t$ the invariant $W(\xi(t))$.

Here is the result that is analogous to Lemma~\ref{lem:w2-forward-difference-quotient}.


\begin{proposition}[Forward difference quotient for $W_\xi$]
\label{prop:w3-forward-difference-quotient}
\label{W_3fordiff}
\uses{prop:eventual-trivial-pi2,def:w-xi-spanning-disk-functional}
Suppose that $({\mathcal M},G)$ is a Ricci flow with surgery as in
Theorem~\ref{MAIN}. Let ${\mathcal X}$ be a path of components of ${\mathcal
M}$ defined for all $t\in[t_0,t_1]$ with $\pi_2({\mathcal X}(t_0))=0$. Suppose
that $\xi\in \pi_3( X(t_0),*)$ is a non-trivial element. Then the function
$W_\xi(t)$ satisfies the following inequality in the sense of forward
difference quotients:
$$\frac{d W_\xi(t)}{dt}\le -2\pi-\frac{1}{2}R_\mathrm{min}(t)W_\xi(t).$$ Also, for every $t\in [t_0,t_1]$ that is not a
surgery time the function $W_\xi(t)$ is continuous at $t$. Lastly,
if $t$ is a surgery time then
$$W_\xi(t)\le \liminf_{t'\rightarrow t^-}W_\xi(t').$$
\end{proposition}

In the next subsection we assume this result and use it to complete
the proof.



\subsection{Proof of Theorem~{\ref{thm:finite-time-extinction-main}} assuming
Proposition~{\ref{prop:w3-forward-difference-quotient}}}

\begin{proof}
\uses{prop:eventual-trivial-pi2,prop:w3-forward-difference-quotient,MAIN,lem:forward-difference-quotient}
 According to Proposition~\ref{prop:eventual-trivial-pi2} there is $T_1$ such that
every component of $M_T$ has trivial $\pi_2$ for every $T\ge T_1$.
Suppose that Theorem~\ref{thm:finite-time-extinction-main} does not hold for this Ricci flow
with surgery. We consider a path of components ${\mathcal X}(t)$ of
${\mathcal M}$ defined for $[T_1,T_2]$.  We shall show that there is
a uniform upper bound to $T_2$.

\begin{claim}[$\mathcal{X}(T_1)$ has non-trivial $\pi_3$]
\label{claim:x-t1-nontrivial-pi3}
 ${\mathcal X}(T_1)$ has non-trivial $\pi_3$.
\end{claim}

\begin{proof}
By hypothesis the fundamental group of $M_0$ is a free product of
infinite cyclic groups and finite groups. This means that the same
is true for the fundamental group of each component of $M_t$ for
every $t\ge 0$, and in particular it is true for ${\mathcal
X}(T_0)$. But we know that $\pi_2({\mathcal X}(T_0))=0$.

\begin{claim}[Non-trivial free product or $\mathbb{Z}$ fundamental group forces $\pi_2\ne 0$]
\label{claim:free-product-implies-pi2-nonzero}
Let $X$ be a compact $3$-manifold. If $\pi_1(X)$ is a non-trivial
free product or if $\pi_1(X)$ is isomorphic to $\Zee$, then
$\pi_2(X)\not=0$.
\end{claim}


\begin{proof}
See \cite{Hempel}, Theorem 5.2 on page 56 (for the case of a copy of $\Zee$)
and \cite{Hempel} Theorem 7.1 on page 66 (for the case of a free product
decomposition).
\end{proof}

Thus, it follows that $\pi_1({\mathcal X}(T_1))$ is a finite group
(possibly trivial). But a $3$-manifold with finite fundamental group
has a universal covering that is a compact $3$-manifold with trivial
fundamental group. Of course, by Poincar\'e duality any simply
connected $3$-manifold  is a homotopy $3$-sphere. It follows
immediately that $\pi_3({\mathcal X}(T_1))\cong \Zee$. This
completes the proof of the claim.
\end{proof}


Now we can apply Proposition~\ref{prop:w3-forward-difference-quotient} to our path of components
${\mathcal X}$ defined for all $t\in [T_1,T_2]$. First recall by
Theorem~\ref{MAIN} that the curvature of $({\mathcal M},G)$ is pinched toward
positive which implies that $R_{\min}(t)\ge (-6)/(1+4t)$. Let $w(t)$ be the
function satisfying the differential equation
$$w'(t)=-2\pi+\frac{3}{1+4t}w(t)$$
with initial condition $w(T_1)=W_\xi(T_1)$. According to
Proposition~\ref{prop:w3-forward-difference-quotient} and Proposition~\ref{lem:forward-difference-quotient} we
see that $W_\xi(t)\le w(t)$ for all $t\in [T_1,T_2]$. But direct
integration shows that
$$w(t)=W_\xi(T_1)\frac{(4t+1)^{3/4}}{(4T_1+1)^{3/4}}+2\pi
(4T_0+1)^{1/4}(4t+1)^{3/4}-2\pi (4t+1).$$ This clearly shows that
$w(t)$ becomes negative for $t$ sufficiently large, how large
depending only on $W_\xi(T_1)$ and $T_1$. On the other hand, since
$W_\xi(t)$ is the infimum of areas of disks, $W_\xi(t)\ge 0$ for all
$t\in [T_1,T_2]$. This proves that $T_2$ is less than a constant
that depends only on $T_1$ and on the component ${\mathcal X}(T_1)$.
Since there are only finitely many connected components of
$M_{T_1}$, this shows that $T_2$ depends only on $T_1$ and the
Riemannian manifold $M_{T_1}$. This completes the proof of
Theorem~\ref{thm:finite-time-extinction-main} modulo Proposition~\ref{prop:w3-forward-difference-quotient}.
\end{proof}

Thus, to complete the argument for Theorem~\ref{thm:finite-time-extinction-main} it remains
only to prove Proposition~\ref{prop:w3-forward-difference-quotient}.


\subsection{Continuity of $W_\xi(t)$}

In this subsection we establish the two continuity conditions for
$W_\xi(t)$ stated in Proposition~\ref{prop:w3-forward-difference-quotient}.

\begin{claim}[Continuity of $W_\xi$ away from surgery times]
\label{claim:w-xi-continuity-nonsurgery}
If $t$ is not a surgery time, then $W_\xi(t)$ is continuous at $t$.
\end{claim}

\begin{proof}
Since $t$ is not a surgery time, a family $\Gamma(t)\colon S^2\to
\Lambda {\mathcal X}(t)$ is also a family $\Gamma(t')\colon
S^2\to\Lambda{\mathcal X}(t')$ for all nearby $t'$. The minimal
spanning disks for the elements of $\Gamma(t)(x)$ are  also spanning
disks in the nearby ${\mathcal X}(t')$ and their areas vary
continuously with $t$. But the maximum of the areas of these disks
is an upper bound for $W_\xi(t)$. This immediately implies that
$W_\xi(t)$ is upper semi-continuous at $t$.

The result for lower semi-continuity is the same as in the case of
$2$-spheres. Given a time $t$ distinct from a surgery time and a
family $\Gamma\colon S^2\to \Lambda {\mathcal X}(t')$ for a time
$t'$ near $t$ we can view the family $\Gamma$ as a map to $\Lambda
{\mathcal X}(t)$. The areas of all minimal spanning disks for the
loops represented by points $\Gamma$ measured in ${\mathcal X}(t)$
are at most $(1+\eta(|t-t'|))$ times their areas measured in
${\mathcal X}(t')$, where $\eta(|t-t'|)$ is a function going to zero
as $|t-t'|$ goes to zero. This immediately implies the lower
semi-continuity at the non-surgery time $t$.
\end{proof}


\begin{claim}[Upper semicontinuity of $W_\xi$ at surgery times]
\label{claim:w-xi-surgery-semicontinuity}
Suppose that $t$ is a surgery time. Then
$$W_\xi(t)\le \mathrm{liminf}_{t'\rightarrow t^-}W_\xi(t').$$
\end{claim}

\begin{proof}
This is immediate from the fact from Proposition~\ref{surgerydist} that for any
$\eta>0$ for every $t'<t$ sufficiently close to $t$ there is a homotopy
equivalence ${\mathcal X}(t')\to {\mathcal X}(t)$ which is a
$(1+\eta)$-Lipschitz map.
\end{proof}

To prove Proposition~\ref{prop:w3-forward-difference-quotient} and hence
Theorem~\ref{thm:finite-time-extinction-main}, it remains to prove the forward difference
quotient statement for $W_\xi(t)$ given in
Proposition~\ref{prop:w3-forward-difference-quotient}.






\subsection{A further reduction of
Proposition~{\ref{prop:w3-forward-difference-quotient}}}

Let $\Gamma\colon S^2\to \Lambda {\mathcal X}(t_0)$ be a family. We
must construct an appropriate deformation of the family of loops
$\Gamma$ in order to establish Proposition~\ref{prop:w3-forward-difference-quotient}. Now we
are ready to state the more technical estimate for the evolution of
$W(\Gamma)$ under Ricci flow that will imply the forward difference
quotient result for $W_\xi(t)$ stated in
Proposition~\ref{prop:w3-forward-difference-quotient}. Here is the result that shows a
deformation as required exists.

\begin{definition}[Comparison ODE solution $w_{a,t'}$]
\label{def:w-a-t-ode-solution}
Let $(M,g(t)),\ t_0\le t\le t_1$, be a Ricci flow on a compact
$3$-manifold. For any $a$ and any $t'\in [t_0,t_1]$ let
$w_{a,t'}(t)$ be the solution to the differential equation
\begin{equation}\label{wdiffeqn}
\frac{dw_{a,t'}}{dt}=-2\pi-\frac{1}{2}R_{\min}(t)w_{a,t'}(t)
\end{equation}
with initial condition $w_{a,t'}(t')=a$. We also denote $w_{a,t_0}$
by $w_a$.
\end{definition}

\begin{proposition}[Existence of the curve-shrinking deformation]
\label{prop:curve-deformation-existence}
\label{XI}
\uses{prop:w3-forward-difference-quotient}
Let $(M,g(t)),\ t_0\le t\le t_1$, be a Ricci flow on a compact
$3$-manifold.  Fix a map $\Gamma$ of $S^2$ to $\Lambda M$
 whose image consists of homotopically trivial loops and $\zeta>0$. Then there
is a continuous family $\widetilde\Gamma(t),\ t_0\le t\le t_1$, of
maps $S^2\to \Lambda M$ whose image consists of homotopically
trivial loops with $[\widetilde\Gamma(t_0)]=[\Gamma]$ in
$\pi_3(M,*)$ such that for each $c\in S^2$ we have
$|A(\widetilde\Gamma(t_0)(c))-A(\Gamma(c))|<\zeta$ and furthermore,
one of the following two alternatives holds:
\begin{enumerate}
\item[(i)] The length of $\widetilde\Gamma(t_1)(c)$ is less than $\zeta$.
\item[(ii)] $A(\widetilde\Gamma(t_1)(c))\le w_{A(\widetilde\Gamma(t_0)(c))}(t_1)+\zeta$.
\end{enumerate}
\end{proposition}

Before proving this result we shall show it implies the forward difference
quotient result in Proposition~\ref{prop:w3-forward-difference-quotient}. Let ${\mathcal X}$ be a path
of components. Suppose that $\pi_2({\mathcal X}(t),x_0)=0$ for all $t$. Fix
$t_0$ and fix a non-trivial element $\xi\in \pi_3({\mathcal X}(t_0),x_0)$,
which we identify with a non-trivial element in $\xi\in \pi_2(\Lambda{\mathcal
X}(t_0),*)$. Fix an interval $[t_0,t_1]$ with the property that there are no
surgery times in the interval $(t_0,t_1]$. Restricting to this interval the
family ${\mathcal X}(t)$ is a Ricci flow on ${\mathcal X}(t_0)$. In particular,
all the ${\mathcal X}(t)$ are identified under the Ricci flow.
 Let
$w(t)$ be the solution to Equation~(\ref{wdiffeqn}) with value
$w(t_0)=W_\xi({\mathcal X},t_0)$. We shall show that $W_\xi(t_1)\le
w(t_1)$. Clearly, once we have this estimate, taking limits as $t_1$
approaches $t_0$ establishes the forward difference quotient result
at $t_0$.

\begin{definition}[Auxiliary integral $A(t)$]
\label{def:a-function-integral}
Let $A(t)=\int_{t'}^t\frac{1}{2}R_{\min}(s)ds$.
\end{definition}

Direct integration shows the following:


\begin{claim}[Explicit formula for $w_{a,t'}$]
\label{claim:w-a-t-explicit-formula}
\label{vclaim}
We have
$$w_{a,t'}(t'')=\exp(-A(t''))\left(a-2\pi\int_{t'}^{t''}\mathrm{exp}(A(t))dt\right).$$ If $a'>a$, then for $t_0\le t'<t''\le t_1$,
we have
$$w_{a',t'}(t'')=w_{a,t'}(t'')+(a'-a)\exp(-A(t'')).$$
\end{claim}







 The next thing to
establish is the following.

\begin{lemma}[Families of short loops represent the trivial $\pi_3$ class]
\label{lem:short-loops-trivial-pi3}
\label{shortpi2triv}
Given a compact Riemannian manifold $(X,g)$ with $\pi_2(X)=0$. Then
there is $\zeta>0$ such that if $\xi\in \pi_3({\mathcal X})$ is
represented by a family $\Gamma\colon S^2\to\Lambda X$ with the
property that for every $c\in S^2$ the length of the loop
$\Gamma(c)$ is less than $\zeta$, then $\xi$ is the trivial homotopy
element.
\end{lemma}

\begin{proof}
We choose $\zeta$ smaller than the injectivity radius of $(X,g)$.
Then any pair of points at distance less than $\zeta$ apart are
joined by a unique geodesic of length less than $\zeta$.
Furthermore, the geodesic varies smoothly with the points. Given a
map $\Gamma\colon S^2\to \Lambda X$ such that every loop of the form
$\Gamma(c)$ has length at most $\zeta$, we consider the map $f\colon
S^2\to X$ defined by $f(c)=\Gamma(c)(x_0)$, where $x_0$ is the base
point of the circle. Then we can join each point $\Gamma(c)(x)$ to
$\Gamma(c)(x_0)$ by a geodesic of length at most $\zeta$ to fill out
a map of the disk $\widehat \Gamma(c)\colon D^2\to X$. This disk is
smooth except at the point $\Gamma(c)(x_0)$. The disks
$\widehat\Gamma(c)$ fit together as $c$ varies to make a continuous
family of disks parameterized by $S^2$ or equivalently a map
$S^2\times D^2$ into $X$ whose boundary is the family of loops
$\Gamma(c)$. Now shrinking the loops $\Gamma(c)$ across the disks
$\hat\Gamma(c)$ to $\Gamma(c)(x_0)$ shows that the family $\Gamma$
is homotopic to a $2$-sphere family of constant loops at different
points of $X$. Since we are assuming that $\pi_2(X)$ is trivial,
this means the family of loops is in fact trivial as an element of
$\pi_2(\Lambda X,*)$, which  means that the original element $\xi\in
\pi_3(X)$ is trivial.
\end{proof}

Notice that this argument also shows the following:

\begin{corollary}[Short loops bound small-area disks]
\label{cor:short-loop-bounds-small-disk}
\label{shlnsharea}
\uses{lem:short-loops-trivial-pi3,claim:w-a-t-explicit-formula}
Let $(X,g)$ be a compact Riemannian manifold. Given $\eta>0$ there is a
$0<\zeta<\eta/2$ such that any $C^1$-loop $c\colon S^1\to X$ of length less
than $\eta$ bounds a disk in $X$ of area less than $\eta$.
\end{corollary}

Now we return to the proof that Proposition~\ref{prop:curve-deformation-existence} implies
Proposition~\ref{prop:w3-forward-difference-quotient}. We consider the restriction of the
path ${\mathcal X}$ to the time interval $[t_0,t_1]$. As we have
already remarked, since there are no surgery times in $(t_0,t_1]$,
this restriction is a Ricci flow and all the ${\mathcal X}(t)$ are
identified with each other under the flow.
 Let $w(t)$ be the solution to Equation~(\ref{wdiffeqn})
with initial condition $w(t_0)=W_\xi(t_0)$. There are two cases to
consider: (i) $w(t_1)\ge 0$ and $w(t_1)< 0$.

Suppose that $w(t_1)\ge 0$. Let $\eta>0$ be given. Then by Claim~\ref{claim:w-a-t-explicit-formula}
and  Corollary~\ref{cor:short-loop-bounds-small-disk}, there is $0<\zeta<\eta/2$ such that the
following two conditions hold:
\begin{enumerate}
\item[(a)] Any loop in ${\mathcal X}(t_1)$ of length less than $\zeta$
bounds a disk of area less than $\eta$.
\item[(b)] For every $a\in [0,W_\xi(t_0)+2\zeta]$
the solution $w_{a}$ satisfies $w_{a}(t_1)<w(t_1)+\eta/2$.
\end{enumerate}

Now fix a map $\Gamma\colon S^2\to \Lambda{\mathcal X}(t_0)$, whose
image consists of homotopically trivial loops, with $[\Gamma]=\xi$,
and with $W(\Gamma)<W_\xi(t_0)+\zeta$. According to
Proposition~\ref{prop:curve-deformation-existence} there is a one-parameter family
$\widetilde\Gamma(t),\ t_0\le t\le t_1$, of maps $S^2\to
\Lambda{\mathcal X}(t)$, whose images consist of homotopically
trivial loops, with $[\widetilde\Gamma(t_0)]=[\Gamma]=\xi$ such that
for every $c\in S^2$ we have
$A(\widetilde\Gamma(t_0)(c))<A(\Gamma(c))+\zeta$ and one of the
following holds
\begin{enumerate}
\item[(i)] the length of $\widetilde\Gamma(t_1)(c)$ is less than $\zeta$, or
\item[(ii)] $$A(\widetilde\Gamma(t_1)(c))<
w_{A(\widetilde\Gamma(t_0)(c))}(t_1)+\zeta.$$
\end{enumerate}
Since $A(\widetilde\Gamma(t_0)(c))<A(\Gamma(c))+\zeta<W_\xi(t_0)+2\zeta$, it
follows from our choice of $\zeta$ that for every $c\in S^2$ either
\begin{enumerate}
\item[(a)] $\widetilde\Gamma(t_1)(c)$ has length less than $\zeta$ and hence bounds a disk
of area less than $\eta$, or
\item[(b)]
$A(\widetilde\Gamma(t_1)(c))<w_{W_\xi(t_0)+2\zeta}(t_1)+\zeta<w(t_1)+\eta/2+\eta/2=
w(t_1)+\eta$.
\end{enumerate}
Since we are assuming that $w(t_1)\ge 0$, it now follows that for
every $c\in S^2$ we have $A(\widetilde\Gamma(t_1)(c))<w(t_1)+\eta$,
and hence $W(\widetilde\Gamma(t_1))< w(t_1)+\eta$. This shows that
for every $\eta>0$ we can find a family $\widetilde\Gamma(t)$ with
$\widetilde\Gamma(t_0)$ representing $\xi$ and with
$W(\widetilde\Gamma(t_1))< w(t_1)+\eta$. This completes the proof of
Proposition~\ref{prop:curve-deformation-existence} when $w(t_1)\ge 0$.

Now suppose that $w(t_1)<0$. In this case, we must derive a contradiction since
clearly it must be the case that for any one-parameter family
$\widetilde\Gamma(t)$ we have $W(\widetilde\Gamma(t_1))\ge 0$. We fix $\eta>0$
such that $w(t_1)+\eta<0$. Then using Lemma~\ref{lem:short-loops-trivial-pi3} and
Claim~\ref{claim:w-a-t-explicit-formula}, we fix $\zeta$ with $0<\zeta<\eta/2$ such that:
\begin{enumerate}
\item[(i)] If $\Gamma\colon S^2\to \Lambda {\mathcal X}(t_1)$ is a family
of loops and each loop in the family is of length less than $\zeta$,
then the family is homotopically trivial.
\item[(ii)] For any $a\in [0,W_\xi(t_0)+2\zeta]$ we have $w_a(t_1)<w(t_1)+\eta/2$.
\end{enumerate}

We fix a map $\Gamma\colon S^2\to {\mathcal X}(t_0)$ with
$[\Gamma]=\xi$ and with $W(\Gamma)<W_\xi(t_0)+\zeta$. Now according
to Proposition~\ref{prop:curve-deformation-existence} there is a family of maps
$\widetilde\Gamma(t)\colon S^2\to \Lambda{\mathcal X}(t)$ with
$[\widetilde\Gamma(t_0)]=[\Gamma]=\xi$ and for every $c\in S^2$ we
have $A(\widetilde\Gamma(t_0)(c))<A(\Gamma(c))+\zeta$ and also
either $A(\widetilde\Gamma(t_1)(c))\le
w_{A(\widetilde\Gamma(t_0)(c))}(t_1)+\zeta$ or the length of
$\widetilde\Gamma(t_1)(c)$ is less than $\zeta$. It follows that for
every $c\in S^2$ we have $A(\widetilde\Gamma(t_0)(c))\le
W(\Gamma)+\zeta<W_\xi(t_0)+2\zeta$. From the choice of $\zeta$ this
means that
$$A(\widetilde\Gamma(t_1)(c))< w(t_1)+\eta/2+\zeta<w(t_1)+\eta<0$$ if
the length of $\widetilde\Gamma(t_1)(c)$ is at least $\zeta$. Of course, by
definition $A(\widetilde\Gamma(t_1)(c))\ge 0$ for every $c\in S^2$. This
implies  that for every $c\in S^2$ the loop $\widetilde\Gamma(t_1)(c)$ has
length less than $\zeta$. By Lemma~\ref{lem:short-loops-trivial-pi3} this implies that
$\widetilde\Gamma(t_1)$ represents the trivial element in
$\pi_2(\Lambda{\mathcal X}(t_1))$, which is a contradiction.

At this point, all that it remains to do in order to complete the
proof of Theorem~\ref{thm:finite-time-extinction-main} is to establish Proposition~\ref{prop:curve-deformation-existence}.
The rest of this chapter is devoted to doing that.






\section{Curve-shrinking flow}\label{sect:csf}

Given $\Gamma$, the idea for constructing the one-parameter family
$\widetilde \Gamma(t)$ required by Proposition~\ref{prop:curve-deformation-existence} is to evolve
an appropriate approximation $\widetilde \Gamma(t_0)$ of $\Gamma$ by
the curve-shrinking flow. Suppose that $(M,g(t)),\ t_0\le t\le t_1$,
is a Ricci flow of compact manifolds and that $c\colon S^1\times
[t_0,t_1]\to (M,g(t_0))$ is a family of parameterized, immersed
$C^2$-curves. We denote by $x$ the parameter on the circle. Let
$X(x,t)$ be the tangent vector $\partial c(x,t)/\partial x$ and let
$S(x,t)=X(x,t)/(|X(x,t)|_{g(t)})$ be the unit tangent vector to $c$.
We denote by $s$ the arc length parameter on $c$. We set
$H(x,t)=\nabla _{S(x,t)}S(x,t)$, the curvature vector of $c$ with
respect to the metric $g(t)$.
 We define the curve-shrinking
flow by
$$\frac{\partial c(x,t)}{\partial t}=H(x,t),$$
where $c(x,t)$ is a one-parameter family of curves and $H(x,t)$ is
the curvature vector of the curve $c(\cdot,t)$ at the point $x$ with
respect to the metric $g(t)$. We denote by $k(x,t)$ the curvature
function: $k(x,t)=|H(x,t)|_{g(t)}$. We shall often denote the
one-parameter family of curves by $c(\cdot,t)$. Notice that if
$c(x,t)$ is a curve-shrinking flow and if $x(y)$ is a
reparameterization of the domain circle then $c'(y,t)=c(x(y),t)$ is
also a curve-shrinking flow.

\begin{claim}[Short-time existence of the curve-shrinking flow]
\label{claim:curve-shrinking-short-time-existence}
\label{curvshrexist}
\uses{prop:curve-deformation-existence}
For any immersed $C^2$-curve $c\colon S^1\to (M,g(t_0))$ there is a
curve-shrinking flow $c(x,t)$ defined for $t\in [t_0,t_1')$ for some
$t_1'>t_0$ with the property that each $c(\cdot,t)$ is an immersion.
Either the curve-shrinking flow extends to a curve-shrinking flow
that is a family of immersions defined at $t_1'$ and beyond, or
$\max_{x\in S^1}k(x,t)$ blows up as $t$ approaches $t_1'$ from
below.
\end{claim}

For a proof of this result, see Theorem 1.13 in \cite{AG}.

\subsection{The proof of Proposition~{\ref{prop:curve-deformation-existence}} in a simple case}

The main technical hurdle to overcome is that in general the curve
shrinking flow may not exist if the original curve is not immersed
and even if the original curve is immersed the curve-shrinking flow
can develop singularities, where the curvature of the curve goes to
infinity. Thus, we may  not be able to define the curve-shrinking
flow as a flow defined on the entire interval $[t_0,t_1]$, even
though the Ricci flow is defined on this entire interval. But to
show the idea of the proof,  let us suppose for a moment that the
starting curve is embedded and that no singularities develop in the
curve-shrinking flow and show how to prove the result.

\begin{lemma}[Area bound for embedded curve-shrinking flow]
\label{lem:embedded-curve-shrinking-area-bound}
\label{embcase}
Suppose that $c\in \Lambda M$ is a homotopically trivial, embedded
$C^2$-loop. and suppose that there is a curve-shrinking flow
$c(x,t)$ defined for all $t\in [t_0,t_1]$ with each $c(\cdot,t)$
being an embedded smooth curve. Consider the function $A(t)$ which
assigns to $t$ the minimal area of a spanning disk for $c(\cdot,t)$.
Then $A(t)$ is a continuous function of $t$ and
$$\frac{dA}{dt}(t)\le -2\pi-\frac{1}{2}R_{\min}(t)A(t)$$
in the sense of forward difference quotients.
\end{lemma}


\begin{proof}
\uses{claim:minimal-sphere-area-variation}
According to results of Hildebrandt and Morrey, \cite{Hildebrandt}
and \cite{Morrey}, for each $t\in [t_0,t_1]$, there is a smooth
minimal disk spanning $c(\cdot,t)$. Fix $t'\in [t_0,t_1)$ and
consider a smooth minimal disk $D\to (M,g(t'))$ spanning
$c(\cdot,t)$. It is immersed, see \cite{HardtSimon} or \cite{GL}.
The family $c(\cdot,t)$ for $t$ near $t'$ is an isotopy of
$c(\cdot,t')$. We can extend this to an ambient isotopy
$\varphi_t\colon M\to M$ with $\varphi_{t'}=\Id$. We impose
coordinates $\{x_\alpha\}$ on $D$; we let $h_{\alpha\beta}(t')$ be
the metric induced  on $\varphi_{t'}(D)$ by $g(t')$, and we let $da$
be the area form induced by the Euclidean coordinates on $D$. We
compute
$$\frac{d}{dt}\bigl|_{t=t'}\bigr.\Area(\varphi_t(D))
=\frac{d}{dt}\bigl|_{t=t'}\bigr.\int_{\varphi_t(D)}\sqrt{\mathrm{det}(h_{\alpha\beta})(t)}da.$$ Of course,
\begin{eqnarray*}
\frac{d}{dt}\bigl|_{t=t'}\bigr.\int_{\varphi_t(D)}\sqrt{\mathrm{det}(h_{\alpha\beta}(t))}da & = & -\int_{\varphi_{t'}(D)}\Bigl(\Tr\,\Ric^T\Bigr)\sqrt{\det(h_{\alpha\beta}(t))}da \\
& & +\int_{\varphi_{t'}(D)}\operatorname{div}\left(\frac{\partial
\varphi_{t'}}{dt}\right)^T\sqrt{\det(h_{\alpha\beta}(t))}da.
\end{eqnarray*}
 Here, $\Ric^T$
denotes the restriction of the Ricci curvature of $g(t')$ to the
tangent planes of $\varphi_{t'}(D)$ and
$\frac{\partial(\varphi_{t'})^T}{\partial t}$ is the component of
$\varphi_{t'}$ tangent to $\varphi_{t'}(D)$. Setting $\hat A$ equal
to the second fundamental form of $\varphi_{t'}(D)$, using the fact
that $\varphi_{t'}(D)$ is minimal and arguing as in the proof of
Claim~\ref{claim:minimal-sphere-area-variation}, we have
\begin{eqnarray*}
\lefteqn{-\int_{\varphi_{t'}(D)}\Bigl(\operatorname{Tr}\,\Ric^T\Bigr)\sqrt{\det(h_{\alpha\beta}(t'))}da } \\
 & = &
-\int_{\varphi_{t'}(D)}K_{\varphi_{t'}(D)}da
-\frac{1}{2}\int_{\varphi_{t'}(D)}(|\hat A|^2+R)da \\
& \le & -\int_{\varphi_{t'}(D)}K_{\varphi_{t'}(D)}\sqrt{\mathrm{det}(h_{\alpha\beta}(t'))}da-\frac{1}{2}\Area{\varphi_{t'}(D)}\min_{x\in M}\{R(x,t')\}.
 \end{eqnarray*}

Integration by parts shows that
$$\int_{\varphi_{t'}(D)}\operatorname{div}\left(\frac{\partial
\varphi_{t'}}{dt}\right)^T\sqrt{\mathrm{det}(h_{\alpha\beta}(t'))}da=-\int_{ \varphi_{t'}(\partial
D)}\left(\frac{d\varphi_t}{dt}\bigl|_{t=t'}\bigr.\right)\cdot n
ds,$$ where $n$ is the inward pointing normal vector to
$\varphi_{t'}(D)$ along $\varphi_{t'}(\partial D)$. Of course, by
definition, if the variation along the boundary is given by the
curve-shrinking flow, then along $\varphi_{t'}(\partial D)$ we have
$$\left(\frac{d\varphi_t}{dt}\bigl|_{t=t'}\bigr.\right)\cdot n=k_\mathrm{geod}.$$

Thus, we have
\begin{eqnarray*}\lefteqn{\frac{d}{dt}\bigl|_{t=t'}\bigr.\int_{\varphi_t(D)}\sqrt{\mathrm{det}(h_{\alpha\beta}(t))}da } \\\ & \le &
-\int_{\varphi_{t'}(D)}K_{\varphi_{t'}(D)}da-\int_{\varphi_{t'}(\partial
D)}k_\mathrm{geod}ds-\frac{1}{2}R_{\min}(t')\Area(\varphi_{t'}(D)).\end{eqnarray*} Of course, the Gauss-Bonnet
theorem allows us to rewrite this as
$$\frac{d}{dt}\bigl|_{t=t'}\bigr.\int_{\varphi_t(D)}\sqrt{\mathrm{det}(h_{\alpha\beta}(t))}da \le -2\pi-\frac{1}{2}R_{\min}(t')\Area(\varphi_{t'}(D)).$$
\end{proof}


Let $\psi(t)$ be the solution to the ODE
$$\psi'(t)=-2\pi-\frac{1}{2}R_{\min}(t)\psi(t)$$
with $\psi(t^-)=A(t^-)$. The following is immediate from the
previous lemma and Lemma~\ref{lem:forward-difference-quotient}.

\begin{corollary}[Area comparison for embedded curve-shrinking flow]
\label{cor:embedded-curve-shrinking-area-comparison}
With notation and assumptions as above, if the curve-shrinking flow
is defined on the interval $[t^-,t^+]$ and if the curves
$c(\cdot,t)$ are embedded for all $t\in [t^-,t^+]$ then
$$A(t^+)\le \psi(t^+).$$
\end{corollary}

Actually, the fact that the loops in the curve-shrinking flow are
embedded is not essential in dimensions $\ge 3$.

\begin{lemma}[Area bound for immersed curve-shrinking flow in dimension at least 3]
\label{lem:immersed-curve-shrinking-area-bound}
\label{nonembeasycase}
Suppose that the dimension of $M$ is at least $3$ and that
$c(\cdot,t)$ is a $C^2$-family of homotopically trivial, immersed
curves satisfying the curve-shrinking equation defined for $t^-\le
t\le t^+$. For each $t$, let $A(t)$ be the infimum of the areas of
spanning disks for  $c(\cdot,t)$. Then $A(t)$ is a continuous
function and, with $\psi$ as above, we have
$$A(t^+)\le \psi(t^+).$$
\end{lemma}


\begin{proof}
\uses{lem:embedded-curve-shrinking-area-bound,lem:forward-difference-quotient}
 We first remark that continuity has already been established.
To show the inequality, we begin  with a claim.

\begin{claim}[Reduction to a $C^2$-approximating family]
\label{claim:curve-shrinking-approximation-suffices}
It suffices to prove the following for every $\delta>0$. There is a
$C^2$-family $\hat c(x,t)$ of immersions within $\delta$ in the
$C^2$-topology to $c(x,t)$ defined on the interval $[t^-,t^+]$ such
that
$$A(t^+)\le \psi_{\delta,\hat c}(t^+)$$
where $\psi_{\delta,\hat c}$ is the solution of the ODE
$$\psi_{\delta,\hat c}'(t)=-2\pi+2\delta L_{\hat c}(t)-\frac{1}{2}R_\mathrm{min}(t)\psi_{\delta,\hat c}(t)$$ with value $A(\hat c(t^-))$ at
$t^-$, and where $L_{\hat c}(t)$ denotes the length of the loop
$\hat c(\cdot,t)$.
\end{claim}

\begin{proof} (of the claim) Suppose that for each $\delta$ there is such
a $C^2$-family as in the statement of the claim. Take a sequence
$\delta_n$ tending to zero, and let $\hat c_n(\cdot,t)$ be a family
as in the claim for $\delta_n$. Then by the continuity of the
infimum of areas of the spanning disk in the $C^1$-topology, we see
that
$$\lim_{n\rightarrow \infty}A(\hat c_n(\cdot, t^\pm))=A(c(\cdot,t^\pm)).$$
Since the $\hat c_n(x,t)$ converge in the $C^2$-topology to
$c(x,t)$, the lengths $L(\hat c_n(t))$ are uniformly bounded and the
$A(\hat c_n(t^-))$ converge to $A(c(t))$. Thus, the
$\psi_{\delta_n,\hat c_n}$ converge uniformly to $\psi$ on
$[t^-,t^+]$, and taking limits shows the required inequality for
$A(c(t))$, thus proving the claim.
\end{proof}

Now we return to the proof of the lemma.  Let $\hat c(x,t)$ be a
generic $C^2$-immersion  sufficiently close to $c(x,t)$ in the
$C^2$-topology so that the following hold:
\begin{enumerate}
\item[(1)] the difference of the curvature of $\hat c$ and of $c$ at
every $(x,t)$ is a vector of length less than $\delta$,
\item[(2)] the difference of $\partial \hat c/\partial t$ and $\partial
c/\partial t$ is a vector of length less than $\delta$,
\item[(3)] the ratio of the arc lengths of $\hat c$ and $c$ at every
$(x,t)$ is between $(1-\delta)$ and $(1+\delta)$.
\end{enumerate}
The generic family $\hat c(x,t)$ consists of embedded curves for all
but a finite number of $t\in [t^-,t^+]$ and at the exceptional $t$
values the curve is immersed. Let $t_1<t_2<\cdots <t_k$ be the
values of $t$ for which $\hat c(\cdot,t)$ is not embedded. We set
$t_0=t^-$ and $t_{k+1}=t^+$. Notice that it suffices to show that
$$A(\hat c(t_{i+1}))-A(\hat c(t_{i}))\le
\psi_{\delta,\hat c}(t_{i+1})-\psi_{\delta,\hat c}(t_i)$$ for
$i=0,\ldots,k$. To establish this inequality for the interval
$[t_i,t_{i+1}]$,  by continuity it suffices to establish the
corresponding inequality for every compact subinterval contained in
the interior of this interval. This allows us to assume that the
approximating family is a family of embedded curves. Let the
endpoints of the parameterizing interval  be denoted $a$ and $b$.
Fix $t'\in [a,b]$ and let $D$ be a minimal disk spanning $\hat
c(\cdot, t')$, and let $\varphi_t$ be an isotopy as in the argument
given the proof of Lemma~\ref{lem:embedded-curve-shrinking-area-bound}. According to this argument
we have
$$\frac{d}{dt}A(\hat c(t))|_{t=t'}\le -2\pi -\frac{1}{2}R_\mathrm{min}(t')A(c(t'))+\int_{c(x,t')}\left[k_\mathrm{geod}(\hat
c)-\left(\frac{d\varphi_t}{dt}|_{t=t'}\right)\cdot n\right] ds$$ in
the sense of forward difference quotients. The restriction of
$\frac{d\varphi_t}{dt}|_{t=t'}$ to the boundary of $D$ agrees with
$\partial \hat c(x,t)/\partial t$. Hence, by our conditions on the
approximating family, and since for $c(\cdot,t)$ the corresponding
quantities are equal,
$$\left|k_\mathrm{geod}(\hat c)-\left(\frac{d\varphi_t}{dt}|_{t=t'}\right)\cdot
n\right|<2\delta.$$ Integrating over the circle implies that
$$\frac{d}{dt}A(\hat c(t))|_{t=t'}\le -2\pi -\frac{1}{2}R_\mathrm{min}(t')A(c(t'))+2\delta L_{\hat c}(t).$$
 The result is then
immediate from Lemma~\ref{lem:forward-difference-quotient}.
\end{proof}


\subsection{Basic estimates for curve-shrinking}

Let us establish some elementary formulas. To simplify the formulas we often
drop the variables $x,t$ from the notation, though they are understood to be
there.

\begin{lemma}[First variation formulas for curve-shrinking flow]
\label{lem:curve-shrinking-first-variation}
\label{1stcurvshr}
Assume that $(M,g(t)),\ t_0\le t\le t_1$, is a Ricci flow and that
$c=c(x,t)$ is a solution to the curve-shrinking flow. We have vector
fields $X=\partial /\partial x$ and $H=\partial /\partial t$ defined
on the domain surface. We denote by $|X|_{c^*g}^2$ the function on
the domain surface whose value at $(x,t)$ is $|(X(x,t))|_{g(t)}^2$.
We define $S=|X|_{c^*g}^{-1}X$, the unit vector in the $x$-direction
measured in the evolving metric. Then,
$$\frac{\partial}{\partial t}(|X|_{c^*g}^2)(x,t)=-2\Ric_{g(t)}(X(x,t),X(x,t))-2k^2|X(x,t)|_{g(t)}^2,$$
and
$$[H,S](x,t)=\left(k^2+\Ric_{g(t)}(S(x,t),S(x,t))\right)S(x,t).$$
\end{lemma}


\begin{proof}
Notice that as $t$ varies $|X|_{c^*g}^2$ is not the norm of the
vector field $X$ with respect to the pullback of a fixed metric
$g(t)$. On the other hand, when we compute $\nabla_H X$ at a point
$(x,t)$ we are taking a covariant derivative with respect to the
pullback of a fixed metric $g(t)$ on the surface. Hence, in
computing $H(|X|_{c^*g}^2)$ the usual Leibniz rule does not apply.
In fact, there are two contributions to $H(|X|_{c^*g}^2)$: one, the
usual Leibniz rule differentiating in a frozen metric $g(t)$ and the
other coming from the effect on $|X|_{c^*g}^2$ of varying the metric
with $t$. Thus, we have
$$H(|X|_{c^*g}^2)(x,t)=-2\Ric_{c^*g(t)}(X(x,t),X(x,t))
+2\langle\nabla_HX(x,t),X(x,t)\rangle_{c^*g(t)}.$$ Since $t$ and $x$
are coordinates on the surface swept out by the family of curves,
$\nabla_HX=\nabla_XH$, and hence the second term on the right-hand
side of the previous equation can be rewritten as
$2\langle\nabla_XH(x,t),X(x,t)\rangle_{c^*g(t)}$. Since $X(x,t)$ and
$H(x,t)$ are orthogonal in $c^*g(t)$ and since $X=|X|_{c^*g}S$,
computing covariant derivatives in the metric $c^*g(t)$, we have
\begin{eqnarray*}
2\langle\nabla_XH,X\rangle_{c^*g(t)} & = & -2\langle H,\nabla_XX\rangle_{c^*g(t)}\\
 & = & -2\langle H,|X|_{g(t)}^2\nabla_SS\rangle_{c^*g(t)}-
 2\langle H,|X|_{g(t)}S(|X|_{c^*g})S\rangle_{c^*g(t)} \\
 & = & -2\langle H,H\rangle_{c^*g(t)}|X|_{g(t)}^2\\
 & = & -2k^2|X|_{c^*g}^2.
 \end{eqnarray*}
This proves the first inequality. As for the second, since $X$ and
$H$ commute we have
$$[H,S]=[H,|X|_{c^*g}^{-1}X]=H\left((|X|_{c^*g}^2)^{-1/2}\right)X=
\frac{-1}{2\left(|X|^2_g\right)^{3/2}}H(|X|_{c^*g})^2)X.$$ According to the
first equation, we can rewrite this as
$$[H,S](x,t)=\left(k^2+\Ric_{c^*g(t)}(S(x,t),S(x,t))\right)S(x,t).$$
\end{proof}


Now let us compute the time derivative of $k^2$. In what follows we drop the
dependence on the metric $c^*g(t)$ from all the curvature terms, but it is
implicitly there.

\begin{lemma}[Evolution equation for the curvature $k^2$]
\label{lem:curve-shrinking-curvature-evolution}
\label{2ndcurvshr}
\begin{eqnarray*}
\frac{\partial }{\partial t}k^2 & = & \frac{\partial ^2}{\partial
s^2}(k^2)-2\langle(\nabla_XH)^\perp,(\nabla_SH)^\perp\rangle_{c^*g}+2k^4\\
& & -2\Ric(H,H) +4k^2\Ric(S,S)+2\Rm(H,S,H,S),\end{eqnarray*} where the superscript $\perp$ means the
image under projection to the orthogonal complement of $X$.
\end{lemma}

\begin{proof}
\uses{lem:curve-shrinking-first-variation}
Using the same conventions as above for the function $|H|_{c^*g}$
and but leaving the metric implicit, we have
\begin{equation}\label{*1}
\frac{\partial }{\partial t}k^2=\frac{\partial }{\partial
t}(|H|^2_{c^*g})=-2\Ric(H,H)+2\langle\nabla_HH,H\rangle_{c^*g}.\end{equation}
 Now we compute
(using the second equation from Lemma~\ref{lem:curve-shrinking-first-variation})
\begin{eqnarray*}
\nabla_HH & = & \nabla_H\nabla_SS \\
& = & \nabla_S\nabla_HS+\nabla_{[H,S]}S+
{\mathcal R}(H,S)S \\
& = & \nabla_S\nabla_SH+\nabla_S([H,S])+\nabla_{[H,S]}S+{\mathcal R}(H,S)S \\
& = & \nabla_S\nabla_SH+\nabla_S\left((k^2+\Ric(S,S))S\right)+(k^2+\Ric(S,S))\nabla_SS+
{\mathcal R}(H,S)S \\
& = & \nabla_S\nabla_SH+2(k^2+\Ric(S,S))H+S(k^2+\Ric(S,S))S+ {\mathcal R}(H,S)S.
\end{eqnarray*}
Using this, and the fact that $\langle H,S\rangle_{c^*g}=0$, we have
\begin{equation}\label{*2}
2\langle\nabla_HH,H\rangle_{c^*g}=2g(\nabla_S\nabla_SH,H)+4k^4+4k^2\Ric(S,S))+2\Rm(H,S,H,S).\end{equation}
 On the other hand,
\begin{equation}\label{*3}
S(S(\langle
H,H\rangle_{c^*g}))=2\langle\nabla_S\nabla_SH,H\rangle_{c^*g}+2\langle\nabla_SH,\nabla_SH\rangle_{c^*g}.
\end{equation}

 We write
$$\nabla_SH=(\nabla_SH)^\perp+\langle\nabla_SH,S\rangle_{c^*g}S.$$
Since $H$ and $S$ are orthogonal, we have
$\langle\nabla_SH,S\rangle_{c^*g}=-\langle
H,\nabla_SS\rangle_{c^*g}=-\langle H,H\rangle$. Thus, we have
$$\nabla_SH=(\nabla_SH)^\perp-\langle H,H\rangle_{c^*g}S. $$
It follows that
$$-2\langle \nabla_SH,\nabla_SH\rangle_{c^*g}=
-2\langle(\nabla_SH)^\perp,(\nabla_SH)^\perp\rangle_{c^*g}-2k^4.$$ Substituting
this into Equation~(\ref{*3}) gives
\begin{equation}\label{*4}
2\langle\nabla_S\nabla_SH,H\rangle_{c^*g}=S(S(|H|_{c^*g}^2))
-2\langle(\nabla_SH)^\perp,(\nabla_SH)^\perp\rangle_{c^*g}-2k^4.
\end{equation}
Plugging this into Equation~(\ref{*2})  and using
Equation~(\ref{*1}) yields
\begin{eqnarray*}
\frac{\partial}{\partial t}k^2 & = & -2\Ric(H,H)+S(S\langle
H,H\rangle_{c^*g})-2\langle(\nabla_SH)^\perp,(\nabla_SH)^\perp\rangle_{c^*g}
\\
& & +2k^4+4k^2\Ric(S,S)+2\Rm(H,S,H,S).\end{eqnarray*} Of course,
$S(S(\langle H,H\rangle_{c^*g}))=(k^2)''$ so that this gives the result.
\end{proof}




Grouping together the last three terms in the statement of the
previous lemma, we can rewrite the result as
\begin{equation}\label{*4'}\frac{\partial }{\partial t}k^2\le
(k^2)''-2\langle(\nabla_SH)^\perp,(\nabla_SH)^\perp\rangle_{c^*g}+2k^4+\widehat
Ck^2,\end{equation} where the primes refer to the derivative with
respect to arc length along the curve and $\widehat C$ is a constant
depending only on an upper bound for the norm of the sectional
curvatures  of the ambient manifolds in the Ricci flow.

\begin{claim}[Differential inequality for the curvature $k$]
\label{claim:curve-shrinking-curvature-inequality}
\label{3rdcurvshr}
 There is a constant $C_1<\infty$ depending only on an
upper bound for the norm of the sectional curvatures of the ambient
manifolds in the Ricci flow $(M,g(t)),\ t_0\le t\le t_1$, such that
$$\frac{\partial}{\partial t} k\le k''+k^3+C_1k.$$
\end{claim}

\begin{proof}
We set $C_1=\widehat C/2$, where $\widehat C$ is as in
Inequality~\ref{*4'}. It follows from Inequality~(\ref{*4'}) that
\begin{equation}
\label{*5} 2k\frac{\partial k}{\partial t}\le
2kk''+2(k')^2+2k^4-2\langle(\nabla_SH)^\perp,(\nabla_SH)^\perp\rangle_{c^*g}+\hat
Ck^2.\end{equation} Since $k^2=\langle H,H\rangle_{c^*g}$, we see
that $(k^2)'=2\langle\nabla_SH,H\rangle_{c^*g}$. Since $H$ is
perpendicular to $S$, this can be rewritten as
$(k^2)'=2\langle(\nabla_SH)^\perp,H\rangle_{c^*g}$. It follows that
$$k'=\frac{\langle(\nabla_SH)^\perp,H\rangle_{c^*g}}{|H|_{c^*g}}.$$
Hence,
$$(k')^2\le \frac{\langle(\nabla_SH)^\perp,H\rangle_{c^*g}^2}{|H|_{c^*g}^2}\le
\langle(\nabla_SH)^\perp,(\nabla_SH)^\perp\rangle_{ c^*g}.$$ Plugging this into
Equation~(\ref{*5}) gives
$$\frac{\partial k}{\partial t}\le k''+k^3+C_1k.$$
\end{proof}

Now we define the {\em total length} of the curve $c(x,t)$,
$$L(t)=\int |X|_{c^*g}dx=\int ds.$$
 We also define the {\em total curvature} of the curve $c(x,t)$,
$$\Theta(t)=\int k|X|_{c^*g}dx=\int kds.$$


\begin{lemma}[Growth bounds for length and total curvature]
\label{lem:curve-shrinking-length-curvature-growth}
\label{1835}
There is a constant $C_2<\infty$ depending only on an
upper bound for the norm of the sectional curvatures of the ambient
manifolds in the Ricci flow such that
\begin{equation}\label{Lformula}
\frac{d}{dt}L\le \int (C_2-k^2)ds\end{equation} and
$$\frac{d}{dt}\Theta\le C_2\Theta.$$
\end{lemma}

\begin{proof}
\uses{lem:curve-shrinking-first-variation,claim:curve-shrinking-curvature-inequality}
$$\frac{d}{dt}L=\int\frac{\partial}{\partial t}\sqrt{|X|_{c^*g}^2}dx.$$
By Lemma~\ref{lem:curve-shrinking-first-variation} we have $$\frac{d}{dt}L=\int
\frac{1}{2|X|_{c^*g}}\left(-2\Ric(X,X)-2k^2|X|_{c^*g}^2\right)dx.$$ Thus,
\begin{equation}\label{dLdteqn}\frac{d}{dt}L=\int (-\Ric(S,S)-k^2)|X|_{c^*g}dx=\int(-\Ric(S,S)-k^2)ds.\end{equation} The first inequality in the lemma
then follows by taking $C_2$ to be an upper bound for the norm of
$\Ric_{g(t)}$.

Now let us consider the second inequality in the statement.
$$\frac{d}{dt}\Theta=\int\frac{\partial}{\partial t}(k|X|_{c^*g})dx=\int\left(\frac{\partial k}{\partial
t}|X|_{c^*g} +k\frac{\partial |X|_{c^*g}}{\partial t}\right)dx.$$ Thus, using
Claim~\ref{claim:curve-shrinking-curvature-inequality} and the first equation in Lemma~\ref{lem:curve-shrinking-first-variation} we have
\begin{eqnarray*}
\frac{d}{dt}\Theta & \le &  \int (k''+k^3+C_1k)ds+\int
\frac{k}{2|X|_{c^*g}}(-2\Ric(X,X)-2k^2|X|_{c^*g}^2)dx \\
& = & \int(k''+k^3+C_1k)ds-\int k(\Ric(S,S)+k^2)ds \\& = &
\int(k''+C_1k-k\,\Ric(S,S))ds.
\end{eqnarray*}
Since $\int k''ds=0$ by the fundamental theorem of calculus, we get
$$\frac{d}{dt}\Theta\le  C_2\Theta,
$$
 for an
appropriate constant $C_2$ depending only on an upper bound for the
norm of the sectional curvatures of the ambient family $(M,g(t))$.
\end{proof}

\begin{corollary}[Exponential bound for length and total curvature]
\label{cor:curve-shrinking-length-curvature-exponential-bound}
\label{lncurvgrowth}
The following holds for the constant $C_2$ as in the previous
lemma. Let $c(x,t)$ be a curve-shrinking flow, let $L(t)$ be the
total length of $c(t)$ and let $\Theta(t)$ be the total curvature of
$c(t)$. Then for any $t_0\le t'<t''\le t_1$ we have $$L(t'')\le
L(t')e^{C_2(t''-t')}$$
$$\Theta(t'')\le \Theta(t')e^{C_2(t''-t')}.$$
\end{corollary}




\subsection{Ramp solutions in $ M\times S^1$}

As we pointed out in the beginning of Section~\ref{sect:csf} the
main obstacle we must overcome is that the curve-shrinking flow does
not always exist for the entire time interval $[t_0,t_1]$. The
reason is the following: Even though, as we shall see, it is
possible to bound the total curvature of the curve-shrinking flow in
terms of the total curvature of the initial curve and the ambient
Ricci flow, there is no pointwise estimate on the curvature for the
curve-shrinking flow. The idea for dealing with this problem, which
goes back to \cite{AG}, is to replace the original situation of
curves in a manifold with graphs by taking the product of the
manifold with a circle and using ramps. We shall see that in this
context the curve-shrinking flow always exists. The problem then
becomes to transfer the information back from the flows of ramps to
the original manifold.


Now suppose that the Ricci flow is of the form $(M, g(t))\times
(S^1_\lambda,ds^2)$ where $(S^1_\lambda,ds^2)$ denotes the  circle
of length $\lambda$. Notice that the sectional curvatures of this
product flow depend only on the sectional curvatures of $(M, g(t))$
and, in particular, are independent of $\lambda$. Let $U$ denote
vector field made up of unit tangent vectors in the direction of the
circle factors. Let $u(x,t)=\langle S,U\rangle_{g(t)}$.

\begin{claim}[Evolution equation for the ramp function $u$]
\label{claim:ramp-function-evolution}
\label{rampcomp}
$$\frac{\partial u}{\partial t}=u''+(k^2+\Ric(S,S))u\ge u''-C'u,$$
where $C'$ is an upper bound for the norm of the Ricci curvature of
$(M,g(t))$.
\end{claim}


\begin{proof}
\uses{lem:curve-shrinking-first-variation}
Since $U$ is a constant vector field and hence parallel along all curves and
since $\Ric(V,U)=0$ for all tangent vectors $V$, by Lemma~\ref{lem:curve-shrinking-first-variation}
we have
\begin{eqnarray*}
\frac{\partial}{\partial t}\langle S,U\rangle_{g(t)} & = & -2\Ric(S,U)
+\langle dc(\nabla_HS),U\rangle_{g(t)} \\
& = & \langle dc(\nabla_HS),U\rangle_{g(t)}
=  \langle dc([H,S]+\nabla_SH),U\rangle_{g(t)} \\
& = & (k^2+\Ric(S,S))u+\langle dc(\nabla_SH),U\rangle_{g(t)} \\
& = & (k^2+\Ric(S,S))u+S(dc(\langle H),U\rangle_{g}) \\
 & = & (k^2+\Ric(S,S))u+S(\langle dc(\nabla_SS),U\rangle_g) \\
&  = & (k^2+\Ric(S,S))u+S(S(u))  =  (k^2+\Ric(S,S))u+u''.
 \end{eqnarray*}
\end{proof}

\begin{definition}[Ramp]
\label{def:ramp-curve}
A curve $c\colon S^1\to  M\times S^1_\lambda$ is said to be a {\em
ramp} if $u$ is strictly positive.
\end{definition}

The main results of this section show that the curve-shrinking flow
is much better behaved for ramps  than for the general smooth curve.
First of all, as the next corollary shows, the curve-shrinking flow
applied to a ramp produces a one-parameter families of ramps. The
main result of this section shows that for any ramp as initial
curve, the curve-shrinking flow does not develop singularities as
long as the ambient Ricci flow does not.


\begin{corollary}[Ramps remain ramps under curve-shrinking flow]
\label{cor:ramp-preserved-under-flow}
\label{ubd}
If $c(x,t),\ t_0\le t< t'_1<\infty$, is a solution of the curve
shrinking flow in $( M, g(t))\times (S^1_\lambda,ds^2)$ and if
$c(t_0)$ a ramp, then $c(t)$ is a ramp for all $t\in [t_0,t'_1)$.
\end{corollary}

\begin{proof}
\uses{claim:ramp-function-evolution}
From  the equation in Claim~\ref{claim:ramp-function-evolution}, we see that for  $C'$ an
upper bound for the norm of the Ricci curvature, we have
$$\frac{\partial}{\partial t}\left(e^{C't}u\right)\ge \left(e^{C't}u\right)''.$$
It now follows from a standard maximum principle argument that the
minimum value of $e^{C't}u$ is a non-decreasing function of $t$.
Hence, if $c(t_0)$ is a ramp then each $c(t)$ is a ramp and in fact
$u(x,t)$ is uniformly bounded away from zero in terms of the minimum
of $u(x,t_0)$ and the total elapsed time $t_1-t_0$.
\end{proof}

\begin{lemma}[Curve-shrinking flow exists for all time for ramps]
\label{lem:curve-shrinking-exists-for-ramps}
\label{csforramps}
Let $(M,g(t)),\ t_0\le t\le t_1$, be a Ricci flow. Suppose that
$c\colon S^1\to ( M\times S^1_\lambda, g(t)\times ds^2)$ is a ramp.
Then there is a curve-shrinking flow $c(x,t)$ defined for all $t\in
[t_0,t_1]$ with $c$ as the initial condition at time $t=t_0$. The
curves $c(\cdot,t)$ are all ramps.
\end{lemma}


\begin{proof}
\uses{claim:curve-shrinking-short-time-existence,claim:curve-shrinking-curvature-inequality,claim:ramp-function-evolution,cor:ramp-preserved-under-flow}
The real issue here is to show that the curve-shrinking flow exists
for all $t\in [t_0,t_1]$. Given this, the second part of the
statement follows from the previous corollary. If the curve
shrinking flow does not exist on all of $[t_0,t_1]$ then by
Claim~\ref{claim:curve-shrinking-short-time-existence} there is a $t'_1\le t_1$ such that the
curve-shrinking flow exists on $[t_0,t_1')$ but $k$ is unbounded on
$S^1\times [t_0,t'_1)$. Thus, to complete the proof we need to see
that for any $t'_1$ for which the curve-shrinking flow is defined on
$[t_0,t'_1)$ we have a uniform bound on $k$ on this region.

Using Claim~\ref{claim:curve-shrinking-curvature-inequality} and Claim~\ref{claim:ramp-function-evolution} we compute
\begin{eqnarray*} \frac{\partial}{\partial t}\left(\frac{k}{u}\right)
& = &\frac{1}{u}\frac{\partial
k}{\partial t}-\frac{k}{u^2}\frac{\partial u}{\partial t} \\
& \le & \frac{k''+k^3+C_1k}{u}-\frac{k}{u^2}\left(u''+(k^2+\Ric(S,S))u\right)
\\
& = & \frac{k''}{u}-\frac{ku''}{u^2}+\frac{C_1k}{u}-\frac{k}{u}\Ric(S,S).
\end{eqnarray*}
On the other hand,
$$\left(\frac{k}{u}\right)''=\frac{k''u-u''k}{u^2}-2\left(\frac{u'}{u}\right)
\left(\frac{k'u-u'k}{u^2}\right).$$ Plugging this in, and using the
curvature bound on the ambient manifolds we get
$$\frac{\partial}{\partial
t}\left(\frac{k}{u}\right)\le
\left(\frac{k}{u}\right)''+\left(\frac{2u'}{u}\right)\left(\frac{k}{u}\right)'+C'\frac{k}{u},$$
for a constant $C'$ depending only on a bound for the norm of the
sectional curvature of the ambient Ricci flow. A standard maximum
principle argument shows that the maximum of $k/u$ at time $t$ grows
at most exponentially rapidly in $t$. Since $u$ stays bounded away
from zero, this implies that for ramp solutions on a finite time
interval, the value of $k$ is bounded.
\end{proof}




Next let us turn to the growth rate of the area of a minimal annulus
connecting two ramp solutions.

\begin{lemma}[Growth rate of the minimal annulus between ramps]
\label{lem:minimal-annulus-area-growth}
\label{mugrowth}
Suppose that the dimension $n$ of $ M$
is at least three. Let $c_1(x,t)$ and $c_2(x,t)$ be ramp solutions
in $( M, g)\times (S^1_\lambda,ds^2)$ with the image under the
projection to $S^1_\lambda$ of each $c_i$ being of degree one. Let
$\mu(t)$ be the infimum of the areas of annuli in $( M\times
S^1_\lambda, g(t)\times ds^2)$ with boundary $c_1(x,t)\cup
c_2(x,t)$. Then $\mu(t)$ is a continuous function of $t$ and
$$\frac{d}{dt}\mu(t)\le (2n-1)\max_{x\in M}|\Rm(x,t)|\mu(t),$$
in the sense of forward difference quotients.
\end{lemma}


\begin{proof}
Fix a time $t'$. First assume that the loops $c_1(\cdot,t')$ and
$c_2(\cdot,t')$ are disjoint. Under Ricci flow the metrics on the
manifold immediately become real analytic (see \cite{Bando}) and
furthermore, under the curve-shrinking flow the curves $c_1$ and
$c_2$ immediately become analytic (see \cite{GageHamilton}).
[Neither of these results is essential for this argument because we
could approximate both the metric and the curves by real analytic
objects.] Establishing the results for these and taking limits would
give the result in general. Since $c_1(\cdot,t')$ and
$c_2(\cdot,t')$ are homotopic and are homotopically non-trivial
there is an annulus connecting them and there is a positive lower
bound to the length of any simple closed curve in any such annulus
homotopic to a boundary component. Hence, there is a minimal annulus
spanning $c_1(\cdot,t')\coprod c_2(\cdot, t')$ According to results
of Hildebrandt (\cite{Hildebrandt}) and Morrey (\cite{Morrey}) any
minimal annulus $A$ with boundary the union of these two curves is
real analytic up to and including the boundary and is immersed
except for finitely many branch points. By shifting the boundary
curves slightly within the annulus, we can assume that there are no
boundary branch points. Again, if we can prove the result for these
perturbed curves taking limits will give the result for the original
ones.  Given the deformation vector $H$ on the boundary of the
annulus, extend it to a deformation vector $\hat H$ on the entire
annulus. The first order variation of the area at time $t'$ of the
resulting deformed family of annuli is given by
$$\frac{d \Area\,A}{d t}(t')=\int_{A}(-\operatorname{Tr}(\Ric^T(g(t'))))da
+\int_{\partial A}-k_\mathrm{geod}ds,$$ where $\Ric^T$ is the
Ricci curvature in the tangent directions to the annulus. (The first
term is the change in the area of the fixed annulus as the metric
deforms. The second term is the change in the area of the family of
annuli in the fixed metric. There is no contribution from moving the
annulus in the normal direction since the original annulus is
minimal.)
 If $A$ is embedded, then by the Gauss-Bonnet theorem, we have
$$\int_{\partial A}-k_\mathrm{geod}ds =\int_{A}Kda$$
where $K$ is the Gaussian curvature of $A$. More generally, if $A$
has interior branch points of orders $n_1,\ldots,n_k$ then there is
a correction term and the formula is
$$\int_{\partial A}-k_\mathrm{geod}ds=\int_{A}Kda-\sum_{i=1}^k2\pi(n_i-1).$$ Thus, we see
$$\frac{d \Area\, A}{d t}(t')\le\int_{A}(-\operatorname{Tr}(\Ric^T(g(t')))+K)da.$$
On the other hand, since $A$ is a minimal surface,  $K$ is at most
the sectional curvature of $( M, g(t'))\times (S^1_\lambda,ds^2)$
along the two-plane tangent to the annulus. Of course, the trace of
the Ricci curvature along $A$ is at most $2(n-1)|\max_{x\in
M}\Rm(x,t')|$. Hence,
$$\frac{d\Area\, A}{dt}(t')\le (2n-1)|\max_{x\in M}\Rm(x,t')|\mu(t').$$

This computation was done assuming that $c_2(\cdot,t')$ is disjoint
from $c_1(\cdot,t')$. In general, since the dimension of $ M$ is at
least three, given $c_1(\cdot,t')$ and $c_2(\cdot,t')$ we can find
$c_3(\cdot,t)$ arbitrarily close to $c_2(\cdot,t')$ in the
$C^2$-sense and disjoint from both $c_1(\cdot,t')$ and
$c_2(\cdot,t')$. Let $A_3$ be a minimal annulus connecting
$c_1(\cdot,t')$ to $c_3(\cdot,t')$ and $A_2$ be an minimal annulus
connecting $c_3(\cdot,t')$ to $c_2(\cdot,t')$. We apply the above
argument to these annuli to estimate the growth rate of minimal
annuli connecting the corresponding curve-shrinking flows. Of course
the sum of these areas (as a function of $t$) is an upper bound for
the area of a minimal annulus connecting the curve-shrinking flows
starting from $c_1(\cdot,t')$ and $c_2(\cdot,t')$. As we choose
$c_3(\cdot,t')$ closer and closer to $c_2(\cdot,t')$, the area of
$A_2$ tends to zero and the area of $A_3$ tends to the area of a
minimal annulus connecting $c_1(\cdot,t')$ and $c_2(\cdot,t')$. This
establishes the continuity of $\mu(t)$ at $t'$ and also establishes
the forward difference quotient estimate in the general case.
\end{proof}

\begin{corollary}[Exponential growth bound for the minimal annulus area]
\label{cor:minimal-annulus-exponential-growth}
\uses{prop:curve-deformation-existence}
Given  curve-shrinking flows $c_1(\cdot,t)$ and $c_2(\cdot,t)$ for
ramps of degree one in $(\bar M,\bar g(t))\times (S^1_\lambda,ds^2)$
the minimal area of an annulus connecting $c_1(\cdot,t)$ and
$c_2(\cdot,t)$ grows at most exponentially with time with an
exponent determined by an upper bound on the sectional curvature of
the ambient flow, which in particular is independent of $\lambda$.
\end{corollary}



\section{Proof of Proposition~{\ref{prop:curve-deformation-existence}}}

Now we are ready to use the curve-shrinking flow for ramps in
$M\times S_\lambda^1$ to establish Proposition~\ref{prop:curve-deformation-existence} for $M$. As
we indicated above, the reason for replacing the flow $(M,g(t))$
that we are studying with its product with $S^1_\lambda$ and
studying ramps in the product is that the curve-shrinking flow
exists for all time $t\in [t_0,t_1]$ for these. By this mechanism we
avoid the difficulty of finite time singularities in the curve
shrinking flow. On the other hand, we have to translate results for
the ramps back to results for the original Ricci flow $(M,g(t))$.
This requires careful analysis.


\subsection{Approximations to the original family}

The first step in the proof of Proposition~\ref{prop:curve-deformation-existence} is to identify
the approximation to the family $\Gamma$ that we shall use. Here is
the lemma that gives the needed approximation together with all the
properties we shall use.

Given a loop $c$ in $M$ and $\lambda>0$ we define a loop $c^\lambda$
in $M\times S^1_\lambda$. The loop $c^\lambda$ is obtained by
setting $c^\lambda(x)=(c(x),x)$ where we use  a standard
identification of the domain circle (the unit circle) for the free
loop space with $S^1_\lambda$, an identification that defines a loop
in $S^1_\lambda$ of constant speed $\lambda/2\pi$.


\begin{lemma}[Good C^2-approximation of a family of loops]
\label{lem:good-approximation-family}
\label{goodapprox}
Given a continuous map $\Gamma\colon S^2\to \Lambda M$ representing
an element of $\pi_3( M,*)$ and $0<\zeta<1$, there is a continuous
map $\widetilde \Gamma\colon S^2\to \Lambda M$ with the following
properties:
\begin{enumerate}
\item[(1)] $[\widetilde \Gamma]=[\Gamma]$ in $\pi_3(M,*)$.
\item[(2)] For each $c\in S^2$ the loop $\widetilde \Gamma(c)$ is a
$C^2$-loop.
\item[(3)] For each $c\in S^2$ the length of $\widetilde \Gamma(c)$ is
within $\zeta$ of the length of $\Gamma(c)$.
\item[(4)] For each $c\in S^2$, we have $|A(\widetilde
\Gamma(c))-A(\Gamma(c))|<\zeta$.
\item[(5)] There is a constant $C_0<\infty$ depending only on $\Gamma$,
on the bounds for the norm of the Riemann curvature operator of the
ambient Ricci flow, and on $\zeta$ such that for each $c\in S^2$ and
each $\lambda\in (0,1)$ the total length and the total curvature of
the ramp $\widetilde \Gamma(c)^\lambda$ are both bounded by $C_3$.
\end{enumerate}
\end{lemma}


Before proving this lemma we need some preliminary definitions and
constructions.

\begin{definition}[Regular n-polygonal approximation]
\label{def:n-polygonal-approximation}
Let $c\colon S^1\to M$ be a $C^1$-map.  Fix a positive integer $n$.
By a {\em regular $n$-polygonal approximation} to $c$ we mean the
following. Let $\xi_n=\exp(2\pi i/n)$, and consider the points
$p_k=c(\xi^k_n)$ for $k=1,\ldots,n+1$. For each $1\le k\le n$, let
$A_k$ be a minimal geodesic in $M$ from $p_k$ to $p_{k+1}$. We
parameterize $A_k$ by the interval $[\xi_n^k,\xi_n^{k+1}]$ in the
circle at constant speed. This gives a  piecewise geodesic map
$c_n\colon S^1\to M$.
\end{definition}

The following is immediate from the definition.

\begin{claim}[Basic properties of n-polygonal approximations]
\label{claim:polygonal-approximation-properties}
\label{onec}
Given $\zeta>0$ and a $C^1$-map $c\colon S^1\to M$ then for all $n$
sufficiently large the following hold for the $n$-polygonal
approximation $c_n$ of $c$.
\begin{enumerate}
\item[(a)] the length of $c_n$ is within $\zeta$ of the length of $c$.
\item[(b)] there is a map of the annulus $S^1\times I$ to $M$ connecting
$c_n$ to $c$ with the property that the image is piecewise smooth
and of area less than $\zeta$.
\end{enumerate}
\end{claim}

\begin{proof}
The length of $c$ is the limit of the lengths of the $n$-polygonal
approximations as $n$ goes to infinity. The first item is immediate
from this. As to the second, for $n$ sufficiently large, the
distance between the maps $c$ and $c_n$ will be arbitrarily small in
the $C^0$-topology, and in particular will be much smaller than the
injectivity radius of $M$. Thus, for each $k$ we can connect $A_k$
to the corresponding part of $c$ by a family of short geodesics.
Together, these form an annulus, and it is clear that for $n$
sufficiently large the area of this annulus is arbitrarily small.
\end{proof}

As the next result shows, for $\zeta>0$, the integer $n(c)$
associated by the previous claim to a $C^1$-map $c$ can be made
uniform as $c$ varies over a compact subset of $\Lambda M$.

\begin{claim}[Uniformity of the polygonal approximation over a compact family]
\label{claim:polygonal-approximation-uniform}
\label{uniformc}
\uses{claim:polygonal-approximation-properties}
Let $X\subset \Lambda M$ be a compact subset and let $\zeta>0$ be
fixed. Then there is $N$ depending only on $X$ and $\zeta$ such the
conclusion of the previous claim holds for every $c\in X$ and every
$n\ge N$.
\end{claim}

\begin{proof}
\uses{claim:polygonal-approximation-properties}
Suppose the result is false. Then for each $N$ there is $c_N\in X$
and $n\ge N$ so that the lemma does not hold for $c_N$ and $n$.
Passing to a subsequence, we can suppose that the $c_N$ converge to
$c_\infty\in X$. Applying Claim~\ref{claim:polygonal-approximation-properties} we see that there is $N$
such that the conclusion of Claim~\ref{claim:polygonal-approximation-properties} holds with $\zeta$
replaced by $\zeta/2$ for $c_\infty$ and all $n\ge N$. Clearly, then
by continuity for all $n\ge N$ the conclusion of Claim~\ref{claim:polygonal-approximation-properties}
holds for the $n$-polygonal approximation for every $c_l$ for all
$l$ sufficiently large. This is a contradiction.
\end{proof}


\begin{corollary}[Polygonal approximation of a family $\Gamma$]
\label{cor:polygonal-approximation-of-family}
\label{Gammaapprox}
Let $\Gamma\colon S^2\to \Lambda M$ be a continuous map  with the
property that $\Gamma(c)$ is homotopically trivial for all $c\in
S^2$. Fix $\zeta>0$. For any $n$ sufficiently large denote by
$\Gamma_n$ the family of loops defined by setting $\Gamma_n(c)$
equal to the  $n$-polygonal approximation to $\Gamma(c)$. There is
$N$ such that for all $n\ge N$ we have
\begin{enumerate}
\item[(1)] $\Gamma_n$ is a continuous family of $n$-polygonal loops in
$M$.
\item[(2)] For each $c\in S^2$, the loop
$\Gamma_n(c)$ is a homotopically trivial loop in $M$ and its length
is within $\zeta$ of the length of $\Gamma(c)$.
\item[(3)] For each $c\in S^2$, we have $|A(\Gamma_n(c))-A(\Gamma(c))|<\zeta$.
\end{enumerate}
\end{corollary}

\begin{proof}
\uses{claim:polygonal-approximation-uniform}
Given $\Gamma$ there is a uniform bound over all $c\in S^2$ on the
maximal speed of $\Gamma(c)$. Hence, for all $n$ sufficiently large,
the lengths of the sides in the $n$-polygonal approximation to
$\Gamma(c)$ will be uniformly small. Once this length is less than
the injectivity radius of $M$, the minimal geodesics between the
endpoints are unique and vary continuously with the endpoints. This
implies that for $n$ sufficiently large the family $\Gamma_n$ is
uniquely determined and itself forms a continuous family of loops in
$M$. This proves the first item. We have already seen that, for $n$
sufficiently large, for all $c\in S^2$ there is an annulus
connecting $\Gamma(c)$ and $\Gamma_n(c)$. Hence, these loops are
homotopic in $M$. The first statement in the second item follows
immediately. The last statement in the second item and third item
follow immediately from Claim~\ref{claim:polygonal-approximation-uniform}.
\end{proof}

The next step is  to turn these $n$-polygonal approximations into
$C^2$-curves. We fix, once and for all, a $C^\infty$ function
$\psi_n$ from the unit circle to $[0,\infty]$ with the following
properties:
\begin{enumerate}
\item[(1)] $\psi_n$ is non-negative and vanishes to infinite order at the point $1$ on the unit circle.
\item[(2)] $\psi_n$ is periodic with period $2\pi i/n$.
 \item[(3)] $\psi_n$ is positive on the interior of the interval $[1,\xi_n]$ on the unit circle,
and the restriction of $\psi_n$ to this interval is symmetric about
$\exp(\pi i/n)$, and is increasing from $1$ to $\exp(\pi
i/n)$.
\item[(4)] $\int_1^{\xi_n}\psi_n(s)ds =2\pi/n$.
\end{enumerate}

 Now we define a map $\tilde\psi_n\colon S^1\to S^1$ by
$$\tilde \psi_n(x)= \int_{1}^x\psi_n(y)dy.$$ It is easy to see that
the conditions on $\psi_n$ imply that this defines a $C^\infty$-map
from $S^1$ to $S^1$ which is a homeomorphism and is a diffeomorphism
on the complement of the $n^{th}$ roots of unity.

Now given an $n$-polygonal loop $c_n$ we define the smoothing
$\tilde c_n$ of $c_n$ by $\tilde c_n=c_n\circ \tilde \psi_n$. This
smoothing $\tilde c_n$ is a $C^\infty$-loop in $M$ with the same
length as the original polygonal loop $c_n$. Notice that the
curvature of $\tilde c_n$ is not itself a continuous function: just
like the polygonal map it replaces, it has a $\delta$-function at
the `corners' of $c_n$.


\begin{proof}
\uses{cor:polygonal-approximation-of-family}
 (of Lemma~\ref{lem:good-approximation-family})
Given a continuous map $\Gamma\colon S^2\to \Lambda M$ and $\zeta>0$
we fix $n$ sufficiently large so that Corollary~\ref{cor:polygonal-approximation-of-family}
holds for these choices of $\Gamma$ and $\zeta$. Let $\widetilde
\Gamma=\widetilde \Gamma_n$ be the family of smoothings of the
family $\Gamma_n$ of $n$-polygonal loops. Since this smoothing
operation changes neither the length nor the area of a minimal
spanning disk, it follows  immediately from the construction and
Corollary~\ref{cor:polygonal-approximation-of-family} that $\widetilde \Gamma$ satisfies the
conclusions of Lemma~\ref{lem:good-approximation-family} except possibly the last one.


To establish the last conclusion we must examine the lengths and
total curvatures of the ramps $\widetilde \Gamma(c)^\lambda$
associated to this family of $C^2$-loops. Fix $\lambda$ with
$0<\lambda<1$, and consider the product Ricci flow $(M,g(t))\times
(S^1_\lambda,ds^2)$ where the metric on $S^1_\lambda$ has length
$\lambda$.

\begin{claim}[Length and curvature bounds for the ramp of the smoothed approximation]
\label{claim:smoothed-approximation-ramp-bounds}
For any $0<\lambda<1$, the length of the ramp $\widetilde
\Gamma(c)^\lambda$ is at most $\lambda$ plus the length of
$\Gamma(c)$. The total curvature of $\widetilde\Gamma(c)^\lambda$ is
at most $n\pi$.
\end{claim}

\begin{proof}
The arc length element for $\widetilde \Gamma(c)^\lambda$ is
$\sqrt{a(x)^2+(\lambda/2\pi)^2}dx\le (a(x)+\lambda/2\pi)dx$ where
$a(x)dx$ is the arc length element for $\widetilde\Gamma(c)$.
Integrating gives the length estimate.

The total curvature of $\widetilde \Gamma(c)^\lambda$ is the sum
over the intervals $[\xi_n^k,\xi_n^{k+1}]$ of the total curvature on
these intervals. On any one of these intervals we have a curve in a
totally geodesic, flat surface: the curve lies in the product of a
geodesic arc in $M$ times $S^1_\lambda$. Let $u$ and $v$ be unit
tangent vectors to this surface, $u$ along the geodesic (in the
direction of increasing $x$) and $v$ along the $S^1_\lambda$ factor.
These are parallel vector fields on the flat surface.
 The tangent
vector $X(x)$ to the restriction of $\widetilde \Gamma(c)^\lambda$
to this interval is $L\psi_n(x) u+(\lambda/2\pi)v$, where $L$ is the
length of the geodesic segment we are considering. Consider the
first-half subinterval $[\xi_n^k,\xi_n^k\cdot\exp(\pi i/n)]$.
The tangent vector $X(x)$ is $(\lambda/2\pi) v$ at the initial point
of this subinterval and is $L\psi_n(\xi_n^k\cdot\exp(\pi
i/2))u+(\lambda/2\pi) v$ at the final point. Throughout this
interval the vector is of the form $a(x)u+(\lambda/2\pi) v$ where
$a(x)$ is an increasing function of $x$. Hence, the tangent vector
is always turning in the same direction and always lies in the first
quadrant (using $u$ and $v$ as the coordinates). Consequently, the
total turning (the integral of $k$ against arc-length) over this
interval is the absolute value of the difference of the angles at
the endpoints. This difference is less than $\pi/2$ and tends to
$\pi/2$ as $\lambda$ tends to zero, unless $L=0$ in which case there
is zero turning for any $\lambda>0$. By symmetry, the total turning
on the second-half subinterval $[\xi_n^k\cdot\exp(\pi
i/n),\xi_n^{k+1}]$ is also bounded above by $\pi/2$. Thus, for any
$\lambda>0$, the total turning on one of the segments is bounded
above by $\pi$. Since there are $n$ segments this gives the upper
bound of $n\pi$ on the total turning of $\widetilde
\Gamma(c)^\lambda$ as required.
\end{proof}

This claim completes the proof of the last property required of
$\widetilde\Gamma=\widetilde \Gamma_n$ and hence completes the proof
of Lemma~\ref{lem:good-approximation-family}.
\end{proof}

Having fixed $\Gamma$ and $\zeta>0$, we fix $n$ and set
$\widetilde\Gamma=\widetilde \Gamma_n$. We choose $n$ sufficiently
large so that $\widetilde \Gamma$ satisfies Lemma~\ref{lem:good-approximation-family}.
 Fix $\lambda\in (0,1)$ and
define $\widetilde \Gamma^\lambda\colon S^2\to (M\times S^1_\lambda)$, by
setting $\widetilde \Gamma^\lambda(c)=\widetilde \Gamma(c)^\lambda$.

Fix  $c\in S^2$, and let $\widetilde\Gamma_c^\lambda(t),\ t_0\le
t\le t_1$, be the curve-shrinking flow given in
Lemma~\ref{lem:curve-shrinking-exists-for-ramps} with initial data the ramp
$\widetilde\Gamma^\lambda(c)$ . As $c$ varies over $S^2$ these fit
together to produce a one-parameter family
$\widetilde\Gamma^\lambda(t)$ of maps $S^2\to\Lambda(M\times
S^1_\lambda)$. Let $p_1$ denote the projection of $M\times
S^1_\lambda$ to $M$. Notice that for any $\lambda$ we have
$\widetilde\Gamma^\lambda_c(t_0)=\widetilde\Gamma^\lambda(c)$, so
that $p_1\widetilde \Gamma_c^\lambda(t_0)=\widetilde \Gamma(c)$. We
shall show that for $\lambda>0$ sufficiently small, the family
$p_1\widetilde\Gamma^\lambda(t)$ satisfies the conclusion of
Proposition~\ref{prop:curve-deformation-existence} for the fixed $\Gamma$ and $\zeta>0$. We do
this in steps. First, we show that fixing one $c\in S^2$, for
$\lambda$ sufficiently small (depending on $c$) an analogue  of
Proposition~\ref{prop:curve-deformation-existence} holds for the one-parameter family of loops
$p_1\widetilde\Gamma_c^\lambda(t)$.  By this we mean that either
$p_1\widetilde\Gamma^\lambda_c(t_1)$ has length less than $\zeta$ or
$A(p_1\widetilde\Gamma^\lambda_c(t_1))$ is at most the value
$v(t_1)+\zeta$, where $v$ is the solution to the
Equation~(\ref{wdiffeqn}) with initial condition
$v(t_0)=A(\widetilde\Gamma(c))$. (Actually, we establish a slightly
stronger result, see Lemma~\ref{lem:curve-deformation-single-point-estimate}.) The next step in the
argument is to take a finite subset ${\mathcal S}\subset S^2$ so
that for every $c\in S^2$ there is $\hat c\in {\mathcal S}$ such
that $\widetilde\Gamma(c)$ and $\widetilde\Gamma(\hat c)$ are
sufficiently close. Then, using the result of a single $c$, we fix
$\lambda>0$ sufficiently small so that the analogue of
Proposition~\ref{prop:curve-deformation-existence} for individual curves (or rather the slightly
stronger version of it) holds for every $\hat c\in {\mathcal S}$.
Then we complete the proof of Proposition~\ref{prop:curve-deformation-existence} using the fact
that for every $c$ the curve $\widetilde\Gamma(c)$ is sufficiently
close to a curve $\widetilde\Gamma(\hat c)$ associated to an element
$\hat c\in {\mathcal S}$.

\subsection{The case of a single $c\in S^2$}



According to Lemma~\ref{lem:good-approximation-family}, for all $\lambda\in (0,1)$  the
lengths and total curvatures of the $\widetilde\Gamma^\lambda(c)$
are uniformly bounded for all $c\in S^2$. Hence, by
Corollary~\ref{cor:curve-shrinking-length-curvature-exponential-bound} the same is true for
$\widetilde\Gamma_c^\lambda(t)$ for all $c\in S^2$ and all $t\in
[t_0,t_1]$.

\begin{claim}[Linear bound on the forward difference quotient of area]
\label{claim:area-forward-difference-linear-bound}
\label{areachange}
There is a constant $C_4$ depending on $t_1-t_0$, on the curvature
bound of the sectional curvature of the Ricci flow $(M,g(t)),\
t_0\le t\le t_1$, on the original family $\Gamma$ and on $\zeta$
such that for any $c\in S^2$ and any $t_0\le t'<t''\le t_1$ we have
$$A(p_1\widetilde \Gamma^\lambda_c(t''))-A(p_1\widetilde
\Gamma^\lambda_c(t'))\le C_4(t''-t').$$
\end{claim}

\begin{proof}
\uses{lem:curve-shrinking-length-curvature-growth}
All the constants in this argument are allowed to depend on
$t_1-t_0$, on the curvature bound of the sectional curvature of the
Ricci flow $(M,g(t)),\ t_0\le t\le t_1$, on the original family
$\Gamma$ and on $\zeta$ but are independent of $\lambda$, $c\in
S^2$, and $t'<t''$ with $t_0\le t'$ and $t''\le t_1$. First, let us
consider the surface $S^\lambda_c[t',t'']$ in  $M\times S^1_\lambda$
swept out by $c(x,t),\ t'\le t\le t''$. We denote by $\Area\,(S^\lambda_c[t',t''])$ the area of this surface with respect
to the metric $g(t'')\times ds^2$. We compute the derivative of this
area for fixed $t'$ as $t''$ varies. There are two contributions to
this derivative: (i) the contribution due to the variation of the
metric $g(t'')$ with $t''$ and (ii) the contribution due to
enlarging the surface. The first is
$$\int_{S^\lambda_c[t',t'']}-\Tr\, \Ric^Tda$$
 where $\Ric^T$ is the restriction of the
Ricci tensor of the ambient metric $g(t'')$ to the tangent planes to
the surface and $da$ is the area form of the surface in the metric
$g(t'')\times ds^2$. The second contribution is
$\int_{c(x,t'')}|H|ds$. According to Lemma~\ref{lem:curve-shrinking-length-curvature-growth} there is a
constant $C'$ (depending only on the curvature bound for the
manifold flow, the initial family $\Gamma(t)$ and $\zeta$ and
$t_1-t_0$) such that the second term is bounded above by $C'$. The
first term is bounded above by $C''\Area\,S^\lambda_c[t',t'']$
where $C''$ depends only on the bound on the sectional curvatures of
the ambient Ricci flow. Integrating we see that there is a constant
$C_1'$ such that the derivative of the area function is at most
$C_1'$. Since its value at $t'$ is zero, we see that $\Area\,S^\lambda_c[t',t'']\le C'_1(t''-t')$. It follows that the
area of $p_1S^\lambda_c[t',t'']$ with respect to the metric $g(t'')$
is at most $C_1'(t''-t')$.

Now we compute an upper bound for the forward difference quotient of
$A(p_1\widetilde \Gamma^\lambda_c(t))$ at $t=t'$. For any $t''>t'$
we have a spanning disk for $p_1\widetilde\Gamma^\lambda_c(t'')$
defined by taking the union of a minimal spanning disk for
$p_1\widetilde \Gamma^\lambda_c(t')$ and the annulus
$p_1S^\lambda_c[t',t'']$. As before, the derivative of the area of
this family of disks has two contributions, one coming from the
change in the metric over the minimal spanning disk at time $t'$ and
the other which we computed above to be at most $C_1'$. Thus, the
derivative is bounded above by
$C'_2A(p_1\widetilde\Gamma^\lambda_c(t'))+C_1'$. This implies that
the forward difference quotient of $A(p_1\widetilde
\Gamma^\lambda_c(t'))$ is bounded above by the same quantity. It
follows immediately that the areas of all the minimal spanning
surfaces are bounded by a constant depending only on the areas of
the minimal spanning surfaces at time $t_0$, the sectional curvature
of the ambient Ricci flow and $t_1-t_0$. Hence, there is a constant
$C_4$ such that the forward difference quotient of
$A(p_1\widetilde\Gamma^\lambda_c(t))$ is bounded above by $C_4$.
This proves the claim.
\end{proof}



Next, by the uniform bounds on total length of all the curves
$\widetilde\Gamma_c^\lambda(t)$, it follows from
Equation~(\ref{Lformula}) that there is a constant $C_5$ (we take
$C_5>1$) depending only on the curvature bound of the ambient
manifolds and the family $\Gamma$ such that for any $c\in S^2$ we
have
\begin{equation}\label{C5eqn}\int_{t_0}^{t_1}\int_{\widetilde\Gamma_c^\lambda(t)}
k^2dsdt\le C_5. \end{equation} Thus, for any constant $1<B<\infty$
there is a subset $I_B(c,\lambda)\subset [t_0,t_1]$ of measure at
least $(t_1-t_0)-C_5B^{-1}$ such that
$$\int_{\widetilde\Gamma_c^\lambda(t)}k^2ds\le B$$
for every $t\in I_B(c,\lambda)$. (Later, we shall fix $B$
sufficiently large depending on $\Gamma$ and $\zeta$.)

Now we need a result for curve-shrinking that in some ways is
reminiscent of Shi's theorem for Ricci flows.

\begin{lemma}[Shi-type curvature estimate for curve-shrinking flow]
\label{lem:curve-shrinking-shi-estimate}
\label{CSSHI}
\uses{prop:curve-deformation-existence}
Let $(M,g(t)),\ t_0\le t\le t_1$, be a Ricci flow. Then there exist
constants $\delta>0$ and $\widetilde C_i<\infty$ for
$i=0,1,2,\ldots$, depending only on $t_1-t_0$  and a bound for the
norm of the curvature of the Ricci flow, such that the following
holds. Let $c(x,t)$ be a curve-shrinking flow that is an immersion
for each $t$. Suppose that at a time $t'$ for some $0<r<1$ such that
$t'+\delta r^2<t_1$, the length of $c(\cdot,t')$ is at least $r$ and
the total curvature of $c(\cdot,t')$ on any subarc of length $r$ is
at most $\delta$. Then for every $t\in [t',t'+\delta r^2)$ the
curvature $k$ and the higher derivatives satisfy
$$k^2\le\widetilde C_0(t-t')^{-1}$$
$$|\nabla_SH|^2\le \widetilde C_1(t-t')^{-2},$$
$$|\nabla_S^iH|^2\le \widetilde C_i(t-t')^{-(i+1)}.$$
\end{lemma}

The first statement follows from arguments very similar to those in
Section 4 of \cite{AG}. Once $k^2$ is bounded by $\widetilde
C_0/(t-t')$ the higher derivative statements are standard, see
\cite{Altschuler}. For completeness we have included the proof of
the first inequality in the last section of this chapter.




We now fix $\delta>0$ (and also $\delta<1$) as described in the last
lemma for the Ricci flow  $(M,g(t)),\ t_0\le t\le t_1$. By
Cauchy-Schwarz it follows that for every $t\in I_B(c,\lambda)$,
  and for any arc $J$ in $\widetilde\Gamma_c^\lambda(\cdot,t)$ of
length at most $\delta^2B^{-1}$ we have
$$\int_{J\times\{t\}}k\le \delta.$$
Applying the previous lemma,  for each $a\in I_B(c,\lambda)$ with
$a\le t_1-B^{-1}-\delta^5B^{-2}$ we set
$J(a)=[a+\delta^5B^{-2}/2,a+\delta^5B^{-2}]\subset
[t_0,t_1-B^{-1}]$. Then for all $t\in \cup_{a\in
I_B(c,\lambda)}J(a)$ for which the length of $\widetilde
\Gamma_c^\lambda(\cdot,t)$ is at least $\delta^2B^{-1}$ we have that
$k$ and all the norms of spatial derivatives of $H$ are pointwise
uniformly bounded. Since $I_B(c,\lambda)$ covers all of $[t_0,t_1]$
except a subset of measure at most $C_5B^{-1}$, it follows that the
union $\widehat J_B(c,\lambda)$ of intervals $J(a)$ for $a\in
I_B(c,\lambda)\cap [t_0,t_1-B^{-1}-\delta^5B^{-2}]$ cover all of
$[t_0,t_1]$ except a subset of measure at most
$C_5B^{-1}+B^{-1}+\delta^5B^{-1}<3C_5B^{-1}$. Now it is
straightforward to pass to a finite subset of these intervals
$J(a_i)$ that cover all of $[t_0,t_1]$ except a subset of measure at
most $3C_5B^{-1}$.
 Once we
have a finite number of $J(a_i)$, we order them along the interval
$[t_0,t_1]$ so that  their initial points form an increasing
sequence. (Recall that they all have the same length.) Then if we
have $J_i\cap J_{i+2}\not=\emptyset$, then $J_{i+1}$ is contained in
the union of $J_i$ and $J_{i+2}$ and hence can be removed from the
collection without changing the union. In this way we reduce to a
finite collection of intervals $J_i$, with the same union, where
every point of $[t_0,t_1]$ is contained in at most $2$ of the
intervals in the collection. Once we have arranged this we have a
uniform bound, independent of $\lambda$ and $c\in S^2$, on the
number of these intervals. We let $J_B(c,\lambda)$ be the union of
these intervals. According to the construction and Lemma~\ref{lem:curve-shrinking-shi-estimate}
these sets $J_B(c,\lambda)$ satisfy the following:
\begin{enumerate}
\item[(1)] $J_B(c,\lambda)\subset [t_0,t_1-B^{-1}]$ is a union of a bounded number of intervals
(the bound being independent of $c\in S^2$ and of $\lambda$) of
length $\delta^5B^{-2}/2$.
\item[(2)] The measure of $J_B(c,\lambda)$ is at least
$t_1-t_0-3C_5B^{-1}$.
\item[(3)] For every $t\in J_B(c,\lambda)$ either the length of $\widetilde\Gamma_c^\lambda(t)$ is
less than $\delta^2 B^{-1}$ or there are uniform bounds, depending
only on the curvature bounds of the ambient Ricci flow and the
initial family $\Gamma$, on the curvature and its higher spatial
derivatives of $\widetilde\Gamma_c^\lambda(t)$.
\end{enumerate}


Now we fix $c\in S^2$ and $1<B<\infty$ and we fix a sequence of
$\lambda_n$ tending to zero. Since the number of intervals in
$J_B(c,\lambda)$ is bounded independent of $\lambda$, by passing to
a subsequence of $\lambda_n$ we can suppose that the number of
intervals in $J_B(c,\lambda_n)$ is independent of $n$, say this
number is $N$, and that their initial points (and hence the entire
intervals since all their lengths are the same) converge as $n$ goes
to infinity. Let $\hat J_1,\ldots, \hat J_N$ be the limit intervals,
and for each $i, 1\le i\le N$, let $J_i\subset \hat J_i$ be a
slightly smaller interval contained in the interior of $\hat J_i$.
We choose the $J_i$ so that they all have the same length. Let
$J_B(c)\subset [t_0,t_1-B^{-1}]$ be the union of the $J_i$. Then an
appropriate choice of the length of the $J_i$ allows us to arrange
the following:
\begin{enumerate}
\item[(1)] $J_B(c)\subset J_B(c,\lambda_n)$ for all $n$ sufficiently
large.
\item[(2)] $J_B(c)$ covers all of $[t_0,t_1]$ except a subset of length
$4C_5B^{-1}$.
\end{enumerate}

Now fix one of the intervals $J_i$ making up $J_B(c)$. After passing
to a subsequence (of the $\lambda_n$), one of the following holds:
\begin{enumerate}
\item[(3)] there are uniform bounds for the curvature and all its
derivatives for the curves $\widetilde\Gamma_c^{\lambda_n}(t)$, for
all $ t\in J_i$ and all $n$, or
\item [(4)] for each $n$ there is $t_n\in J_i$ such that the length of
$\widetilde\Gamma_c^{\lambda_n}(t_n)$ is less than $\delta^2B^{-1}$.
\end{enumerate}

By passing to a further subsequence, we arrange that the same one of
the Alternatives (3) and (4) holds for every  one of the intervals
$J_i$ making up $J_B(c)$.

The next claim is the statement that a slightly stronger version of
Proposition~\ref{prop:curve-deformation-existence} holds for $p_1\widetilde\Gamma_c^\lambda(t)$.


\begin{lemma}[Sharper single-point estimate for the curve deformation]
\label{lem:curve-deformation-single-point-estimate}
\label{xistronger}
Given $\zeta>0$, there is   $1<B<\infty$, with $B>(t_1-t_0)^{-1}$,
depending only on $\Gamma$ and the curvature bounds on the ambient
Ricci flow $(M,g(t)),\ t_0\le t\le t_1$, such that the following
holds. Let $t_2=t_1-B^{-1}$. Fix $ c\in S^2$. Let $v_c$ be the
solution to Equation~(\ref{wdiffeqn}) with initial condition
$v_{c}(t_0)=A(p_1(\widetilde\Gamma(c)))$, so that in our previous
notation $v_c=w_{A(p_1(\widetilde\Gamma(c)))}$. Then for all
$\lambda>0$ sufficiently small, either
$A(p_1\widetilde\Gamma_c^\lambda(t_1))<v_c(t_1)+\zeta/2$ or the
length of $\widetilde\Gamma_c^\lambda(t)$ is less than $\zeta/2$ for
all $t\in [t_2,t_1]$.
\end{lemma}



\begin{proof}
\uses{claim:area-forward-difference-linear-bound,claim:w-a-t-explicit-formula,cor:curve-shrinking-length-curvature-exponential-bound,claim:piecewise-domination-estimate,claim:curve-shrinking-short-time-existence,lem:immersed-curve-shrinking-area-bound,claim:curve-shrinking-flow-limit,rem:short-loop-length-bound,claim:approximate-piecewise-domination-estimate}
In order to establish this lemma  we need a couple of claims about
functions on $[t_0,t_1]$ that are approximately dominated by
solutions to Equation~(\ref{wdiffeqn}). In the first claim the
function in question is dominated on a finite collection of
subintervals by solutions to these equations and the subintervals
fill up most of the interval. In the second, we also allow the
function to only be approximately dominated by the solutions to
Equation~(\ref{wdiffeqn}) on these sub-intervals. In both claims the
result is that on the entire interval the function is almost
dominated by the solution to the equation with the same initial
value.


\begin{claim}[Piecewise domination by the comparison ODE]
\label{claim:piecewise-domination-estimate}
\label{C1C4delta}
\uses{claim:area-forward-difference-linear-bound}
Fix $C_4$ as in Claim~\ref{claim:area-forward-difference-linear-bound} and fix a constant
$\widetilde A>0$. Given $\zeta>0$ there is $\delta'>0$ depending on
$C_4$, $t_1-t_0$, and $\widetilde A$ as well as the curvature bound
of the ambient Ricci flow such that the following holds. Suppose
that $f\colon [t_0,t_1]\to \Ar$ is a function and suppose that
$J\subset [t_0,t_1]$ is a finite union of intervals. Suppose that on
each interval $[a,b]$ of $J$ the function $f$ satisfies
$$f(b)\le w_{f(a),a}(b).$$ Suppose further that for any $t'<t''$ we
have
$$f(t'')\le f(t')+C_4(t''-t').$$
Then, provided that the total length of $[t_0,t_1]\setminus J$ is at
most $\delta'$ and $0\le f(t_0)\le \widetilde A$, we have
$$f(t_1)\le w_{f(t_0),t_0}(t_1)+\zeta/4.$$
\end{claim}

\begin{proof}
\uses{claim:w-a-t-explicit-formula}
We write $J$ as a union of disjoint intervals $J_1,\ldots,J_k$ so
that  $J_i<J_{i+1}$ for every $i$. Let $a_i$, resp. $b_i$, be the
initial, resp. final, point of $J_i$. For each $i$ let $\delta_i$ be
the length of the interval between $J_i$ and $J_{i+1}$. (Also, we
set $\delta_0=a_1-t_0$, and $\delta_{k}=t_1-b_k$.) Let $C_6\ge 0$ be
such that $R_{\min}(t)\ge -2C_6$ for all $t\in [t_0,t_1]$. Let
$V(a)$ be the maximum value of $|w_{a,t_0}|$ on the interval
$[t_0,t_1]$ and let $V=\max_{a\in [0,\widetilde A]}V(a)$. let
$C_7=C_4+2\pi+C_6V$.
 We shall prove by
induction that
$$f(a_i)-w_{f(t_0),t_0}(a_i)\le \sum_{j=0}^{i-1}\left(C_7\delta_j
\prod_{\ell=j+1}^{i-1}e^{C_6|J_\ell|}\right)$$ and
$$f(b_i)-w_{f(t_0),t_0}(b_i)\le
\sum_{j=0}^{i-1}\left(C_7\delta_j\prod_{\ell=j+1}^ie^{
C_6|J_\ell|}\right).$$

We begin the induction by establishing the result at $a_1$. By
hypothesis we know that
$$f(a_1)\le f(t_0)+C_4\delta_0.$$
 On the other
hand, from the defining differential equation for $w_{f(t_0),t_0}$
and the definitions of $C_6$ and $V$ we have
$$w_{f(t_0),t_0}(a_1)\ge f(t_0) -(C_6V+2\pi)\delta_0.$$
Thus,
$$f(a_1)-w_{f(t_0),t_0}(a_1)\le (C_4+2\pi+C_6V)\delta_0=C_7\delta_0,$$
which is exactly the formula given in the case of $a_1$.

Now suppose that we know the result for $a_i$ and let us establish
it for $b_i$. Let $\alpha_i=f(a_i)-w_{f(t_0),t_0}(a_i)$, and let
$\beta_i=f(b_i)-w_{f(t_0),t_0}(b_i)$. Then by Claim~\ref{claim:w-a-t-explicit-formula} we
have
$$\beta_i\le e^{C_6|J_i|}\alpha_i.$$
Given the inductive inequality for $\alpha_i$, we immediately get
the one for $\beta_i$.

Now suppose that we have the inductive inequality for $\beta_i$.
Then
$$f(a_{i+1})\le f(b_i)+C_4\delta_i.$$
 On the other hand, by
 the definition of $C_6$ and $V$ we have
$$w_{f(t_0),t_0}(a_{i+1})-w_{f(t_0),t_0}(b_i)\ge
-(C_6V+2\pi)\delta_{i}.$$ This yields
$$f(a_{i+1})-w_{f(t_0),t_0}(a_{i+1})\le \beta_i+C_7\delta_i.$$
Hence, the inductive result for $\beta_i$ implies the result for
$\alpha_{i+1}$. This completes the induction.

Applying this to $a_{k+1}=t_1$ gives
$$f(t_1)-w_{f(t_0),t_0}(t_1)\le
\sum_{j=0}^k\left(C_7\delta_j\prod_{\ell=j+1}^{k}e^{
C_6|J_\ell|}\right)\le C_7\sum_{j=0}^k\delta_je^{C_6(t_1-t_0)}.$$

Of course $\sum_{j=1}^k\delta_j=t_1-t_0-\ell(J)\le\delta'$, and
$C_7$ only depends on $ C_6,C_4$ and $V$, while $V$  only depends on
$\widetilde A$ and $C_6$ only depends on the sectional curvature
bound on the ambient Ricci flow. Thus, given $C_4, \widetilde A$ and
$t_1-t_0$ and the bound on the sectional curvature of the ambient
Ricci flow, making $\delta'$ sufficiently small makes
$f(t_1)-w_{f(t_0),t_0}(t_1)$ arbitrarily small. This completes the
proof of the claim.
\end{proof}

Here is the second of our claims:

\begin{claim}[Approximate piecewise domination by the comparison ODE]
\label{claim:approximate-piecewise-domination-estimate}
\label{deltaprime}
\uses{claim:piecewise-domination-estimate}
Fix $\zeta>0$, $A$ and $C_6, C_4$ as in the last claim, and let
$\delta'>0$ be as in the last claim. Suppose that we have $J\subset
[t_0,t_1]$ which is a finite disjoint union of intervals with
$t_1-t_0-|J|\le \delta'$.
 Then there
is $\delta''>0$ ($\delta''$ is allowed to depend on $J$) such that
the following holds. Suppose that we have a function $f\colon
[t_0,t_1]\to \Ar$ such that:
\begin{enumerate}
\item[(1)] For all $t'<t''$ in $[t_0,t_1]$ we have $f(t'')-f(t')\le
C_4(t''-t')$.
\item[(2)] For any interval $[a,b]\subset J$ we have
$f(b)\le w_{f(a),a}(b)+\delta''$.
\end{enumerate}
Then $f(t_1)\le w_{f(t_0),t_0}(t_1)+\zeta/2$.
\end{claim}

\begin{proof}
\uses{claim:piecewise-domination-estimate}
We define $C_7$ as in the previous proof. We use the notation
$J=J_1\coprod\cdots\coprod J_k$ with $J_1<J_2<\cdots<J_k$ and let
$\delta_i$ be the length of the interval separating $J_{i-1}$ and
$J_i$. The arguments in the proof of the previous claim work in this
context to show that
$$f(a_i)-w_{f(t_0),t_0}(a_i)\le
\sum_{j=0}^{i-1}\left(C_7\delta_j\prod_{\ell=j+1}^{i-1}(e^{C_6|J_\ell|}+\delta'')\right).$$
Applying this to $a_{k+1}$ and taking the limit as $\delta''$ tends
to zero, the right-hand side tends to a limit smaller than
$\zeta/4$. Hence, for $\delta''$ sufficiently small the right-hand
side is less than $\zeta/2$.
\end{proof}

Now let us return to the proof of Lemma~\ref{lem:curve-deformation-single-point-estimate}. Recall
that $c\in S^2$ is fixed. We shall apply the above claims to the
curve-shrinking flow $\widetilde\Gamma_c^\lambda(t)$ and thus prove
Lemma~\ref{lem:curve-deformation-single-point-estimate}. Now it is time to fix $B$. First, we fix
$\widetilde A=W(\Gamma)+\zeta$,  we let $C_2$ be as in
Corollary~\ref{cor:curve-shrinking-length-curvature-exponential-bound}, $C_4$ be as in Claim~\ref{claim:area-forward-difference-linear-bound},
$C_5$ be as in Equation~(\ref{C5eqn}), and $C_6$ be as in the proof
of Claim~\ref{claim:piecewise-domination-estimate}. Then we have $\delta'$ depending on
$C_6,C_4,\widetilde A$ as in Claim~\ref{claim:piecewise-domination-estimate}. We fix $B$ so
that:
\begin{enumerate}
\item[(1)]  $B\ge 3C_5(\delta')^{-1}$,
\item[(2)] $B\ge 3e^{C_2(t_1-t_0)}\zeta^{-1}$, and
\item[(3)] $B>C_2/(\log 4-\log 3)$.
\end{enumerate}
The first step in the proof of Lemma~\ref{lem:curve-deformation-single-point-estimate} is the
following:

\begin{claim}[Limit of the curve-shrinking flows]
\label{claim:curve-shrinking-flow-limit}
\label{limitcs}
After passing to a subsequence of $\{\lambda_n\}$, either:
\begin{enumerate}
\item[(1)] for each $n$ sufficiently large there is $t_n\in J_B(c)$ with
the length of $\widetilde\Gamma_c^{\lambda_n}(t_n)<\delta^2B^{-1}$,
or
\item[(2)] for each component $J_i=[t_i^-,t_i^+]$ of $ J_B(c)$, after
composing $\widetilde\Gamma_c^{\lambda_n}(x,t)$ by a
reparameterization of the domain circle (fixed in $t$ but a
different reparameterization for each $n$) so that the $\widetilde
\Gamma_c^{\lambda_n}(t_i^-)$ have constant speed, there is a smooth
limiting curve-shrinking flow denoted $\widetilde\Gamma_c(t)$, for
$t\in J_i$ for the sequence $p_1\widetilde\Gamma_c^{\lambda_n}(t),\
t_i^-\le t\le t_i^+$. The limiting flow consists of immersions.
\end{enumerate}
\end{claim}

\begin{proof}
\uses{claim:curve-shrinking-short-time-existence}
Suppose that the first case does not hold for any subsequence. Fix a
component $J_i$ of $J_B(c)$. Then, by passing to a subsequence, by
the fact that $J_B(c)\subset J_B(c,\lambda_n)$ for all $n$, the
curvatures and all the derivatives of the curvatures of
$\widetilde\Gamma_c^{\lambda_n}(t)$ are uniformly bounded
independent of $n$ for all $t\in J_B(c)$. We reparameterize the
domain circle so that the $\widetilde\Gamma_c^{\lambda_n}(t_i^-)$
have constant speed. By passing to a subsequence we can suppose that
the lengths of the $\widetilde\Gamma_c^{\lambda_n}(t_i^-)$ converge.
The limit is automatically positive since we are assuming that the
first case does not hold for any subsequence. Denote by
$S_n=S_c^{\lambda_n}(t_i^-)$ the unit tangent vector to $\widetilde
\Gamma_c^{\lambda_n}(t_i^-)$ and by $u_n$ the inner product $\langle
S_n,U\rangle$. Now we have a family of loops with tangent vectors
and all higher derivatives bounded. Since $u_n$ is everywhere
positive, since $\int u_nds=\lambda_n$, since the length of the loop
$\widetilde \Gamma_c^{\lambda_n}(t_i^-)$ is bounded away from $0$
independent of $n$, and since
$|(u_n)'|=|\langle\nabla_{S_n}S_n,U\rangle|$ is bounded above
independent of $n$, we see that $u_n$ tends uniformly to zero as $n$
tends to infinity. This means that the $|p_1(S_n)|$ converge
uniformly to one as $n$ goes to infinity.
 Since
the ambient manifold is compact, passing to a further subsequence we
have a smooth limit of the
$p_1\widetilde\Gamma_c^{\lambda_n}(t_i^-)$. The result is an
immersed curve in $(M,g(t_i^-))$ parameterized at unit speed. Since
all the spatial and time derivatives of the $p_1\widetilde
\Gamma_c^{\lambda_n}(t)$ are uniformly bounded, by passing to a
further subsequence, there is a  smooth map $f\colon S^1\times
[t_i^-,t_i^+]\to M$ which is a smooth limit of the sequence
$\widetilde\Gamma_c^{\lambda_n}(t),\ t_i^-\le t\le t_i^+$. If for
some $t\in [t_i^-,t_i^+]$ the curve $f|_{S^1\times\{t\}}$ is
immersed, then this limiting map along this curve agrees to first
order with the curve-shrinking flow. Thus, for some $t>t_i^-$ the
restriction of $f$ to the interval $[t_i^-,t]$ is a curve-shrinking
flow. We claim that $f$ is a curve-shrinking flow on the entire
interval $[t_i^-,t_i^+]$. Suppose not. Then there is a first $t'\le
t_i^+$ for which $f|_{S^1\times \{t'\}}$ is not an immersion.
According to Lemma~\ref{claim:curve-shrinking-short-time-existence} the maximum of the norms of
the curvature of the curves $f(t)$ must tend to infinity as $t$
approaches $t'$ from below. But the curvatures of $f(t)$ are the
limits of the curvatures of the family $p_1\widetilde
\Gamma_c^{\lambda_n}(t)$ and hence are uniformly bounded on the
entire interval $[t_i^-,t_i^+]$.
 This contradiction shows that the
entire limiting surface
$$f\colon S^1\times [t_i^-,t_i^+]\to (M,g(t))$$
is a curve-shrinking flow of immersions.
\end{proof}

\begin{remark}[Length bound in the short-loop alternative]
\label{rem:short-loop-length-bound}
\label{xicomp}
\uses{lem:immersed-curve-shrinking-area-bound,claim:curve-shrinking-flow-limit,claim:area-forward-difference-linear-bound,claim:piecewise-domination-estimate,claim:approximate-piecewise-domination-estimate,cor:curve-shrinking-length-curvature-exponential-bound,lem:curve-deformation-single-point-estimate,prop:curve-deformation-existence}
Notice that if the first case holds then by the choice of $B$ we
have a point $t_n\in J_B(c)$ for which the length of
$\widetilde\Gamma_c^{\lambda_n}(t_n)$ is less than
$e^{-C_2(t_1-t_0)}\zeta/3$.
\end{remark}


For each $n$,  the family of curves
$p_1\widetilde\Gamma_c^{\lambda_n}(t)$  in $M$ all have
$p_1\widetilde\Gamma_c^{\lambda_n}(t_0)=\widetilde\Gamma(c)$ as
their initial member. Thus, these curves are all homotopically
trivial. Hence, for each $t\in J_B(c)$ the limiting curve
$\widetilde\Gamma(c)(t)$ of the
$p_1\widetilde\Gamma_c^{\lambda_n}(t)$ is then also homotopically
trivial. It now follows from Lemma~\ref{lem:immersed-curve-shrinking-area-bound},
Claim~\ref{claim:curve-shrinking-flow-limit} and Remark~\ref{rem:short-loop-length-bound} that one of the
following two conditions holds:
\begin{enumerate}
\item[(1)] for some $t\in J_B(c)$ the length of $\widetilde\Gamma(c)(t)$ is less than or equal to
 $e^{-C_2(t_1-t_0)}\zeta/3$ or
\item[(2)] the function $A(t)$ that assigns to each $t\in
J_B(c)$ the area of the minimal spanning disk for
$p_1\widetilde\Gamma(c)(t)$ satisfies
$$\frac{dA(t)}{dt}\le -2\pi-\frac{1}{2}R_{\min}(t)A(t)$$
in the sense of forward difference quotients.
\end{enumerate}
By continuity,  for any $\delta''>0$ then for all $n$
 sufficiently large one of the following two conditions holds:
\begin{enumerate}
\item[(1)] there is $t_n\in J_B(c)$ such that the length of $\widetilde\Gamma_c^{\lambda_n}(t_n)$
is less than $e^{-C_2(t_1-t_0)}\zeta/2$, or
\item[(2)]  for every $t\in J_B(c)$, the areas of the minimal
spanning disks for $p_1(\widetilde\Gamma_c^{\lambda_n}(t))$ satisfy
$$\frac{dA(p_1\widetilde\Gamma_c^{\lambda_n}(t))}{dt}\le -2\pi-\frac{1}{2}R_\mathrm{min}(t)A(p_1\widetilde\Gamma_c^{\lambda_n}(t))+\delta''$$ in the
sense of forward difference quotients.
\end{enumerate}

Suppose  that for every $n$ sufficiently large, for every  $t\in
J_B(c)$ the length of $\widetilde\Gamma_c^{\lambda_n}(t)$ is at
least $e^{-C_2(t_1-t_0)}\zeta/2$. We have already seen in
Claim~\ref{claim:area-forward-difference-linear-bound} that for every $t'<t''$ in $[t_0,t_1]$ the
areas satisfy
$$A(p_1(\widetilde\Gamma_c^{\lambda_n}(t'')))-A(p_1(\widetilde\Gamma_c^{\lambda_n}(t')))\le
C_4(t''-t').$$ Since the total length of the complement $J_B(c)$ in
$[t_0,t_1]$ is at most $3C_5B^{-1}$, it follows from our choice of
$B$ that this total length is at most the constant $\delta'$ of
Claim~\ref{claim:piecewise-domination-estimate}. Invoking Claim~\ref{claim:approximate-piecewise-domination-estimate} and the fact
that $A(\widetilde\Gamma(c))\le W(\Gamma(c))+\zeta\le
W(\Gamma)+\zeta=\widetilde A$, we see that for all $n$ sufficiently
large we have
$$A(p_1(\widetilde\Gamma^c_{\lambda_n}(t_1)))-v_c(t_1)<\zeta/2.$$

The other possibility to consider is that for each $n$ there is
$t_n\in J_B(c)$ such that the length of
$\widetilde\Gamma_c^{\lambda_n}(t_n)<e^{-C_2(t_1-t_0)}\zeta/2$.
Since $J_B(c)\subset [t_0,t_1-B^{-1}]$, in this case we invoke the
first inequality in Corollary~\ref{cor:curve-shrinking-length-curvature-exponential-bound} to see that the
length of $\widetilde\Gamma_c^{\lambda_n}(t)<\zeta/2$ for every
$t\in [t_1-B^{-1},t_1]$. This completes the proof of
Lemma~\ref{lem:curve-deformation-single-point-estimate}.
\end{proof}

\subsection{The completion of the proof of
Proposition~{\ref{prop:curve-deformation-existence}}}

Now we wish to pass from  Lemma~\ref{lem:curve-deformation-single-point-estimate} which deals with an
individual $c\in S^2$ to a proof of Proposition~\ref{prop:curve-deformation-existence} which deals
with the entire family $\widetilde\Gamma$. Let us introduce the
following notation. Suppose that $\omega\subset S^2$ is an arc. Then
$\widetilde\Gamma(\omega)=\cup_{c\in \omega}\widetilde \Gamma(c)$ is
an annulus in $M$ and for each $t\in [t_0,t_1]$ we have
 the annulus $\widetilde\Gamma_\omega^\lambda(t)$ in $M\times
S^1_\lambda$.



A finite set ${\mathcal S}\subset S^2$ with the property that for
$c\in S^2$ there is $\hat c\in {\mathcal S}$ and an arc $\omega$ in
$S^2$ joining $c$ to $\hat c$ so that the area of the annulus
$\widetilde\Gamma(\omega)$ is less than $\nu$ is called {\em a
$\nu$-net} for $\widetilde\Gamma$. Similarly, if for every $c\in
S^2$ there is $\hat c\in {\mathcal S}$ and an arc $\omega$
connecting them for which the area of the annulus
$\widetilde\Gamma_\omega^\lambda(t_0)$ is less than $\nu$, we say
that ${\mathcal S}$ is {\em a $\nu$-net} for
$\widetilde\Gamma^\lambda$. Clearly, for any $\nu$ there is a subset
${\mathcal S}\subset S^2$ that is a $\nu$-net for $\widetilde\Gamma$
and for $\widetilde\Gamma^\lambda$ for all $\lambda$ sufficiently
small.



\begin{lemma}[Area comparison across a mu-net arc]
\label{lem:mu-net-area-comparison}
\label{muarea}
There is a $\mu>0$ such that the following holds. Let $c,\hat c\in
S^2$. Suppose that there is an arc $\omega$ in $S^2$ connecting $c$
to $\hat c$ with the area of the annulus
$\widetilde\Gamma_\omega^\lambda(t_0)$ in $M\times S^1_\lambda$ less
than $\mu$. Let $v_{\hat c}$, resp., $v_c$, be the solution to
Equation~(\ref{wdiffeqn}) with initial condition $v_{\hat
c}(t_0)=A(\widetilde\Gamma(\hat c))$, resp.,
$v_c(t_0)=A(\widetilde\Gamma(c))$. If
 $$A(p_1\widetilde\Gamma^\lambda_{\hat c}(t_1))\le v_{\hat c}+\zeta/2,$$
then
$$A(p_1(\widetilde\Gamma^\lambda_c(t_1))\le v_c+\zeta.$$
\end{lemma}


\begin{proof}
\uses{lem:minimal-annulus-area-growth}
 First of all we require that $\mu<e^{-(2n-1)C'(t_1-t_0)}\zeta/4$
where $C'$ is an upper bound for the norm of the Riemann curvature
tensor at any point of the ambient Ricci flow.  By
Lemma~\ref{lem:minimal-annulus-area-growth} the fact that the area of the minimal annulus
between the ramps $\widetilde\Gamma_c^\lambda(t_0)$ and
$\widetilde\Gamma_{\hat c}^\lambda(t_0)$ is less than $\mu$ implies
that the area of the minimal annulus between the ramps
$\widetilde\Gamma_c^\lambda(t_1)$ and $\widetilde\Gamma_{\hat
c}^\lambda(t_1)$ is less than $\mu e^{(2n-1)C'(t_1-t_0)}=\zeta/4$.
The same estimate also holds for the image under the projection
$p_1$ of this minimal annulus. Thus, with this condition on $\mu$,
and for $\lambda$ sufficiently small, we have
$$\left|A(p_1\widetilde\Gamma_c^\lambda(t_1))-A(p_1\widetilde\Gamma_{\hat c}^\lambda(t_1))\right|<\zeta/4.$$
The other condition we impose upon $\mu$ is that if $a,\hat a$ are
positive numbers at most $W(\Gamma)+\zeta$ and if $a<\hat a+\mu$
then
$$w_{a,t_0}(t_1)<w_{\hat a,t_0}(t_1)+\zeta/4.$$
Applying this with $a=A(\widetilde\Gamma(c))$ and $\hat
a=A(\widetilde\Gamma(\hat c))$ (both of which are at most
$W(\widetilde \Gamma)<W(\Gamma)+\zeta$), we see that these two
conditions on $\mu$ together imply the result.
\end{proof}

 We must also examine what
happens if the second alternative holds for $\widetilde\Gamma_{\hat
c}^\lambda$. We need the following lemma to treat this case.

\begin{lemma}[Length comparison for ramps joined by a small annulus]
\label{lem:ramp-length-comparison-small-annulus}
\label{BARMU}
\uses{prop:curve-deformation-existence}
There is  $\delta>0$ such that for any $r>0$  there is $\bar\mu>0$,
depending on $r$ and on the curvature bound for the ambient Ricci
flow such that the following holds. Suppose that $\gamma$ and $\hat
\gamma$ are ramps in $(M,g(t))\times S^1_\lambda$. Suppose that the
length of $\gamma$ is at least $r$ and suppose that on any
sub-interval $I$ of $\gamma$ of length $r$ we have
$$\int_Ikds<\delta.$$
 Suppose also that there is an annulus connecting $\gamma$ and
$\hat\gamma$ of area less than $\bar\mu$. Then the length of
$\hat\gamma$ is at least $3/4$ the length of $\gamma$.
\end{lemma}

We give a proof of this lemma in the next section. Here we finish
the proof of Proposition~\ref{prop:curve-deformation-existence} assuming it.

\begin{claim}[Length comparison across a mu-net arc]
\label{claim:mu-net-length-comparison}
\label{mulength}
There is $\mu>0$ such that the following holds. Suppose that $c,\hat
c\in S^2$ are such that there is an arc $\omega$ in $S^2$ connecting
$c$ and $\hat c$ such that the area of the annulus
$\widetilde\Gamma^\lambda(\omega)$ is at most $\mu$.  Set
$t_2=t_1-B^{-1}$. If the length of $\widetilde\Gamma_{\hat
c}^\lambda(t)$ is less than $\zeta/2$ for all $t\in [t_2,t_1]$, then
the length of $p_1\widetilde\Gamma_c^\lambda(t_1)$ is less than
$\zeta$.
\end{claim}

\begin{proof}
\uses{cor:curve-shrinking-length-curvature-exponential-bound,lem:ramp-length-comparison-small-annulus,lem:minimal-annulus-area-growth}
The proof is by contradiction: Suppose that the length of
$p_1\widetilde\Gamma_c^\lambda(t_1)$ is at least $\zeta$ and the
length of $\widetilde\Gamma_{\hat c}^\lambda(t)$ is less than
$\zeta/2$ for all $t\in [t_2,t_1]$. Of course, it follows that the
length of $\widetilde\Gamma_c^\lambda(t_1)$ is also at least
$\zeta$. The third condition on $B$ is equivalent to
$$e^{C_2B^{-1}}<4/3.$$
It then follows from
 Corollary~\ref{cor:curve-shrinking-length-curvature-exponential-bound} that for every $t\in
 [t_2,t_1]$ the length of $\widetilde\Gamma^c_\lambda(t)$ is at least $3\zeta/4$.
On the other hand, by hypothesis for every such $t$, the length of
$\widetilde\Gamma_{\hat c}^\lambda(t)$ is less than $\zeta/2$. It
follows from Equation~(\ref{Lformula}) that
$$\int_{t_2}^{t_1}\left(\int k^2ds\right)dt\le
C_2\left(\int_{t_2}^{t_1}L(\widetilde\Gamma_c^\lambda(t))dt\right)-
L(\widetilde\Gamma_c^\lambda(t_1))+L(\widetilde\Gamma_c^\lambda(t_2)).$$
(Here $L$ is the length of the curve.) From this and
Corollary~\ref{cor:curve-shrinking-length-curvature-exponential-bound} we see that there is a constant $C_8$
depending on the original family $\Gamma$ and on the curvature of
the ambient Ricci flow such that
$$\int_{t_2}^{t_1}\left(\int_{\widetilde\Gamma_c^\lambda(t)} k^2ds\right)dt\le C_8.$$
Since $t_1-t_2=B^{-1}$, this implies that there is $t'\in [t_2,t_1]$
with
$$\int_{\widetilde\Gamma_c^\lambda(t')}k^2ds\le C_8B.$$
By Cauchy-Schwarz, for any subinterval $I$ of length $\le r$ in
$\widetilde\Gamma_c^\lambda(t')$ we have
$$\int_Ikds\le \sqrt{C_8Br}.$$
We choose $0<r\le \zeta$ sufficiently small so that $\sqrt{C_8Br}$
is less than or equal to the constant $\delta$ given in
Lemma~\ref{lem:ramp-length-comparison-small-annulus}. Then we set $\bar \mu$ equal to the constant
given by that lemma for this value of $r$.

Now suppose that $\mu$ is sufficiently small so that the solution to
the equation
$$\frac{d\mu(t)}{dt}=(2n-1)|\Rm_{g(t)}|\mu(t)$$
with initial condition $\mu(t_0)\le \mu$ is less than $\bar\mu$ on
the entire interval $[t_0,t_1]$. With this condition on $\mu$,
Lemma~\ref{lem:minimal-annulus-area-growth} implies that for every $t\in [t_0,t_1]$ the
ramps $\widetilde\Gamma_c^\lambda(t)$ and $\widetilde\Gamma_{\hat
c}^\lambda(t)$ are connected by an annulus of area at most
$\bar\mu$. In particular, this is true for
$\widetilde\Gamma_c^\lambda(t')$ and $\widetilde\Gamma_{\hat
c}^\lambda(t')$. Now we have all the hypotheses of Lemma~\ref{lem:ramp-length-comparison-small-annulus}
at time $t'$. Applying this lemma we conclude that
$$L(\widetilde\Gamma_{\hat c}^\lambda(t'))\ge \frac{3}{4}L(\widetilde\Gamma_c^\lambda(t')).$$
But this is a contradiction since by assumption
$L(\widetilde\Gamma_{\hat c}^\lambda(t'))<\zeta/2$ and the
supposition that $L(p_1\widetilde\Gamma_c^\lambda(t_1))\ge \zeta$
led to the conclusion that $L(\widetilde\Gamma_c^\lambda(t'))\ge
3\zeta/4$. This contradiction shows that our supposition that
$L(p_1\widetilde\Gamma_c^\lambda(t_1))\ge \zeta$ is false.
\end{proof}

Now we complete the proof of Proposition~\ref{prop:curve-deformation-existence}.

\begin{proof}
\uses{lem:mu-net-area-comparison,lem:curve-deformation-single-point-estimate}
 (of Proposition~\ref{prop:curve-deformation-existence}.)
Fix $\mu>0$ sufficiently small so that Lemma~\ref{lem:mu-net-area-comparison} and
Claim~\ref{claim:mu-net-length-comparison} hold. Then we choose a $\mu/2$-net $X$ for
$\widetilde\Gamma$. We take $\lambda$ sufficiently small so that
Lemma~\ref{lem:curve-deformation-single-point-estimate} holds for every $\hat c\in {\mathcal S}$. We
also choose $\lambda$ sufficiently small so that $X$ is a $\mu$-net
for $\widetilde\Gamma^\lambda$. Let $c\in S^2$. Then there is $\hat
c\in {\mathcal S}$ and an arc $\omega$ connecting $c$ and $\hat c$
such that the area of $\widetilde\Gamma^\lambda(\omega)<\mu$. Let
$v_{\hat c}$, resp., $v_c$ be the solution to
Equation~\ref{wdiffeqn} with initial condition $v_{\hat
c}(t_0)=A(\widetilde\Gamma({\hat c}))$, resp.,
$v_c(t_0)=A(\widetilde\Gamma(c))$. According to
Lemma~\ref{lem:curve-deformation-single-point-estimate} either $A(p_1\widetilde\Gamma_{\hat
c}^\lambda(t_1))<v_{\hat c}(t_1)+\zeta/2$ or the length of
$\widetilde\Gamma_{\hat c}^\lambda(t)$ is less than $\zeta/2$ for
every $t\in [t_2,t_1]$ where $t_2=t_1-B^{-1}$. In the second case,
Claim~\ref{claim:mu-net-length-comparison} implies that the length of
$p_1\widetilde\Gamma_c^\lambda(t_1)$ is less than $\zeta$. In the
first case, Lemma~\ref{lem:mu-net-area-comparison} tells us that
$A(p_1\widetilde\Gamma_c^\lambda(t_1))<v_c(t_1)+\zeta$. This
completes the proof of Proposition~\ref{prop:curve-deformation-existence}.
\end{proof}


\section{Proof of Lemma~{\ref{lem:ramp-length-comparison-small-annulus}}: annuli of small area}\label{18.5}

Except for the brief comments that follow, our proof involves
geometric analysis that takes place on an abstract annulus with
bounds on its area, upper bounds on its Gaussian curvature, and on
integrals of the geodesic curvature on the boundary.
Proposition~\ref{prop:annulus-boundary-length-estimate} below gives the precise result along
these lines. Before stating that proposition, we show that its
hypotheses hold in the situation that arises in Lemma~\ref{lem:ramp-length-comparison-small-annulus}.
Let us recall  the situation of Lemma~\ref{lem:ramp-length-comparison-small-annulus}. We have ramps
$\gamma$ and $\hat \gamma$ in which are real analytic embedded
curves in the real analytic Riemannian manifold $(M\times
S^1_\lambda,g\times ds^2)$. By a slight perturbation we can assume
they are disjointly embedded. These curves that are connected by an
annulus $A_0\to M\times S^1_\lambda$ of small area, an area bounded
above by, say,  $\mu$. We take an energy minimizing map of an
annulus $\psi\colon A\to M\times S^1_\lambda$ spanning
$\gamma\coprod \hat \gamma$. According to \cite{Hildebrandt}, $\psi$
is a real analytic map and the only possible singularities
(non-immersed points) of the image come from the {\em branch points}
of $\psi$, i.e., points where $d\psi$ vanishes. There are finitely
many branch points. If there are branch points on the boundary, then
the restriction of $\psi$ to $\partial A$ will be a homeomorphism
rather than a diffeomorphism onto $\gamma\coprod \hat \gamma$.
Outside the branch points, $\psi$ is a conformal map onto its image.
The image is an area minimizing annulus spanning $\gamma\coprod \hat
\gamma$. Thus, the area of the image is at most $\mu$. According to
\cite{White} the only branch points on the boundary are false branch
points, meaning that a local smooth reparameterization of the map on
the interior of $A$ near the boundary branch point removes the
branch point. These reparameterizations produce a new smooth
structure on $A$, identified with the original smooth structure on
the complement of the boundary branch points. Using this new smooth
structure on $A$ the map $\psi$ is an immersion except at finitely
many interior branch points. From now on the domain surface $A$ is
endowed with this new smooth structure. Notice that, after this
change, the domain is no longer real analytic; it is  only smooth.
Also, the original annular coordinate is not smooth at the finitely
many boundary branch points.

The pullback of the metric $g\times ds^2$ is a smooth symmetric
two-tensor on $A$. Off the finite set of interior branch points it
is positive definite and hence a Riemannian metric, and in
particular, it is a Riemannian metric near the boundary. It vanishes
at each interior branch point.
  Since the geodesic
curvature $k_\mathrm{geod}$ of the boundary of the annulus is given by $k\cdot n$
where $n$ is the unit normal vector along the boundary pointing into $A$, we
see that the restriction of the geodesic curvature to $\gamma$, $k_\mathrm{geod}\colon \gamma\to \Ar$ has the property that for any sub-arc $I$ of
$\gamma$ of length $r$ we have
$$\int_I|k_\mathrm{geod}|ds<\delta.$$ Lastly, because the map of $A$ into $M\times
S^1_\lambda$ is minimal, off  the set of interior branch points, the
Gaussian curvature of the pulled back metric is bounded above by the
upper bound for the sectional curvature of $M\times S^1_\lambda$,
which itself is bounded independent of $\lambda$ and $t$, by say
$C'>0$.


Next, let us deal with the singularities of the pulled back metric on $A$
caused by the interior branch points. As the next claim shows, it is an easy
matter to deform the metric slightly near each branch point without increasing
the area much and without changing the upper bound on the Gaussian curvature
too much. Here is the result:

\begin{claim}[Smoothing the metric near interior branch points]
\label{claim:branch-point-metric-deformation}
Let $\psi\colon A\subset M\times S^1_\lambda$ be an area-minimizing
annulus of area at most $\mu$ with smoothly embedded boundary as
constructed above. Let $h$ be the induced (possibly singular) metric
on $A$ induced by pulling back $g\times ds^2$ by $\psi$, and let
$C''>0$ be an upper bound on the Gaussian curvature of $h$ (away
from the branch points). Then  there is a deformation $\tilde h$ of
$h$, supported near the interior branch points, to a smooth metric
with the property that the area of the deformed smooth metric is at
most $2\mu$ and where the upper bound for the curvature of $\tilde
h$ is $2C''$.
\end{claim}


\begin{proof} Fix an interior branch point $p$.
 Since $\psi$  is smooth and
conformal onto its image, there is a disk in $A$ centered at $p$ in
which $h=f(z,\bar z)|dz|^2$ for a smooth function $f$ on the disk.
The function $f$ vanishes at the origin and is positive on the
complement of the origin. Direct computation shows that the Gaussian
curvature $K(h)$ of $h$ in this disk is given by
$$K(h)=\frac{-\triangle f}{2f^2}+\frac{|\nabla f|^2}{2f^3}\le C,$$
where $\triangle$ is the usual Euclidean Laplacian on the disk and $|\nabla
f|^2=(\partial f/\partial x)^2+(\partial f/\partial y)^2$.  Now consider the
metric $(f+\epsilon)|dz|^2$ on the disk. Its Gaussian curvature is
$$\frac{-\triangle f}{2(f+\epsilon)^2}+\frac{|\nabla f|^2}{2(f+\epsilon)^3}.$$
\begin{claim}[Curvature bound for the perturbed conformal metric]
\label{claim:perturbed-metric-curvature-bound}
For all $\epsilon>0$ the Gaussian curvature of $(f+\epsilon)|dz|^2$
is at most $2C''$. \end{claim}

\begin{proof}
We see that $-\triangle f\le C''f^2$, so that
\begin{eqnarray*}\frac{-\triangle
f}{(f+\epsilon)^2}+\frac{|\nabla f|^2}{(f+\epsilon)^3} & = &
\frac{(f+\epsilon)(-\triangle
f)+|\nabla f|^2}{(f+\epsilon)^3} \\
& \le & \frac{C''f^3-\epsilon\triangle f}{(f+\epsilon)^3}\le
C''+\frac{\epsilon f^2C''}{(f+\epsilon)^3}\le 2C''.
\end{eqnarray*}
\end{proof}

Now we fix a smooth function $\rho(r)$ which is identically one on a subdisk
$D'$ of $D$ and vanishes near $\partial D$ and we replace the metric $h$ on the
disk by
$$h_\epsilon=(f+\epsilon \rho(r))|dz|^2.$$ The above computation shows that the
Gaussian curvature of $h_\epsilon$ on $D'$ is bounded above by
$2C''$. As $\epsilon$ tends to zero the restriction of the metric
$h_\epsilon$ to $D\setminus D'$ converges uniformly in the
$C^\infty$-topology to $h$. Thus, for all $\epsilon>0$ sufficiently
small the Gaussian curvature of $h_\epsilon$ on $D\setminus D'$ will
also be bounded by $2C''$. Clearly, as $\epsilon$ tends to zero the
area of the metric $h_\epsilon$ on $D$ tends to the area of $h$ on
$D$.

Performing this construction near each of the finite number of interior branch
points and taking $\epsilon$ sufficiently small gives the perturbation $\tilde
h$ as required.
\end{proof}


Thus, if $\gamma$ and $\hat \gamma$ are ramps as in
Lemma~\ref{lem:ramp-length-comparison-small-annulus}, then replacing $\hat \gamma$ by a close $C^2$
approximation we have an abstract smooth annulus with a Riemannian
metric connecting $\gamma$ and $\hat \gamma$. Taking limits shows
that  establishing the conclusion of Lemma~\ref{lem:ramp-length-comparison-small-annulus} for a
sequence of better and better approximations to $\hat \gamma$ will
also establish it for $\hat \gamma$. This allows us to assume that
$\gamma$ and $\hat \gamma$ are disjoint. The area of this annulus is
bounded above by a constant arbitrarily close to $\mu$. The Gaussian
curvature of the Riemannian metric is bounded above by a constant
depending only on the curvature bounds of the ambient Ricci flow.
Finally, the integral of the absolute value of the geodesic
curvature over any interval of length $r$ of $\gamma$ is at most
$\delta$.

With all these preliminary remarks, we see that Lemma~\ref{lem:ramp-length-comparison-small-annulus}
follows from:


\begin{proposition}[Boundary length estimate for annuli of small area]
\label{prop:annulus-boundary-length-estimate}
\label{annlengthest}
Fix $0<\delta<1/100$.  For each $0<r$ and $C''<\infty$ there is a
$\mu>0$ such that the following holds. Suppose that $A$ is an
annulus with boundary components $c_0$ and $c_1$. Denote by $l(c_0)$
and $l(c_1)$ the lengths of $c_0$ and $c_1$, respectively. Suppose
that the Gaussian curvature of $A$ is bounded above by $C''$.
Suppose that  $l(c_0)>r$ and that for each sub-interval $I$ of $c_0$
of length $r$, the integral of the absolute value of the geodesic
curvature along $I$ is less than $\delta$. Suppose that the area of
$A$ is less than $\mu$. Then
$$l(c_1)\ge \frac{3}{4}l(c_0).$$
\end{proposition}

To us, this statement was intuitively extremely reasonable but we
could not find a result along these lines stated in the literature.
Also, in the end, the argument we constructed is quite involved,
though elementary.

The intuition is that we exponentiate in from the boundary component
$c_0$ using the family of geodesics perpendicular to the boundary.
The bounds on the Gaussian curvature and local bounds on the
geodesic curvature of $c_0$ imply that the exponential mapping will
be an immersion out to some fixed distance $\delta$ or until the
geodesics  meet the other boundary, whichever comes first.
Furthermore, the metric induced by  this immersion will be close to
the product metric. Thus, if there is not much area, it must be the
case that, in the measure sense, most of the geodesics in this
family must meet the other boundary before distance $\delta$. One
then deduces the length inequality. There are two main difficulties
with this argument that must be dealt with. The first is due to the
fact that we do not have a pointwise bound on the geodesic curvature
of $c_0$, only an integral bound of the absolute value over all
curves of short length. There may be points of arbitrarily high
geodesic curvature. Of course, the length of the boundary where the
geodesic curvature is large is very small. On these small intervals
the exponential mapping will not be an immersion out to any fixed
distance. We could of course, simply omit these regions from
consideration and work on the complement. But these small regions of
high geodesic curvature on the boundary can cause focusing (i.e.,
crossing of the nearby geodesics). We must estimate out to what
length along the boundary this happens. Our first impression was
that the length along the boundary where focusing occurred would be
bounded in terms of the total turning along the arc in $c_0$. We
were not able to establish this. Rather we found a weaker estimate
where this focusing length is bounded in terms of the total turning
 and the area bounded by the triangle cut out by the two geodesics
that meet. This is a strong enough result for our application. Since
the area is small and the turning on any interval of length $r$ is
small, a maximal collection of focusing regions will meet each
interval of length $r$ in $c_0$ in a subset of small total length.
Thus, on the complement (which is most of the length of $c_0$) the
exponential mapping will be an immersion out to length $\delta$ and
will be an embedding when restricted to each interval of length one.
The second issue to face is to show that the exponential mapping on
this set is in fact an embedding, not just an immersion. Here one
uses standard arguments invoking the Gauss-Bonnet theorem to rule
out various types of pathologies, e.g., that the individual
geodesics are not embedded or geodesics that end on $c_0$ rather
than $c_1$, etc. Once these are ruled out, one has established that
the exponential map on this subset is an embedding and the argument
finishes as indicated above.

\subsection{First reductions}


Of course, if the hypothesis of the proposition holds for $r>0$ then
it holds for any $0<r'<r$. This allows us to assume that $r<\mathrm{min}((C'')^{-1/2},1)$. Now let us scale the metric by $4r^{-2}$. The
area of $A$ with the rescaled metric is  $4r^{-2}$ times the area of
$A$ with the original metric. The Gaussian curvature of $A$ with the
rescaled metric is less than $(r^2C''/4)\le 1$. Furthermore, in the
rescaled metric $c_0$ has length greater than $2$ and the total
curvature along any interval of length $1$ in $c_0$ is at most
$\delta$. This allows us to assume (as we shall) that  $r=1$, that
$C''\le 1$, and that $l(c_0)\ge 2$. We must find a $\mu>0$ such that
the proposition holds provided that the area of the annulus is less
than $\mu$.



The function $k_\mathrm{geod}\colon c_0\to \Ar$ is smooth. We choose a
regular value $\alpha$ for $k_\mathrm{geod}$ with $1<\alpha<1.1$. In
this way we divide $c_0$ into two disjoint subsets, $Y$ where
$k_\mathrm{geod}>\alpha$, and $X$ where $k_\mathrm{geod}\le \alpha$. The
subset $Y$ is a union of finitely many disjoint open intervals and
$X$ is a disjoint union of finitely many closed intervals.

\begin{remark}[Length bound on the high-curvature subset Y]
\label{rem:high-curvature-set-length-bound}
\label{Ylengthrem}
The condition on $k_\mathrm{geod}$ implies that for any arc $J$ in
$c_0$ of length $1$ the total
 length of $J\cap Y$ is less than $\delta$.
\end{remark}



Fix $\delta'>0$. For each $x\in X$ there is a geodesic $D_x$ in $A$
whose initial point is $x$ and whose initial direction is orthogonal
to $c_0$. Let $f(x)$ be the minimum of $\delta$ and the distance
along $D_x$ to the first point (excluding $x$) of its intersection
with $\partial A$. We set
$$S_X(\delta') =\{(x,t)\in X\times [0,\delta']\bigl|\bigr.\  t\le f(x)\}.$$
The subset $S_X(\delta')$ inherits a Riemannian metric from the
product of the  metric on $X$ induced by the embedding $X\subset
c_0$ and the standard metric on the interval $[0,\delta']$.


\begin{claim}[Exponential map is a local diffeomorphism near c_0]
\label{claim:exponential-map-local-diffeomorphism}
\label{expcomp}
There is $\delta'>0$ such that the following holds.
 The
exponential mapping defines a map $\exp\colon S_X(\delta')\to
A$ which is a local diffeomorphism and the pullback of the metric on
$A$ defines a metric on $S_X(\delta')$ which is at least
$(1-\delta)^2$ times the given product metric.
\end{claim}


\begin{proof}
This is a standard computation using the Gaussian curvature upper
bound and the geodesic curvature bound.
\end{proof}

Now we fix $0<\delta'<1/10$ so that Claim~\ref{claim:exponential-map-local-diffeomorphism} holds, and
we set $S_X=S_X(\delta')$. We define
$$\partial_+S_X=\{(x,t)\in S_X \bigl|\bigr.\, t=f(x)\}.$$
Then the boundary of $S_X$ is made up of $X$, the arcs $\{x\}\times
[0,f(x)]$ for $x\in \partial X$ and $\partial_+(S_X)$. For any
subset $Z\subset X$ we denote by $S_Z$ the intersection $(Z\times
[0,\delta])\cap S_X$, and we denote by $\partial_+(S_Z)$ the
intersection of $S_Z\cap \partial_+S_X$.

Lastly, we fix $\mu>0$ with $\mu<(1-\delta)^2(\delta')/10$. Notice
that this implies that $\mu<1/100$. We now assume that the area of
$A$ is less than this value of $\mu$ (and recall that $r=1$, $C''=1$
and $l(c_0)\ge 2$). We must show that $l(c_1)>3l(c_0)/4$.



\subsection{Focusing triangles}

By a {\em focusing triangle} we mean the following. We have distinct
points $x,y\in X$ and sub-geodesics $D'_x\subset D_x$ and
$D'_y\subset D_y$ that are embedded arcs with $x$, respectively $y$,
as an endpoint. The intersection $D'_x\cap D'_y$ is a single point
which is the other endpoint of each of $D'_x$ and $D'_y$. Notice
that since $D'_x\subset D_x$ and $D'_y\subset D_y$, by construction
both $D'_x$ and $D'_y$ have lengths at most $\delta'$. We have an
arc $\xi$ in $c_0$ with endpoints $x$ and $y$ and the loop
$\xi*D'_y*(D'_x)^{-1}$ bounds a disk $B$ in $A$. The arc $\xi$ is
called the {\em base} of the focusing triangle and with, respect to
an orientation of $c_0$, if $x$ is the initial point of $\xi$ then
$D'_x$ is called the {\em left-hand side} of the focusing triangle
and $D'_y$ is called its {\em right-hand side}.

Our main goal here is the following lemma which gives an upper bound
for the length of the base, $\xi$, of a focusing triangle in terms
of the turning along the base and the area of the region $B$
enclosed by the triangle.

\begin{lemma}[Base-length bound for a focusing triangle]
\label{lem:focusing-triangle-base-length-bound}
\label{laest}
 Suppose that we have a focusing triangle ${\mathcal T}$ with base $\xi$ bounding a disk $B$
 in $A$.
 Suppose that the length of $\xi$ is at most one.
 Then
$$l(\xi)\le \left(\int_\xi k_\mathrm{geod}ds+\Area(B)\right).$$
\end{lemma}


\begin{proof}
\uses{claim:focusing-triangle-angle-curvature-bound,claim:focusing-triangle-vertex-geodesic}
We begin with a preliminary computation. We denote by $a(B)$ the
area of $B$. We define
$$t_{\xi}=\int_{\xi}k_\mathrm{geod}ds \ \ \ \text{and} \ \ \ T_\xi=\int_{\xi}|k_\mathrm{geod}|ds.$$

Recall that given a piecewise smooth curve, its {\em total turning}
is the integral of the geodesic curvature over the smooth part of
the boundary plus the sum over the break points of $\pi$ minus the
interior angle at the break point.  The Gauss-Bonnet theorem tells
us that for any compact surface with piecewise smooth boundary the
integral of the Gaussian curvature over the interior of the surface
plus the total turning around the boundary equals $2\pi$ times the
Euler characteristic of the surface.



\begin{claim}[Angle and curvature bounds for a focusing triangle]
\label{claim:focusing-triangle-angle-curvature-bound}
\label{theta/Kbd}
The angle $\theta_B$ between $D'_x$ and $D'_y$ at the vertex $v$
satisfies $$\theta_B\le t_\xi+a(B)$$ and for any measurable subset
$B'\subset B$ we have
$$\theta_B-t_\xi-a(B)\le \int_{B'}Kda<a(B).$$
\end{claim}

\begin{proof}
 Since $D'_x$ and
$D'_y$ meet $\partial A$ in right angles, the total turning around
the boundary of $B$  is
$$t_\xi+2\pi-\theta_B.$$
Thus, by Gauss-Bonnet, we have
$$\theta_B=\int_BKda+ t_\xi.$$
But $K\le 1$, giving the first stated inequality. On the other hand
$\int_BKda=\int_BK^+da+\int_BK^-da$, where $K^+=\max(K,0)$ and
$K^-=K-K^+$. Since $0\le \int_BK^+da\le a(B)$ and $\int_BK^-\le 0$,
the second string of inequalities follows.
\end{proof}

In order to make the computation we need to know that this triangle
is the image under the exponential mapping of a spray of geodesics
out of the vertex $v$. Establishing that requires some work.


\begin{claim}[Existence of a shortest geodesic to the focusing-triangle vertex]
\label{claim:focusing-triangle-vertex-geodesic}
\label{avgeo}
Let $a\in \mathrm{int}\,\xi$. There is a shortest path in $B$ from $a$
to $v$. This shortest path is a geodesic meeting $\partial B$ only
in its end points. It has length $\le(1/2)+\delta'$.
\end{claim}

\begin{proof}
The length estimate is obvious: Since $\xi$ has length at most $1$,
a path along $\partial A\cap B $ from $a$ to the closest of $x$ and
$y$ has length at most $1/2$. The corresponding side has length at
most $\delta'$. Thus, there is a path from $a$ to $v$ in $B$ of
length at most $(1/2)+\delta'$.


Standard convergence arguments show that there is a shortest path in
$B$ from $a$ to $v$. Fix $a\in \mathrm{int}(\partial A\cap B)$. It is
clear that the shortest path cannot meet either of the  `sides'
$D'_x$ and $D'_y$ at any point other than $v$. If it did, then there
would be an angle at this point and a local shortcut, cutting off a
small piece of the angle, would provide a shorter path. We must rule
out that the shortest path from $a$ to $v$ meets $\partial A\cap B$
in another point. If it does, let $a'$ be the last such point
(parameterizing the geodesic starting at $a$). The shortest path
from $a$ then leaves $\partial A$ at $a'$ in the direction tangent
at $a'$ to $\partial A$. (Otherwise, we would have an angle which
would allow us to shorten the path just as before.) This means that
we have a geodesic $\gamma$ from $v$ to $a'$ whose interior is
contained in the interior of $B$ and which is tangent to $\partial
A$ at $a'$. We label the endpoints of $\partial A\cap B$ so that the
union of $\gamma$ and the interval on $\partial A\cap B$ from $a'$
to $y$ gives a $C^1$-curve. Consider the disc $B'$ bounded
 by $\gamma$, the arc of $\partial A$ from $a'$
to $y$, and $D'_y$. The total turning around the boundary is at most
$3\pi/2+\delta$, and the integral of the Gaussian curvature over
$B'$ is at most the area of $B$, which is less than
$\mu<1/20<(\pi/4)-\delta$. This contradicts the Gauss-Bonnet
theorem.
\end{proof}

\begin{claim}[Uniqueness of the minimal geodesic to the vertex]
\label{claim:focusing-triangle-vertex-geodesic-unique}
For any $a\in (\partial A\cap B)$ there is a unique minimal geodesic
in $B$ from $a$ to $v$.
\end{claim}


\begin{proof}
Suppose not; suppose there are two $\gamma$ and $\gamma'$ from $v$
to $a$. Since they are both minimal in $B$, each is embedded, and
they must be disjoint except for their endpoints. The upper bound on
the curvature and the Gauss-Bonnet theorem implies that the angles
that they make at each endpoint are less than $\mu<\pi/2$. Thus,
there is a spray of geodesics (i.e. geodesics determined by an
interval $\beta$ in the circle of directions at $v$) coming out of
$v$ and moving into $B$ with extremal members of the spray being
$\gamma$ and $\gamma'$. The geodesics $\gamma$ and $\gamma'$ have
length at most $(1/2)+\delta'$, and hence the exponential mapping
from $v$ is a local diffeomorphism on all geodesics of length at
most the length of $\gamma$. Since the angle they make at $a$ is
less than $\pi/2$ and since the exponential mapping is a local
diffeomorphism near $\gamma$, as we move in from the $\gamma$ end of
the spray we find geodesics from $v$ of length less than the length
of $\gamma$ ending on points
 of $\gamma'$. The same
Gauss-Bonnet argument shows that the angles that each of these
shorter geodesics makes with $\gamma'$ is at most $\mu$. Consider
the subset $\beta'$ of $\beta$ which are directions of geodesics in
$B$ of length $<(1/2)+\delta'$ that end on points of $\gamma'$ and
make an angle less than $\mu$ with $\gamma'$. We have just seen that
$\beta'$ contains an open neighborhood of the end of $\beta$
corresponding to $\gamma$. Since the Gaussian curvature is bounded
above by $1$, and these geodesics all have length at most
$1/2+\delta$, it follows that the exponential map is a local
diffeomorphism near all such geodesics. Thus, $\beta'$ is an open
subset of $\beta$. On the other, hand if the direction of
$\gamma''\not=\gamma'$ is a point $b''\in \beta$ which is an
endpoint of an open interval  $ \beta'$, and if this interval
separates $b''$ from the direction of $\gamma$ then the length of
$\gamma''$ is less than the length of each point in the interval.
Hence, the length of $\gamma''$ is less than $(1/2)+\delta'$.
Invoking Gauss-Bonnet again we see that the angle between $\gamma''$
and $\gamma'$ is $<\mu$.

This proves that if $U$ is an open  interval in $\beta'$ then the
endpoint of $U$ closest to the direction of $\gamma'$ is also
contained in $\beta'$ (unless that endpoint is the direction of
$\gamma'$). It is now elementary to see that $\beta'$ is all of
$\beta$ except the endpoint corresponding to $\gamma'$. But this is
impossible. Since the exponential mapping is a local diffeomorphism
out to distance $(1/2)+\delta'$, and since $\gamma'$ is embedded,
any geodesic from $v$ whose initial direction is sufficiently close
to that of $\gamma'$ and whose length is at most $(1/2)+\delta'$
will not cross $\gamma'$.
\end{proof}

\begin{remark}[Uniqueness of the embedded geodesic from the vertex]
\label{rem:focusing-triangle-vertex-geodesic-embedded}
The same argument shows that from any $a\in (\partial A\cap B)$
there is a unique embedded geodesic in $B$ from $v$ to $a$ with
length at most $(1/2)+\delta'$. (Such geodesics may cross more than
once, but the argument given in the lemma applies to sub-geodesics
from $v$ to the first point of intersection along $\gamma$.)
\end{remark}


Let $E$ be the sub-interval of the circle of tangent directions at
$v$ consisting of all tangent directions of geodesics pointing into
$B$ at $v$. The endpoints of $E$ are the tangent directions for
$D'_x$ and $D'_y$. We define a function from $\xi$ to the interval
$E$ by assigning to each $a\in \xi$ the direction at $v$ of the
unique minimal geodesic in $B$ from $v$ to $a$. Since the minimal
geodesic is unique, this function is continuous and, by the above
remark, associates to $x$ and $y$ the endpoints of $E$. Since
geodesics are determined by their initial directions, this function
is one-to-one. Hence it is a homeomorphism from $\xi$ to $E$. That
is to say the spray of geodesics coming out of $v$ determined by the
interval $E$ produces a diffeomorphism between a wedge-shaped subset
of the tangent space at $v$ and $B$. Each of the geodesics in
question ends when it meets $\xi$.


Now that we have shown that the region enclosed by the triangle is
the image under the exponential map from the vertex $v$ of a
wedge-shaped region in the tangent space at $v$, we can make the
usual computation relating length and geodesic curvature. To do this
we pull back to the tangent space at $v$, and, using polar
coordinates, we write $\xi$ as $\{s=h(\psi); \psi\in E\}$ where $s$
is the radial coordinate and $\psi$ is the angular coordinate.
Notice that $h(\psi)\le (1/2)+\delta'$ for all $\psi\in E$. (In
fact, because the angles of intersection at the boundary are all
close to $\pi/2$ we can give a much better estimate on $h$ but we do
not need it.)
 We consider the one-parameter family of arcs $\lambda(t)$ defined
 to be the graph of the function $t\mapsto s(t)=th(\psi)$, for $0\le t\le
 1$. We set $l(t)$ equal to the length of
 $\lambda(t)$.

 \begin{claim}[Derivative bound for the length of the radial family]
\label{claim:arc-length-derivative-bound}
$$\frac{dl}{dt}(t)\le \max_{\psi\in E}h(\psi)\int_{\lambda(t)}k_\mathrm{geod}ds.$$
 \end{claim}

\begin{proof}
First of all notice that, by construction,  the curve $\xi$, which
is defined by $\{s=h(\psi)\}$, is orthogonal to the radial geodesics
to the endpoints. As a consequence, $h'(\psi)=0$ at the endpoints.
Thus, each of the curves $\lambda(t)$ is orthogonal to the radial
geodesics through its end points. Therefore, as we vary the family
$\lambda(t)$ the formula for the derivative of the length is
$$l'(t)=\int_{\lambda(t)}k_\mathrm{geod}(\psi)h(\psi)|\mathrm{cos}(\theta(\psi,t))|ds$$ where $\theta(\psi,t)$ is the angle at
$(th(\psi),\psi)$ between the curve $s=th(\psi)$ and the radial
geodesic. The result follows immediately.
\end{proof}

Next, we must bound the turning of $\lambda(t)$. For this we invoke the
Gauss-Bonnet theorem once again. Applying this to the wedge-shaped disk $W(t)$
cut out by $\lambda(t)$ gives
$$\int_{W(t)}Kda+\int_{\lambda(t)}k_\mathrm{geod}ds=\theta_B.$$
From Claim~\ref{claim:focusing-triangle-angle-curvature-bound} we conclude that
$$\int_{\lambda(t)}k_\mathrm{geod}ds\le t_\xi+a(B).$$


Of course, by Claim~\ref{claim:focusing-triangle-vertex-geodesic} we have $\max_{\psi\in
E}h(\psi)\le (1/2)+\delta'$. Since $l(0)=0$, this implies that
$$l(\xi)=l(1)\le (a(B)+t_B)((1/2)+\delta')<a(B)+t_B.$$
This completes the proof of Lemma~\ref{lem:focusing-triangle-base-length-bound}.
\end{proof}

\begin{corollary}[Bound on total base length of a disjoint collection of focusing triangles]
\label{cor:focusing-triangle-base-bound-collection}
\label{muepsbd}
Suppose that ${\mathcal T}$ is a focusing triangle with base $\xi$
of length at most one. Then the length of $\xi$ is at most
$\delta+\mu$. More generally, suppose we have a collection of
focusing triangles ${\mathcal T}_1,\ldots,{\mathcal T}_n$ whose
bases all lie in a fixed interval of length one in $c_0$. Suppose
also that the interiors of disks bounded by these focusing triangles
are disjoint. Then the sum of the lengths of the bases is at  most
$\delta+\mu$.
\end{corollary}

\begin{proof}
\uses{lem:focusing-triangle-base-length-bound}
The first statement is immediate from the previous lemma. The second
comes from the fact that the sum of the areas of the disks bounded
by the ${\mathcal T}_i$ is at most $\mu$ and the sum of the total
turnings of the $\xi_i$ is at most $\delta$.
\end{proof}


This completes our work on the local focusing issue. It remains to
deal with global pathologies that would prevent the exponential
mapping from being an embedding out to distance $\delta'$.

\subsection{No $D_x$ is an embedded arc with both ends in $c_0$}

One thing that we must show  is that the geodesics $D_x$ are
embedded. Here is a special case that will serve some of our
purposes.


\begin{lemma}[No sub-geodesic of D_x is a homotopically trivial embedded loop]
\label{lem:geodesic-no-trivial-embedded-loop}
\label{noloops}
For each $x\in X$, there is no non-trivial sub-geodesic of $D_x$
which is a homotopically trivial embedded loop in $A$.
\end{lemma}


\begin{proof}
Were there such a loop, its  total turning would be $\pi$ minus the
angle it makes when the endpoints of the arc meet. Since $K\le 1$
and the area of the disk bounded by this loop is less than the area
of $A$ which in turn is less than $\mu<\pi$,  one obtains a
contradiction to the Gauss-Bonnet theorem.
\end{proof}

Next, we rule out the possibility that  one of the geodesics $D_x$
has both endpoints contained in $c_0$. This is the main result of
this section. In a sense, what the argument we give here shows that
if there is a $D_x$ with both ends on $c_0$, then under the
assumption of small area, $D_x$ cuts off a thin tentacle of the
annulus. But out near the end of this thin tentacle there must be a
short arc with large total turning, violating our hypothesis on the
integrals of the geodesic curvature over arcs of length at most one.


\begin{lemma}[No D_x spans between two points of c_0]
\label{lem:no-geodesic-spans-c0-to-c0}
\label{Dxacross}
There is no $x$ for which $D_x$ is an embedded arc  with both
endpoints on $c_0$ and otherwise disjoint from $\partial A$.
\end{lemma}

\begin{proof}
\uses{lem:geodesic-no-trivial-embedded-loop,claim:local-embedding-of-geodesic-family,claim:focusing-triangle-limit-existence,cor:focusing-triangle-base-bound-collection,rem:high-curvature-set-length-bound}
Suppose that there were such a $D_x$. Then $D_x$ separates $A$ into
two components, one of which, $B$, is a topological disk. Let $c_0'$
be the intersection of $c_0$ with $B$. We consider two cases: Case
(i): $l(c'_0)\le 1$ and Case (ii): $l(c_0')>1$.

Let us show that the first case is not possible. Since $D_x$ is a
geodesic and $D_x$ is perpendicular to $c_0$ at one end, the total
turning around the boundary of $B$ is at most
$$3\pi/2+\int_{c_0'}k_\mathrm{geod}ds<3\pi/2+\delta,$$
where the last inequality uses the fact that the length of $c_0'$ is
at most one. On the other hand, $\int_{B}Kda<\mu$, and
$\mu<1/20<(\pi/2)-\delta$. This contradicts the Gauss-Bonnet
theorem.


Now let us consider the second case. Let $J$ be the subinterval of
$c_0'$ with one end point being $x$ and with the length of $J$ being
one. We orient $J$ so that $x$ is its initial point. We set
$X_J=J\cap X$. We define $S_{X_J}(B)\subset S_{X_J}$ as follows. For
each $y\in X_J$ we let $f_B(y)$ be the minimum of $\delta'$ and the
distance along $D_y$ to the first point (excluding $y$) of $D_y$
contained in $\partial B$ and let $D_y(B)$ be the sub-geodesic of
$D_y$ of this length starting at $y$. Then $S_{X_J}(B)\subset
S_{X_J}$ is the union over $y\in X_J$ of $[0,f_B(y)]$.  Clearly, the
exponential mapping defines an immersion of $S_{X_J}(B)$ into $B$.
We need to replace $X_ J$ by a slightly smaller subset in order to
make the exponential mapping be an embedding. To do this we shall
remove bases of a maximal focusing triangles in $B$.

First notice that for each $y\in  X_J$ the exponential mapping is an
embedding on $D_y(B)$. The reason is that the image of $D_y(B)$ is a
geodesic contained in the ball $B$. Lemma~\ref{lem:geodesic-no-trivial-embedded-loop} then shows
that this geodesic is embedded. This leads to:

\begin{claim}[Local embedding of the geodesic family near a component of X_J]
\label{claim:local-embedding-of-geodesic-family}
\label{locemb}
For any component $c$ of $X_J$, the restriction of the exponential
mapping to $S_c(B)=\left(c\times [0,\delta')\right)\cap S_{X_J}(B)$
is an embedding.
\end{claim}

\begin{proof}
Since the geodesics that make up $S_c(B)$ have length at most
$\delta'<1/10$ and since the curvature of the annulus is bounded
above by $1$, the restriction of the exponential mapping to $S_c(B)$
is a local diffeomorphism. The restriction  to each $\{y\}\times
[0,f_B(y)]$ is an embedding onto $D_y(B)$. If the restriction of the
exponential mapping to $S_c(B)$ is not an embedding, then there are
$y\not= y'$ in $c$ such that the geodesics $D_y(B)$ and $D_{y'}(B)$
meet. When they meet, they meet at a positive angle and by the
Gauss-Bonnet theorem this angle is less than $\mu+\delta$.  Thus,
all the geodesics starting at points sufficiently close to $y'$ and
between $y$ and $y'$ along $c$ must also meet $D_y(B)$. Of course,
if a sequence of $D_{y_i}(B)$ meet $D_y(B)$, then the same is true
for the limit. It now follows that $D_{y''}(B)$ meets $D_y(B)$ for
all $y''$ between $y$ and $y'$. This contradicts the fact that
$D_y(B)$ is embedded.
\end{proof}




\begin{claim}[Focusing triangles contain components of J minus X_J and admit limits]
\label{claim:focusing-triangle-limit-existence}
\label{limittri}
Any focusing triangle for $J$ must contain a component of $J\setminus X_J $. If
$\{{\mathcal T}_n\}$ is an infinite sequence of focus triangles for $J$, then,
after passing to a subsequence, there is a limiting focusing triangle for $J$.
\end{claim}

\begin{proof}
\uses{claim:local-embedding-of-geodesic-family}
The first statement is immediate from Claim~\ref{claim:local-embedding-of-geodesic-family}.
 Since $X\cap J$ is compact, it is clear that after passing to a
 subsequence each of the sequence of left-hand sides and the
 sequence of right-hand sides  converge to a geodesic arc orthogonal to
 $J$ at points of $X$. Furthermore, these limiting geodesics meet in
 a point at distance at most $\delta'$ from the end of each. The only
 thing remaining to show is that the limiting left- and right-hand
 sides do not begin at the same point of $X$.
This is clear since each focusing triangle contains one of the finitely many
components of $J\setminus X_J$.
\end{proof}

Using Claim~\ref{claim:focusing-triangle-limit-existence} we see that if there is  a  focusing
triangle for $J$ there is a first point $x_1$ in $ X_J$ whose
associated geodesic contains the left-hand side of a focusing
triangle for $J$. Then since the base length of any focusing
triangle is bounded by a fixed constant, invoking again
Claim~\ref{claim:focusing-triangle-limit-existence}, that there is a focusing triangle ${\mathcal
T}_1$ for $J$ that has left-hand side contained in the geodesic
$D_{x_1}$ and has a maximal base among all such focusing triangles,
Maximal in the sense that the base of this focusing triangle
contains the base of any other focusing triangle with left-hand side
contained in $D_{x_1}$. Denote its base by $\xi_1$ and denote the
right-hand endpoint of $\xi_1$ by $y_1$. For the triangle we take
the geodesic arcs to the first point of intersection measured along
$D_{y_1}$. Set $J_1=\overline{J\setminus \xi}$, and repeat the
process for $J_1$. If there is a focusing triangle for $J_1$ we find
the first left-hand side of such focusing triangles and then find
the maximal focusing triangle ${\mathcal T}_2$ with this left-hand
side.

\begin{claim}[Disjointness of successive maximal focusing triangles]
\label{claim:maximal-focusing-triangles-disjoint}
The interior of ${\mathcal T}_2$ is disjoint from the interior of
${\mathcal T}_1$.
\end{claim}

\begin{proof}
Since by construction the interiors of the bases of ${\mathcal T}_1$
and ${\mathcal T}_2$ are disjoint, if the interior of ${\mathcal
T}_2$ meets ${\mathcal T}_1$, then one of the sides of ${\mathcal
T}_2$ crosses the interior of one of the sides of ${\mathcal T}_1$.
But since ${\mathcal T}_1$ is a maximal focusing triangle with its
left-hand side, neither of the sides of ${\mathcal T}_2$ can cross
the interior of the left-hand side of ${\mathcal T}_1$. If one of
the sides of ${\mathcal T}_2$ crosses the interior of the right-hand
side of ${\mathcal T}_1$, then the right-hand side of ${\mathcal
T}_1$ is the left-hand side of a focusing triangle for $J_1$. Since
by construction the left-hand side of ${\mathcal T}_2$ is the first
such, this means that the left-hand side of ${\mathcal T}_2$ is the
right-hand side of ${\mathcal T}_1$. This means that the right-hand
side of ${\mathcal T}_2$ terminates when it meets the right-hand
side of ${\mathcal T}_1$ and hence the right-hand side of ${\mathcal
T}_2$ ends the first time that it meets the right-hand side of
${\mathcal T}_1$ and hence does not cross it.
\end{proof}

We continue in this way constructing focusing triangles for $J$ with
disjoint interiors. Since each focusing triangle for $J$ contains a
component of $J\setminus X_J$, and as there are only finitely many
such components, this process must terminate after a finite number
of steps. Let ${\mathcal T}_1,\ldots,{\mathcal T}_k$ be the focusing
triangles so constructed, and denote by $\xi_i$ the base of
${\mathcal T}_i$. Let $X_J'= X_J\setminus \cup_{i=1}^k\xi_i$.

\begin{definition}[Maximal set of focusing triangles]
\label{def:maximal-focusing-triangle-set}
We call the triangles ${\mathcal T}_1,\ldots,{\mathcal T}_k$
constructed above, the {\em maximal set of focusing triangles} for
$J$ relative to $B$.
\end{definition}



\begin{claim}[Length bound on the complement of the focusing-triangle bases]
\label{claim:remaining-set-length-bound}
The length of $X'_J$ is at least $1-2\delta-\mu$.
\end{claim}

\begin{proof}
\uses{cor:focusing-triangle-base-bound-collection,rem:high-curvature-set-length-bound}
Since the interiors of the ${\mathcal T}_i$ are disjoint, according
to Corollary~\ref{cor:focusing-triangle-base-bound-collection}, we have $\sum_il(\xi_i)<\delta+\mu$. We
also know by Remark~\ref{rem:high-curvature-set-length-bound} that the length of $X_J$ is at
least $(1-\delta)$. Putting these together gives the result.
\end{proof}

We define $S_{X'_J}(B)$ to be the intersection of $S_{X_J}(B)$ with
$S_{X'_J}$.

\begin{claim}[Embedding of the exponential map on the remaining set]
\label{claim:exponential-map-embedding-on-remainder}
The restriction of the exponential mapping to $S_{X'_J}(B)$ is an
embedding.
\end{claim}

\begin{proof}
Suppose that we have distinct points $x',y'$ in $X'_J$ such that
$D_{x'}(B)\cap D_{y'}(B)\not=\emptyset$. We assume that $x'<y'$ in
the orientation on $J$. Then there is a focusing triangle for $J$
whose base is the sub-arc of $J$ with endpoints $x'$ and $y'$, and
hence the left-hand side of the focusing triangle is contained in
$D_{x'}(B)$. Since $x'$ is not a point of $\cup_{i}\xi_i$ either it
lies between two of them, say $\xi_j$ and $\xi_{j+1}$ or it lies
between the initial point $x$ of $J$ and the initial point of
$\xi_1$ or it lies between the last $\xi_n$ and the final point of
$J$.

But $x'$ cannot lie before $\xi_1$, for this would contradict the
construction which took as the left-hand endpoint of $\xi_1$ the
first point of $J$ whose geodesic contained the left-hand side of a
focusing triangle for $J$. Similarly, $x'$ cannot lie between
$\xi_j$ and $\xi_{j+1}$ for any $j$ since the left-hand endpoint of
$\xi_{j+1}$ is the first point at or after the right-hand endpoint
of $\xi_j$ whose geodesic contains the left-hand side of a focusing
triangle for $J$. Lastly, $x'$ cannot lie to the right of the last
$\xi_k$, for then we would not have finished the inductive
construction.
\end{proof}

We pull back the metric of $A$ to the $S_{X'_J}(B)$ by the
exponential mapping. Since this pullback metric is at least
$(1-\delta)^2$ times the product of the metric on $X'_J$ induced
from $c_0$ and the usual metric on the interval, and since the map
on this subset is an embedding, we see that the area of the region
of the annulus which is the image under the exponential mapping of
this subset is at least
$$(1-\delta)^2\int_{X'_J}f_B(x)ds,$$
where $s$ is arc length along $X'_J$. Of course, the area of this
subset is at most $\mu$. This means that, setting $Z$ equal to the
subset of $ X'_J(\delta')$ given by
$$Z=\{z\in X'_J|f_B(z)<\delta'\},$$ the total length of $X'_J\setminus Z$
satisfies
$$l(X'_J\setminus Z)\le (1-\delta)^{-2}(\delta')^{-1}\mu<\frac{1}{10},$$
where the last inequality is an immediate consequence of our choice
of $\mu$. Thus, the length of $Z$ is at least $(0.9-2\delta-\mu)\ge
0.87$. Let $\partial_+S_Z(B)$ be the union of the final endpoints
(as opposed to the initial points) of the $D_x(B)$ as $x$ ranges
over $Z$.
 Of course, since $f_B(z)<\delta'$
for all $z\in Z$, it must be the case that  the exponential mapping
embeds $\partial_+S_Z(B)$ into $\partial B$. Furthermore, the total
length of the image of $\partial_+S_Z(B)$ is at least
$(1-\delta)l(Z)\ge 0.86$. The boundary of $B$ is made up of two
pieces: $D_x$ and an arc on $c_0$. But the length of $D_x$ is at
most $\delta'<1/20$ so that not all of $\partial_+S_Z(B)$ can be
contained in $D_x$. Thus, there is $z\in Z$, distinct from $x$ such
that $D_z$ has both endpoints in $c_0$. It then follows  that all
points of $Z$ that are separated (along $J$) from $x$ by $z$ have
the same property. Since the length of $Z$ is at least $0.86$, it
follows that there is a point $z\in X$ at least distance $0.85$
along $J$ from $x$ with the property that $D_z$ has both endpoints
in $c_0$. The complementary component of $D_z$ in $A$, denoted $B'$,
 is a disk that is contained in $B$ and the length of $B'\cap c_0$ is at
least $0.85$ less than the length of $B\cap c_0$.

The length of $B'\cap c_0$ cannot be less than $1$, for that gives a
contradiction. But if the length of $B'\cap c_0$ is greater than one, we now
repeat this construction replacing $B$ by $B'$. Continuing in this way we
eventually we cut down the length of $B\cap c_0$ to be less than one and hence
reach a contradiction.
\end{proof}


\subsection{For every $x\in X$, the geodesic $D_x$ is embedded}

The steps in the above argument inductively constructing disjoint
maximal focusing triangles and showing that their bases have a small
 length and that off of them the map is an embedding will be
repeated in two other contexts. The next context is to rule out the
case when a sub-arc of $D_x$  forms a homotopically non-trivial loop
in $A$.


\begin{lemma}[No sub-geodesic of D_x is an embedded loop in A]
\label{lem:geodesic-no-embedded-loop-general}
For any $x\in X$ there is no sub-geodesic of $D_x$ that is an
embedded loop in $A$.
\end{lemma}

\begin{proof}
\uses{lem:no-geodesic-spans-c0-to-c0}
We have already treated the case when the loop bounds a disk. Now we
need to treat the case when the loop is homotopically non-trivial in
$A$. Let $D'_x\subset D_x$ be the minimal compact sub-geodesic
containing $x$ that is not an embedded arc. Let $\mathrm{int}\,B$ be
the complementary component of $D'_x$ in $A$ that contains
$c_0\setminus\{x\}$. There is a natural compactification of $\mathrm{int}\, B$ as a disk and an immersion of this disk into $A$, an
immersion that is two-to-one along the shortest sub-geodesic of
$D'_x$ from $x$ to the point of intersection of $D_x$ with itself.
We do exactly the same construction as before. Take a sub-arc $J$ of
length one with $x$ as an endpoint and construct $S_{X_J}(B)$
consisting of the union of the sub-geodesics of $D_z$, for $z\in
J\cap X$ that do not cross the boundary of $B$. We then construct a
sequence of maximal focusing triangles along $J$ relative to $B$
just as in the previous case. In this way we construct a subset $Z$
of $X\cap J$ of total length at least $0.87$ with the property that
for every $z\in Z$ the final end of $D'_z(B)$ lies in $\partial B$.
Furthermore, the length of the arcs that these final ends sweep out
is at least $0.86$. Hence, since the total length of the part of the
boundary of $B$ coming from $D_x'$ is at most $2\delta'<0.2$, there
must be a $z\in Z$ for which $D_z(B)$ has both ends in $c_0$. This
puts us back in the case ruled out in Lemma~\ref{lem:no-geodesic-spans-c0-to-c0}.
\end{proof}

\subsection{Far apart $D_x$'s don't meet}

Now the last thing that can prevent the exponential mapping in the
complement of the focusing triangles from being an embedding is that
geodesics $D_x$ and $D_y$ might meet even though $x$ and $y$ are far
apart along $c_0$. Our next goal is to rule this out.

\begin{lemma}[Geodesics D_x, D_y from far-apart points do not meet]
\label{lem:far-apart-geodesics-disjoint}
Let $x,y$ be distinct points of $X$. Suppose that there are
sub-geodesics $D'_x\subset D_x$ and $D'_y\subset D_y$ with a common
endpoint. Then the arc $D'_x*(D'_y)^{-1}$ cuts $A$ into two
complementary components, exactly one of which is a disk, denoted
$B$. Then it is not possible for  $B\cap c_0$ to contain an arc of
length $1$.
\end{lemma}

\begin{proof}
\uses{lem:no-geodesic-spans-c0-to-c0}
The proof is exactly the same as in Lemma~\ref{lem:no-geodesic-spans-c0-to-c0} except that
the part of the boundary of $B$ that one wants to avoid has length
at most $2\delta'<0.2$ instead of $\delta'$. Still, since (in the
notation of the proof of Lemma~\ref{lem:no-geodesic-spans-c0-to-c0}) the total length of
$Z$ is at least $0.87$ so that the lengths of the other ends of the
$D_z$ as $z$ ranges over $Z$ is at least $0.86$, there is $z\in Z$
for which both ends of $D_z$ lie in $c_0$. Again this puts us back
in the case ruled out by Lemma~\ref{lem:no-geodesic-spans-c0-to-c0}.
\end{proof}


As a special case of this result we have the following.

\begin{corollary}[Short arcs with meeting geodesics bound a focusing disk]
\label{cor:focusing-triangle-from-short-arc}
\label{noholes}
Suppose that we have an arc $\xi$ of length at most $1$ on $c_0$.
Denote the endpoints of $\xi$ by $x$ and $y$ and  suppose that
$D_x\cap D_y\not=\emptyset$. Let $D'_x$ and $D'_y$ be sub-geodesics
containing $x$ and $y$ respectively ending at the same point, $v$,
and otherwise disjoint. Then the loop $\xi*D'_y*(D'_x)^{-1}$ bounds
a disk in $A$.
\end{corollary}


\begin{proof}
\uses{lem:no-geodesic-spans-c0-to-c0}
If not, then it is homotopically non-trivial in $A$ and replacing
$\xi$ by its complement, $c_0\setminus\mathrm{int}\,\xi$, gives us
exactly the situation of the previous lemma. (The length of
$c_0\setminus\mathrm{int}\,\xi$ is at least one since the length of
$c_0$ is at least $2$.)
\end{proof}


Let us now summarize what we have established so far about the intersections of
the geodesics $\{D_x\}_{x\in X}$.

\begin{corollary}[Summary of intersection properties of the geodesics D_x]
\label{cor:geodesic-family-intersection-summary}
\label{focussummary}
\uses{prop:annulus-boundary-length-estimate,cor:focusing-triangle-from-short-arc,lem:no-geodesic-spans-c0-to-c0}
For each $x\in X$, the geodesic $D_x$ is an embedded arc in $A$.
Either it has length $\delta'$ or its final point lies on $c_1$.
Suppose there are   $x\not=x'$ in $X$ with $D_x\cap
D_{x'}\not=\emptyset$. Then there is an arc $\xi$ on $c_0$
connecting $x$ to $x'$ with the length of $\xi$ at most
$\delta+\mu$. Furthermore, for sub-geodesics $D'_x\subset D_x$,
containing $x$, and $D'_{x'}\subset D_{x'}$, containing $x'$, that
intersect exactly in an endpoint of each, the loop
$\xi*D'_{x'}*(D'_x)^{-1}$ bounds a disk $B$ in $A$, and the length
of $\xi$ is at most the turning of $\xi$ plus the area of $B$.
\end{corollary}

\subsection{Completion of the proof}

We have now completed all the technical work on focusing and we have
also shown that  the restriction of the exponential mapping to the
complement of the bases of the focusing regions is an embedding. We
are now ready to complete the proof of
Proposition~\ref{prop:annulus-boundary-length-estimate}.


Let $J$ be an interval of length one in $c_0$. Because of
Corollary~\ref{cor:geodesic-family-intersection-summary} we can  construct the maximal focusing
triangles for $J$ as follows. Orient $J$, and begin at the initial
point of $J$. At each step we consider the first $x$ (in the
subinterval of $J$ under consideration) which intersects a $D_y$ for
some later $y\in J$. If we have such $y$, then we can construct the
sides of the putative triangle for sub-geodesics of $D_x$ and $D_y$.
But we need to know that we have a focusing triangle. This is the
content of Corollary~\ref{cor:focusing-triangle-from-short-arc}. The same reasoning works when we
construct the maximal such focusing triangle with a given left-hand
side, and then when we show that in the complement of the focusing
triangles the map is an embedding. Thus, as before, for an interval
$J$  of length $1$, we construct a subset $X_J'\subset X\cap J$ of
length at least $0.97$ such that the restriction of the exponential
mapping to $S_{X'_J}$ is an embedding. Again the area estimate shows
that there is a subset $Z\subset X'_J$ whose length is at least
$0.87$ with the property that for every $z\in Z$ the geodesic $D_z$
has both endpoints in $\partial A$. By Lemma~\ref{lem:no-geodesic-spans-c0-to-c0}, the
only possibility for the final endpoints of all these $D_z$'s is
that they lie in $c_1$.


In particular, there are $x\in X$ for which $D_x$ spans from $c_0$
to $c_1$. We pick one such, $x_0$, contained in the interior of $X$,
and use it as the starting point for a construction of maximal
focusing triangles all the way around $c_0$. What we are doing at
this point actually is cutting the annulus open along $D_{x_0}$ to
obtain a disk and we construct a maximal family of focusing
triangles of the interval $[x'_0,x_0'']$ obtained by cutting $c_0$
open at $x_0$ relative to this disk. Here $x_0'$ and $x_0''$ are the
points of the disk that map to $x_0$ when the disk is identified to
form $A$. Briefly, having constructed a maximal collection of
focusing triangles for a subinterval $[x'_0,x]$, we consider the
first point $y$ in the complementary interval $[x,x''_0]$ with the
property that there is $y'$ in this same interval, further along
with $D_y\cap D_{y'}\not=0$. Then, using
Corollary~\ref{cor:geodesic-family-intersection-summary} we construct the maximal focusing
triangle on $[x,x''_0]$ with left-hand side being a sub-geodesic of
$D_y$. We then continue the construction inductively until we reach
$x''_0$. Denote by $\xi_1,\ldots,\xi_k$ the bases of these focusing
triangles and let $X'$ be $X\setminus\cup_i\xi_i$.

The arguments above show that the exponential mapping is an
embedding of $S_{X'}$ to the annulus.




\begin{claim}[Total base length of focusing triangles meeting a unit interval]
\label{claim:focusing-triangle-total-base-length-bound}
For every subinterval $J$ of length one in $c_0$ the total length of
the bases $\xi_i$ that meet $J$ is at most $2\delta+\mu<0.03$.
\end{claim}

\begin{proof}
\uses{cor:focusing-triangle-base-bound-collection}
Since, by Corollary~\ref{cor:focusing-triangle-base-bound-collection}, every base of a focusing triangle
has length at most $\delta+\mu$, we see that the union of the bases
of focusing triangles meeting $J$ is contained in an interval of
length $1+2(\delta+\mu)<2$. Hence, the total turning of the bases of
these focusing triangles is at most $2\delta$ whereas the sum of
their areas is at most $\mu$. The result now follows from
Corollary~\ref{cor:focusing-triangle-base-bound-collection}.
\end{proof}

By hypothesis there is an integer $n\ge 1$ such that the length
$l(c_0)$ of $c_0$ is greater than $n$ but less than or equal to
$n+1$. Then it follows from the above that the total length of the
bases of all the focusing triangles in our family is at most
$$(n+1)(2\delta+\mu)< 0.03(n+1)\le 0.06n\le  0.06 l(c_0).$$


Since the restriction of the exponential mapping to $S_{X'}$ is an
embedding, it follows from Claim~\ref{claim:exponential-map-local-diffeomorphism} and the choice of
$\delta'$ that, for any open subset $Z$ of $X'$, the area of the
image under the exponential mapping of $S_Z$ is at least
$(1-\delta)^2\int_Zf(x)ds$, where $ds$ is the arc length along $Z$.
Also, the image under the exponential mapping of $\partial_+(S_Z)$
is an embedded arc in $A$ of length at least $(1-\delta)l(Z)$. Since
the length of $X'$ is at least $(0.94)l(c_0)$ and since the area of
$A$ is less than $\mu<(1-\delta)^2\delta'/10$, it follows that the
subset of $X'$ on which $f$ takes the value $\delta'$ has length at
most $0.10<(0.10)l(c_0)$. Hence, there is a subset $X''\subset X'$
of total length at least $(0.84)l(c_0)$ with the property that
$f(x)<\delta'$ for all $x\in X''$. This means that for every $x\in
X''$ the geodesic $D_x$ spans from $c_0$ to $c_1$, and hence the
exponential mapping embeds $\partial_+S_{X''}$ into $c_1$. But we
have just seen that the length of the image under the exponential
mapping of $\partial_+S_{X''}$  is at least
$$(1-\delta)l(X'')>(0.99)l(X'')>(0.83)l(c_0).$$
It follows  that the length of $c_1$ is at least
$(0.83)l(c_0)>3(l(c_0))/4$.

This completes the proof.


\section{Proof of the first inequality in Lemma~{\ref{lem:curve-shrinking-shi-estimate}}}\label{18.6}

Here is the statement that we wish to establish when the manifold
$(W,h(t))$ is the product of $(M,g(t))\times (S^1_\lambda,ds^2)$.


\begin{lemma}[First curvature bound in the curve-shrinking Shi estimate]
\label{lem:curve-shrinking-shi-estimate-first-inequality}
\label{csshi1}
Let $(W,h(t)),\ t_0\le t\le t_1$, be a Ricci flow and fix
$\Theta<\infty$. Then there exist constants $\delta>0$ and $0<r_0\le
1$   depending only on the curvature bound for the ambient Ricci
flow and $\Theta$ such that the following holds. Let $c(x,t),\
t_0\le t\le t_1$, be a curve-shrinking flow with $c(\cdot,t)$
immersed for each $t\in [t_0,t_1]$ and with the total curvature of
$c(\cdot,t_0)$ being at most $\Theta$. Suppose that there is $0<r\le
r_0$ and at a time $t'\in [t_0,t_1-\delta r^2]$ such that the length
of $c(\cdot,t')$ is at least $r$ and the total curvature of
$c(\cdot,t')$ on any sub-arc of length $r$ is at most $\delta$. Then
for every $t\in [t',t'+\delta r^2]$ the curvature $k$ satisfies
$$k^2\le  \frac{2}{(t-t')}.$$
\end{lemma}

The rest of this section is devoted to the proof of this lemma. In
\cite{AG} such a local estimate was established when the ambient
manifold was Euclidean space and the curve in question is a graph. A
related result for hypersurfaces that are graphs appears in
\cite{EckerHuisken}. The passage from Euclidean space to a general
Ricci flow is straightforward, but it is more delicate to use the
bound on total curvature on initial sub-arcs of length $r$ to show
that in appropriate coordinates the evolving curve can be written as
an evolving graph, so that the analysis in [2] can be applied.

We fix $\delta>0$ sufficiently small. We fix $t'\in [t_0,t_1-\delta
r^2]$ for which the hypotheses of the lemma hold.  The strategy of
the proof is to first restrict to the maximum subinterval of
$[t',t_2]$ of $[t',t'+\delta r^2]$ on which $k$ is bounded by
$\sqrt{2/(t-t')}$. If $t_2<t'+\delta r^2$, then $k$ achieves the
bound $\sqrt{2/(t-t')}$ at time $t_2$. We show that in fact on this
subinterval $k$ never achieves the bound. The result then follows.
To show that $k$ never achieves the bound, we show that on a
possibly smaller interval of time $[t',t_3]$ with $t_3\le t_2$ we
can write the restriction of the curve-shrinking flow to any
interval whose length at time $t'$ is $(0.9)r$ as a family of graphs
in a local coordinate system so that the function $f$ (of arc and
time) defining the graph has derivative along the arc bounded in
norm by $1/2$. We take $t_3\le t_2$ maximal with respect to these
conditions. Then with both the bound on $k$ and the bound on the
derivative of $f$ one shows that the spatial derivative of $f$ never
reaches $1/2$ and also that the curves do not move too much so that
they always remain in the coordinate patch. The only way that this
can happen is that if $t_2=t_3$, that is to say, on the entire time
interval where we have the curvature bound, we also can write the
curve-shrinking flow as a flow of graphs with small spatial
derivatives. Then it is convenient to replace the curve-shrinking
flow by an equivalent flow, introduced in \cite{AG}, called the
graph flow. Applying a simple maximum principle argument to this
flow we see that $k$ never achieves the value $\sqrt{2/(t-t')}$ on
the time interval $[t',t_2]$ and hence the curvature estimate
$k<\sqrt{2/(t-t')}$ holds throughout the interval $(t',t'+\delta
r^2]$.




\subsection{A bound for $\int kds$}

Recall that $k$ is the norm of the curvature vector $\nabla_SS$, and
in particular, $k\ge 0$. For any sub-arc $\gamma_{t'}$ of
$c(\cdot,t')$  at time $t'$ we let $\gamma_t$ be the result at time
$t$ of applying the curve-shrinking flow to $\gamma_{t'}$. The
purpose of this subsection is to show that $\int_{\gamma_t} kds$ is
small for all $t\in [t',t'+\delta r^2]$ and all initial arcs
$\gamma_{t'}$ of length at most $r$.



\begin{claim}[Uniform bound on length and total curvature on short sub-arcs]
\label{claim:curve-shrinking-length-curvature-uniform-bound}
\label{C0claim}
There is a constant $D_0<\infty$, depending only on $\Theta$ and the
curvature bound of the ambient Ricci flow such that for every $t\in
[t',t'+\delta r^2]$ and every sub-arc $\gamma_{t'}$ whose length is
at most $r$, we have $\int_{\gamma_t} kds<D_0$ and $l(\gamma_t)\le
D_0r$, where $l(\gamma_t)$ is the length of $\gamma_t$.
\end{claim}



\begin{proof}
\uses{cor:curve-shrinking-length-curvature-exponential-bound}
This is immediate from Corollary~\ref{cor:curve-shrinking-length-curvature-exponential-bound} applied to all
of $c(\cdot,t)$.
\end{proof}

Now we fix $t_2\le t'+\delta r^2$ maximal subject to the condition
that $k(x,t)\le \sqrt{\frac{2}{t-t'}}$ for all $x$ and all $t\in
[t',t_2]$. If $t_2<t'+\delta r^2$ then there is $x$ with
$k(x,t_2)=\sqrt{\frac{2}{(t_2-t')}}$.

Now consider a curve $\gamma_{t'}$ of length $r$. From the integral estimate in
the previous claim and the assumed pointwise estimate on $k$, we see that
$$\int_{\gamma_t}k^2 ds\le \max_{x\in\gamma_t}k(x,t)\int_{\gamma_{t}} kds
<\sqrt{\frac{2}{t-t'}}\cdot D_0.$$ Using Equation~(\ref{dLdteqn}),
it follows easily that, provided that $\delta>0$ is sufficiently
small, the length of $\gamma_t$ is at least $(0.9)r$ for all $t\in
[t',t_2]$, and more generally for any subinterval $\gamma'_{t'}$ of
$\gamma_{t'}$ and for any $t\in [t',t_2]$ the length of the
corresponding interval $\gamma'_t$ is at least $(0.9)$ times the
length of $\gamma'_{t'}$. We introduce a cut-off function on
$\gamma_{t'}\times [t',t_2]$ as follows. First, fix a smooth
function $\psi\colon [-1/2,1/2]\to [0,1]$ which is identically zero
on $[-0.50,-0.45]$ and on $[0.45,0.50]$, and is identically $1$ on
$[-3/8,3/8]$. There is a constant $D'$ such that $|\psi'|\le D'$ and
$|\psi''|\le D'$. Now we fix the midpoint $x_0\in \gamma_{t'}$ and
define the signed distance from $(x_0,t)$, denoted
$$s\colon \gamma_{t'}\times [t',t_2]\to \Ar,$$
as follows:
$$s(x,t)=\int _{x_0}^x|X(y,t)|dy.$$

We define the cut-off function
$$\varphi(x,t)=\psi\left(\frac{s(x,t)}{r}\right).$$



\begin{claim}[Bound on the time derivative of the cutoff function]
\label{claim:cutoff-function-time-derivative-bound}
\label{1887}
There is a constant $D_1$ depending only on the curvature bound for
 the ambient Ricci flow such that for any sub-arc $\gamma_{t'}$ of
 length $r$, defining $\varphi(x,t)$ as above, for all $x\in \gamma_t$
 and all $t\in [t',t_2]$ we have
$$\left|\frac{\partial \varphi(x,t)}{\partial t}\right|\le
\frac{D_1}{r\sqrt{t-t'}}+D_1.$$
\end{claim}


\begin{proof}
\uses{lem:curve-shrinking-first-variation,claim:curve-shrinking-length-curvature-uniform-bound}
Clearly,
$$\frac{\partial \varphi(x,t)}{\partial
t}=\psi'\left(\frac{s(x,t)}{r}\right)\cdot \frac{1}{r}\frac{\partial
s(x,t)}{\partial t}.$$ We know that $|\psi'|\le D'$ so that
$$\left|\frac{\partial \varphi(x,t)}{\partial
t}\right|\le \frac{D'}{r}\left|\frac{\partial s(x,t)}{\partial
t}\right|.$$ On the other hand,
$$s(x,t)=\int_{x_0}^x|X(y,t)|dy,$$
so that  $$ \left|\frac{\partial s(x,t)}{\partial t}\right| =
\left|\int_{x_0}^x\frac{\partial |X(y,t)|}{\partial t}dy\right|,$$
By Lemma~\ref{lem:curve-shrinking-first-variation} we have $$\frac{\partial
|X(y,t)|}{\partial y}dy  = \left(-\Ric(S(y,t),S(y,t))-k^2(y,t)\right)ds,$$
 so that there is a constant $D$ depending only on the bound of the
sectional curvatures of the ambient Ricci flow with
\begin{eqnarray*}\left|\frac{\partial s(x,t)}{\partial t}\right| &
\le &  \int_{x_0}^x(D+k^2)ds\le
Dl(\gamma_t)+\int_{x_0}^xk^2(y,t)ds(y,t), \end{eqnarray*} and hence
by Claim~\ref{claim:curve-shrinking-length-curvature-uniform-bound}
$$\left|\frac{\partial s(x,t)}{\partial t}\right|\le
  DD_0r+
\int_{x_0}^xk^2(y,t)ds(y,t).$$
  Using
the fact that $k^2\le 2/(t-t')$,  we have
$$\int_{x_0}^xk^2(y,t)ds(y,t)\le\sqrt{\frac{2}{t-t'}}\int_{x_0}^xkds\le
\frac{\sqrt{2}D_0}{\sqrt{t-t'}}.$$ Putting all this together, we see
that there is a constant $D_1$ such that
$$\left|\frac{\partial \varphi(x,t)}{\partial t}\right|\le
D_1(\frac{1}{r\sqrt{t-t'}}+1).$$
\end{proof}



\begin{claim}[Bound on the cutoff-weighted curvature integral]
\label{claim:cutoff-weighted-curvature-integral-bound}
\label{claim1885}
There is a constant $D_2$ depending only on the curvature bound of
the ambient Ricci flow and $\Theta$ and a constant $D_3$ depending
only on the curvature bound of the ambient Ricci flow, such that for
any $t\in [t',t_2]$ and any sub-arc $\gamma_{t'}$ of length $r$, we
have
$$\left|\frac{d}{dt}\int_{\gamma_t}\varphi
kds\right|\le
D_2\left(1+\frac{1}{r\sqrt{t-t'}}\right)+\frac{D_2}{r^2}+D_3\int_{\gamma_t}\varphi
kds.$$
\end{claim}


\begin{proof}
\uses{claim:cutoff-function-time-derivative-bound,claim:curve-shrinking-curvature-inequality,lem:curve-shrinking-length-curvature-growth,claim:curve-shrinking-length-curvature-uniform-bound}
We have
$$\left|\frac{d}{dt}\int_{\gamma_t}\varphi
kds\right|\le\left|\int_{\gamma_t}\frac{\partial
\varphi(x,t)}{\partial
t}kds\right|+\left|\int_{\gamma_t}\varphi\frac{\partial
(kds)}{\partial t}\right|.$$ Using Claim~\ref{claim:cutoff-function-time-derivative-bound} for the first
term and Claim~\ref{claim:curve-shrinking-curvature-inequality} and arguing as in the proof of
Lemma~\ref{lem:curve-shrinking-length-curvature-growth} for the second term, we have
$$\left|\frac{d}{dt}\int_{\gamma_t}\varphi
kds\right|\le
D_1\left(1+\frac{1}{r\sqrt{t-t'}}\right)\int_{\gamma_t}kds+\left|\int_{\gamma_t}
\varphi k''ds\right|+\int_{\gamma_t}C'_1\varphi kds,$$ where $C_1'$
depends only on the ambient curvature bound. We bound the first term
by
$$D_1D_0\left(\frac{1}{r\sqrt{t-t'}}+1\right),$$
where $D_0$ is the constant depending on $\Theta$ and the ambient
curvature bound from Claim~\ref{claim:curve-shrinking-length-curvature-uniform-bound}. Since the ends of
$\gamma_t$ are at distance at least $(0.45)r$ from $x_0$ all $t\in
[t',t_2]$, we see that for all $t\in [t',t_2]$
$$\int_{\gamma_t}\varphi k''=\int_{c(\cdot,t)}\varphi k''.$$
Integrating by parts we have
$$\int_{c(\cdot,t)} \varphi k'' ds=\int_{c(\cdot,t)}\varphi'' k ds,$$
where the prime here refers to the derivative along $c(\cdot,t)$
with respect to arc length. Of course $|\varphi''|\le
\frac{D'}{r^2}$. Thus, we see that
$$\left|\int_{\gamma_t}\varphi k''ds\right|\le
\frac{D'}{r^2}\int_{c(\cdot,t)}kds\le \frac{D'D_0}{r^2}.$$ Putting
all this together, we have
$$\left|\frac{d}{dt}\int_{\gamma_t}\varphi
kds\right|\le
D_2\left(1+\frac{1}{r\sqrt{t-t'}}\right)+\frac{D_2}{r^2}+D_3\int_{\gamma_t}\varphi
kds$$ for $D_2=D_0\,\max(D',D_1)$ and $D_3=C_1'$. This gives
the required estimate.
\end{proof}

\begin{corollary}[Square-root bound on the cutoff-weighted curvature integral]
\label{cor:cutoff-weighted-curvature-integral-sqrt-bound}
For any $t\in [t',t_2]$  and any sub-arc $\gamma_{t'}$ of length $r$
we have
$$\int_{\gamma_t}\varphi kds
\le  D_4\sqrt{\delta}$$ for a constant $D_4$ that depends only on
the sectional curvature bound of the ambient Ricci flow and
$\Theta$.
\end{corollary}

\begin{proof}
This is immediate from the previous result by integrating from $t'$
to $t_2\le t'+\delta r^2$, and using the fact that $\delta<1$ and
$r<1$, and using the fact that $$\int_{\gamma_{t'}}\varphi kds\le
\int_{\gamma_{t'}}kds<\delta$$ since $\gamma_{t'}$ has length at
most $r$.
\end{proof}

This gives:

\begin{corollary}[Uniform bound on the curvature integral over short sub-arcs]
\label{cor:curvature-integral-uniform-bound}
\label{D_4cor}
For $\gamma_{t'}\subset c(\cdot,t')$ a sub-arc of length at most $r$
 and for any $t\in [t',t_2]$, we have
$$\int_{\gamma_t}kds\le
2D_4\sqrt{\delta}.$$ For any $t\in [t',t_2]$ and any sub-arc
$J\subset c(\cdot,t)$ of length at most $r/2$ with respect to the
metric $h(t)$, we have
$$\int_Jk(x,t)ds(x,t)\le 2D_4\sqrt{\delta}.$$
\end{corollary}

\begin{proof}
\uses{cor:curve-shrinking-length-curvature-exponential-bound}
We divide an  interval $\gamma_{t'}\subset c(\cdot,t')$ of length at
most $r$ into two subintervals $\gamma'_{t'}$ and $\gamma''_{t'}$ of
lengths at most $r/2$. Let $\hat \gamma'_{t'}$ and $\hat
\gamma''_{t'}$ be intervals of length $r$ containing $\gamma'_{t'}$
and $\gamma''_{t'}$ respectively as  middle subintervals. We then
apply the previous corollary to  $\hat \gamma'_{t'}$ and $\hat
\gamma''_{t'}$ using the fact that $\varphi k\ge 0$ everywhere and
$\varphi k=k$ on the middle subintervals of $\hat \gamma'_{t'}$ and
$\hat \gamma ''_{t'}$. For an interval $J\subset \gamma_t$ of length
$r/2$, according to Lemma~\ref{cor:curve-shrinking-length-curvature-exponential-bound} the length of
$\gamma_{t'}|_J$ with respect to the metric $h(t')$ is at most $r$,
and hence this case follows from the previous case.
\end{proof}


\subsection{Writing the curve flow as a graph}

Now we restrict attention to $[t',t_2]$, the maximal interval in
$[t',t'+\delta r^2]$ where $k^2\le 2/(t-t')$. Let $\gamma_{t'}$ be
an arc of length $r$ in $c(\cdot,t')$ and let $x_0$ be the central
point of $\gamma_{t'}$. Denote $\gamma_{t'}(x_0)=p\in W$. We take
the $h(t')$-exponential mapping from $T_pW\to W$. This map will be a
local diffeomorphism out to a distance determined by the curvature
of $h(t')$.   For an appropriate choice of the ball (depending on
the ambient curvature bound) the metric on the ball induced by
pulling back $h(t)$ for all $t\in [t',t_2]$ will be within $\delta$
in the $C^1$-topology to the Euclidean metric $h'=h(t')_p$. By this
we mean that
\begin{enumerate}
\item[(1)] $\left|\langle X,Y\rangle_{h(t)}-\langle
X,Y\rangle_{h'}\right|<\delta |X|_{h'}|Y|_{h'}$ for all tangent vectors in the
coordinate system, and \item[(2)] viewing the connection $\Gamma$ as a bilinear
map on the coordinate space with values in the coordinate space we have $
|\Gamma(X,Y)|_{h'}<\delta |X|_{h'}|Y|_{h'}$.
\end{enumerate}

We choose $0<r_0\le 1$ so that it is much smaller than this
distance, and hence  $r$ is also much smaller than this distance. We
lift to the ball in $T_pW$.

 We fix orthonormal coordinates with respect to the metric $h'$ so
that the tangent vector of $\gamma_{t'}(x_0)$ points in the positive
$x^1$-direction. Using these coordinates we decompose the coordinate
patch as a product of an interval in the $x^1$-direction and an open
ball, $B$, spanned by the remaining Euclidean coordinates. From now
on we shall work in this coordinate system using this product
structure.  To simplify the notation in the coming computations, we
rename the $x^1$-coordinate the $z$-coordinate. Ordinary derivatives
of a function $\alpha$ with respect to $z$ are written $\alpha_z$.
When we write norms and inner products without indicating the metric
we implicitly mean that the metric is $h(t)$. When we use the
Euclidean metric on these coordinates we denote it explicitly. Next,
we wish to understand how $\gamma_t$ moves in the Euclidean
coordinates under the curve-shrinking flow. Since we have
$|\nabla_SS|_{h}=k$, it follows that $|\nabla_SS|_{h'}\le
\sqrt{1+\delta}k\le 2/\sqrt{t-t'}$, and hence, integrating tells us
that for any $x\in \gamma_{t'}$ we have
$$\left|\gamma_t(x)-\gamma_{t'}(x)\right|_{h'}\le 4\sqrt{t-t'}\le 4\sqrt{\delta}r.$$
This shows that for every $t\in [t',t_2]$, the curve $\gamma_t$ is
contained in the coordinate patch that we are considering. This
computation also implies that the $z$-coordinate of $\gamma_t$
changes by at most $4\sqrt{\delta} r$ over this time interval.



Because the total curvature of $\gamma_{t'}$ is small and the metric
is close to the Euclidean metric, it follows that the tangent vector
at every point of $\gamma_{t'}$ is close to the positive
$z$-direction. This means that we can write $\gamma_{t'}$ as a graph
of a function $f$ from a subinterval in the $z$-line to $Y$ with
$|f_z|_{h'}<2\delta$. By continuity, there is $t_3\in (t',t_2]$ such
that all the curves $\gamma_{t}$ are written as graphs of functions
(over subintervals of the $z$-axis that depend on $t$) with
$|f_z|_{h'}\le 1/10$. That is to say, we have an open subset $U$ of
the product of the $z$-axis with $[t',t_3]$, and the evolving curves
define a map $\tilde\gamma$ from $U$ into the coordinate system,
where the slices at constant time are graphs $z\mapsto (z,f(z,t))$
and are the curves $\gamma_t$. Using the coordinates $(z,t)$ gives a
new flow of curves by moving in the $t$-direction. This new flow is
called the {\sl graph-flow}. It is
 a reparameterization of the curve shrinking flow in such a way that
the $z$-coordinate is preserved. We denote by $Z=Z(z,t)$ the image
under the differential of the map $\widetilde\gamma$ of the tangent
vector in the $z$-direction and by $Y(z,t)$ the image under the
differential of $\widetilde \gamma$ of the tangent vector in the
$t$-direction. Notice that $Z$ is the tangent vector along the
parameterized curves in the graph flow. Since we are now using a
different parameterization of the curves from the one determined by
the curve-shrinking flow, the tangent vector $Z$ has the same
direction but not necessarily the same length as the tangent vector
$X$ from the curve-shrinking parameterization. Also, notice that in
the Euclidean norm we have $|Z|_{h'}^2=1+|f_z|^2_{h'}$. It follows
that on $U$ we have
$$(1-\delta)(1+|f_z|_{h'}^2)\le |Z(z,t)|^2_{h(t)}\le (1+\delta)(1+|f_z|_{h'}^2).$$
In particular, because of our restriction to the subset where $|f_z|_{h'}\le
1/10$ we have  $(1-\delta)\le |Z(z,t)|^2_{h(t)}\le (1.01)(1+\delta)$.




Now we know that $\gamma_{t'}$ is a graph of a function $f(z,t')$
defined on some  interval $I$ along the $z$-axis.  Let $I'$ be the
subinterval of $I$ centered in $I$ with $h'$-length $(0.9)$ times
the $h'$-length of $I$. By the above estimate on $|Z|$ it follows
that  the restriction of $\gamma_{t'}$ to $I'$ has length between
$(0.8)r$ and $r$, and also that the $h'$-length of $I'$ is between
$(0.8)r$ and $r$. The above estimate means that, provided that
$\delta>0$ is sufficiently small, for every $t\in [t',t_3]$ there is
a subinterval of $\gamma_t$ that is the graph of a function defined
on all of $I'$. We now restrict attention to the family of curves
parameterized by $I'\times [t',t_3]$. For every $t\in[t',t_3]$ the
curve $\gamma_t|_{I'}$ has length between $(0.8)r$ and $r$. The
curve-shrinking flow is not defined on this product because under
the curve-shrinking flow the $z$-coordinate of any given point is
not constant. But the graph flow defined above, and studied in
\cite{AG} (in the case of Euclidean background metric), is defined
on $I'\times [t',t_3]$ since this flow preserves the $z$-coordinate.
The time partial derivative in the curve-shrinking flow is given by
\begin{equation}\label{1814}
\nabla_SS=\frac{\nabla_ZZ}{|Z|^2}-\frac{1}{|Z|^4}\langle \nabla_ZZ,Z\rangle Z.
\end{equation} The time partial derivative in the graph-flow is given by
$Y=\partial\widetilde \gamma/\partial t$. The tangent vector $Y$ is
characterized by being $h'$-orthogonal to the $z$-axis and differing from
$\nabla_SS$ by a functional multiple of $Z$.

\begin{claim}[Comparison of the graph-flow and curve-shrinking velocities]
\label{claim:graph-flow-velocity-comparison}
\label{curvflowcompare}
$$Y=\frac{\nabla_ZZ-\langle\Gamma(Z,Z),\partial_z\rangle_{h'} Z}{|Z|^2}
=\nabla_SS+\left(\frac{\langle \nabla_ZZ,Z\rangle}{|Z|^4}
-\frac{\langle\Gamma(Z,Z),\partial_z\rangle_{h'} }{|Z|^2}\right)Z.$$
\end{claim}

\begin{proof}
In our Euclidean coordinates, $Z=(1,f_z)$ so that
$\nabla_ZZ=(0,f_{zz})+\Gamma(Z,Z)$. Thus,
$$\langle \nabla_ZZ,\partial_z\rangle_{h'}=\langle
\Gamma(Z,Z),\partial_z\rangle_{h'}.$$ Since $\langle
Z,\partial_z\rangle_{h'}=1$, it follows that
$$\frac{ \nabla_ZZ-\langle\Gamma(Z,Z),\partial_z\rangle_{h'} Z}{|Z|^2}$$ is
$h'$-orthogonal to the $z$-axis and hence is a multiple of $Y$. Since it
differs by a multiple of $Z$ from $\nabla_SS$, it follows that it is $Y$. This
gives the first equation; the second follows from this and
Equation~(\ref{1814}).
\end{proof}

To simplify the notation we set
$$\psi(Z)=\frac{\langle\Gamma(Z,Z),\partial_z\rangle_{h'}}{|Z|^2}.$$
Notice that from the conditions on $\Gamma$ and $h'$ it follows
immediately that $|\psi(Z)|<(1.5)\delta$.

\subsection{Proof that $t_3=t_2$}

At this point we have a product coordinate system on which the
metric is almost the Euclidean metric in the $C^1$-sense, and we
have the graph flow given by
$$Y=\frac{\partial\widetilde \gamma}{\partial
t}=\frac{\nabla_ZZ}{|Z|^2}-\psi(Z)Z$$ defined on $[t',t_3]$ with
image always contained in the given coordinate patch and written as
a graph over a fixed interval $I'$ in the $z$-axis. For every $t\in
[t',t_3]$ the length of $\gamma_{t'}|_{I'}$ in the metric $h(t')$ is
between $(0.8)r$ and $r$. The function $f(z,t)$ whose graphs give
the flow satisfies $|f_z|_{h'}\le 1/10$. Our next goal is to
estimate $|f_z|_{h'}$ and show that it is always less than  $1/10$
as long as $k^2\le 2/(t-t')$ and $t-t'\le \delta r^2$ for a
sufficiently small $\delta$, i.e., for all $t\in [t',t_2]$; that is
to say, our next goal is to prove that $t_3=t_2$. In all the
arguments that follow $C'$ is a constant that depends only on the
curvature bound for the ambient Ricci flow, but the value of $C'$ is
allowed to change from line to line.


The first step in doing this is to consider the angle between
$\nabla_ZZ$ and $Z$.

\begin{claim}[Angle and norm bounds for the graph flow]
\label{claim:graph-flow-angle-and-norm-bounds}
\label{1891}
Provided that $\delta>0$ is sufficiently small,
the angle (measured in $h(t)$) between $Y$ and $Z=(1,f_z)$ is
greater than $\pi/4$. Also,
\begin{enumerate}
\item[(1)] $$k\le |Y|<\sqrt{2}k.$$
\item[(2)] $$\left|\langle \nabla_ZZ,Z\rangle\right|<(k+ 2\delta) |Z|^3.$$
\item[(3)] $$\left|\nabla_ZZ\right|<2(|Y|+\delta).$$
\item[(4)] $$|\langle Y,Z\rangle |\le |Y||f_z|(1+3\delta).$$
\end{enumerate}
\end{claim}

\begin{proof}
Under the hypothesis that $|f_z|_{h'}\le 1/10$, it is easy to see
that the Euclidean angle between $(0,f_{zz})$ and $(1,f_z)$ is at
most $\pi/2-\pi/5$. From this, the first statement follows
immediately provided that $\delta$ is sufficiently small. Since $Y$
is the sum of $\nabla_SS$ and a multiple of $Z$ and since
$\nabla_SS$ is $h(t)$-orthogonal to $Z$, it follows that
$|Y|=|\nabla_SS|\left(\cos(\theta)\right)^{-1}$, where $\theta$
is the angle between $\nabla_SS$ and $Y$. Since $Y$ is a multiple of
$(0,f_{zz})$, it follows from the first part of the claim that the
$h(t)$-angle between $Y$ and $\nabla_SS$ is less than $\pi/4$. Item
(1) of the claim then follows from the fact that by definition
$|\nabla_SS|=k$.

Since
$$\frac{\nabla_ZZ}{|Z|^2}=Y+\psi(Z)Z,$$ and $|Z|^2\le (1.01)(1+\delta)$, the third item is
immediate. For Item (4), since $Y$ is $h'$-orthogonal to the $z$-axis, we have
$$\left|\langle Y,Z\rangle_{h'}\right|=\left|\langle Y,(0,f_z)\rangle_{h'}\right|\le |Y|_{h'}|f_z|_{h'}.$$
From this and the comparison of $h(t)$ and $h'$, the Item (4) is
immediate. Lastly, let us consider Item (2). We have
$$\langle Y,Z\rangle =\frac{\langle \nabla_ZZ,Z\rangle}{|Z|^2}-\langle
\Gamma(Z,Z),\partial_z\rangle_{h'}.$$ Thus, from Item (4) we have
$$\frac{\langle \nabla_ZZ,Z\rangle}{|Z|^2}\le |Y||f_z|(1+3\delta)+(1.5)\delta.$$
Since $Y<\sqrt{2}k$ and $|f_z|<1/10$, the Item (2) follows.
\end{proof}

\begin{claim}[Bounds for the connection correction term]
\label{claim:graph-flow-connection-term-bounds}
\label{psiclaim}
The following hold provided that $\delta>0$ is sufficiently small:
\begin{enumerate}
\item[(1)] $|Z(\psi(Z))|<C'(1+\delta |Y|)$, and
\item[(2)] $|Y(\psi(Z))|<C'(|Y|+\delta|\nabla_ZY|)$.
\end{enumerate}
(Recall that $C'$ is a constant depending only on the curvature bound of the
ambient Ricci flow.)
\end{claim}


\begin{proof}
For the first item, we write $Z(\psi(Z))$ as a sum of terms where the
differentiation by $Z$ acts on the various. When the $Z$-derivative acts on
$\Gamma$ the resulting term has norm bounded by a constant depending only on
the curvature of the ambient Ricci flow. When the $Z$-derivative acts on one of
the $Z$-terms in $\Gamma(Z,Z)$ the norm of the result is bounded by
$2\delta|\nabla_ZZ||Z|$. Action on each of the other $Z$-terms gives a term
bounded in norm by the same expression. Lastly, when the $Z$-derivative acts on
the constant metric $h'$ the norm of the result is bounded by $2\delta^2$.
Since $|\nabla_ZZ|\le 2(|Y|+\delta)$, the first item follows.

We compute $Y(\psi(Z))$ in a similar fashion. When the $Y$-derivative acts on
the $\Gamma$, the norm of the result is bounded by $C'|Y|$. When the
$Y$-derivative acts on one of the $Z$-terms the norm of the result is bounded
by $2\delta|\nabla_YZ|$. Lastly, when the $Y$-derivative acts on the constant
metric $h'$, the norm of the result is bounded by $\delta^2|Y|$. Putting all
these terms together establishes the second inequality above.
\end{proof}







Now we wish to compute $\int_{I'\times\{t\}}|Z|^2dz.$ To do this we
first note that using the definition of $Y$, and arguing as in the
proof of the first equation in of Lemma~\ref{lem:curve-shrinking-first-variation} we have we
have
\begin{eqnarray*}\frac{\partial}{\partial t}|Z|^2
& = & -2\Ric(Z,Z)+2\langle \nabla_YZ,Z\rangle \\
& = & -2\Ric(Z,Z)+2\langle \nabla_ZY,Z\rangle \end{eqnarray*}
Direct computation shows that
$$2\langle
\nabla_Z\left(\frac{\nabla_ZZ}{|Z|^2}\right),Z\rangle
=Z\left(\frac{Z(|Z|^2)}{|Z|^2}\right)-2\frac{|\nabla_ZZ|^2}{|Z|^4}|Z|^2.$$ Thus
from the Claim~\ref{claim:graph-flow-velocity-comparison}, we have
\begin{eqnarray}\label{Z_t}
\frac{\partial}{\partial t}|Z|^2 & = & 2\langle\nabla_ZY,Z\rangle -2\Ric(Z,Z)\nonumber \\
& = &
Z\left(\frac{Z(|Z|^2)}{|Z|^2}\right)-2\frac{|\nabla_ZZ|^2}{|Z|^4}|Z|^2-2\langle
\nabla_Z(\psi(Z)Z),Z\rangle-2\Ric(Z,Z) \nonumber \\
& = & Z\left(\frac{Z(|Z|^2)}{|Z|^2}\right)-2|Y|^2|Z|^2 +V,
\end{eqnarray}
where $$V=-4|Z|^2\langle Y,\psi(Z)Z\rangle-2\psi^2(Z)|Z|^4-2\langle
\nabla_Z(\psi(Z)Z),Z\rangle-2\Ric(Z,Z).$$ By Item (1) in
Claim~\ref{claim:graph-flow-connection-term-bounds} and Item (4) in Claim~\ref{claim:graph-flow-angle-and-norm-bounds} we have
\begin{equation}\label{Vest}
|V|<C'(1+\delta|Y|).
\end{equation}





Using this and the fact that $|Y|\le \sqrt{2}k$ we compute:
\begin{eqnarray*}
\frac{d}{dt}\int_{I'\times\{t\}}|Z|^2dz & \le & \int_{I'\times\{t\}}
Z\left(\frac{Z(|Z|^2)}{|Z|^2}\right)dz +\int_{I'\times \{t\}}\left(C'(1+\delta
k)\right) dz
 \\
& = & \frac{Z(|Z|^2)}{|Z|^2}\Bigl|_0^a\Bigr.+\int_{I'\times\{t\}}\left(C'(1+\delta k)\right) dz \\
& = & 2\frac{\langle \nabla_ZZ,Z\rangle}{|Z|^2}\Bigl|_0^a\Bigr.
+\int_{I'\times\{t\}}\left(C'(1+\delta k)\right) dz,
\end{eqnarray*}
where we denote the endpoints of $I'$ by $\{0\}$ and $\{a\}$. By Item (2) in
Claim~\ref{claim:graph-flow-angle-and-norm-bounds}, the first term is at most
$2(k+2\delta)\sqrt{(1.01)(1+\delta)}$, which is at most $\frac{8}{\sqrt{t-t'}}$
and the second term is at most $C'(1+\delta k) r$. Now integrating from $t'$ to
$t$ we see that for any $t\in [t',t_3]$ we have
$$\int_{I'\times\{t\}}|Z|^2dz\le\int_{I'\times\{t'\}}|Z|^2dz+16\sqrt{\delta}r+C'\delta
r^3+C'\delta^{3/2}r^2.$$  Since $|f_z(z,t')|_{h'}\le 2\delta$ and $|Z|^2$ is
between $(1-\delta)(1+|f_z|_{h'}^2)$ and $(1+\delta)(1+|f_z|_{h'}^2)$,  we see
that $\int_{I'\times\{t'\}}|Z|^2dz\le (1+3\delta)\ell_{h'}(I')$. It follows
that for any $t\in [t',t_3]$ we have
$$\int_{I'\times\{t\}}|Z|^2dz\le (1+3\delta)l_{h'}(I')+C'(\sqrt{\delta}r+\delta r^3+\delta^{3/2}r^2).$$
Since $|Z|^2$ is between $(1-\delta)(1+|f_z|_{h'}^2)$ and
$(1+\delta)(1+|f_z|_{h'}^2)$, we see that there is a constant
$C_1''$ depending only on the ambient curvature bound such that for
any $t\in [t',t_3]$, denoting by $\ell_{h'}(I')$ the length of $I'$
with respect to $h'$, we have
$$\int_{I'\times\{t\}}|f_z|_{h'}^2dz\le 4\delta\ell_{h'}(I')
+C_1''(\sqrt{\delta}r+\delta r^3+\delta^{3/2}r^2).$$ Since
$(0.8)r\le \ell_{h'}(I')\le r<1$, we see  that provided that
$\delta$ is sufficiently small, for each $t\in [t',t_3]$ there is
$z(t)\in I'$ such $|f_z(z(t),t)|^2_{h'}\le 2C_1''\sqrt{\delta}$. If
we have chosen $\delta$ sufficiently small, this means that for each
$t\in [t',t_3]$ there is $z(t)$ such that $|f_z(z(t),t)|_{h'}\le
1/20$. Since by Corollary~\ref{cor:curvature-integral-uniform-bound}
$\int_{I'\times\{t\}}kds<2D_4\sqrt{\delta}$,
 provided that $\delta$ is sufficiently small, it follows that for all
 $t\in[t',t_3]$ the
curve $\gamma_t|_{I'}$ is a graph of $(z,t)$ and $|f_z|_{h'}<1/10$. But by
construction either $t_3=t_2$ or there is a point in $(z,t_3)\in I'\times
\{t_3\}$ with $|f_z(z,t_3)|_{h'}=1/10$. Hence, it must be the case that
$t_3=t_2$, and thus our graph curve flow is defined for all $t\in [t',t_2]$ and
satisfies the derivative bound $|f_z|_{h'}< 1/10$ throughout the interval
$[t',t_2]$.

\subsection{Proof that $t_2=t'+\delta r^2$}

The last step is to show that the inequality $k^2< 2/(t-t')$ holds
for all $t\in [t',t'+\delta r^2]$.

 We fix a point $x_0$. We continue all the notation, assumptions and results of the previous section.
 That is to say, we lift the evolving family of curves to the tangent space $T_{x_0}M$
 using the exponential mapping, which is a
 local diffeomorphism. This tangent space is split as the
product of the $z$-axis and $B$. On this coordinate system we have
the evolving family of Riemannian metrics $h(t)$ pulled back from
the Ricci flow and also we have the Euclidean metric $h'$ from the
metric $h(t')$ on $T_{x_0}M$. We fix an interval $I'$ on the
$z$-axis
  of $h'$-length between $(0.8)r$ and $r$.    We choose $I'$ to be centered at
 $x_0$ with respect to the $z$-coordinate. On $I'\times [t',t_2]$
  we have the graph-flow  which is reparameterization
  of the pull back of the curve-shrinking flow. The graph-flow is given as the graph
  of a function $f$ with $|f_z|_{h'}<1/10$. The vector fields $Z$
  and $Y$ are as in the last section.

 We follow closely the discussion in Section 4 of \cite{AG} (pages 293 -294).
 Since we are not working in a flat background, there are two
 differences:
 (i) we take covariant derivatives instead of ordinary derivatives
 and (ii) there are various correction terms from curvature, from covariant derivatives,
 and from the fact that $Y$ is not equal to $\nabla_ZZ/|Z|^2$.

Notice that
\begin{eqnarray*}
Z\left(\frac{Z(|Z|^2)}{|Z|^2}\right) & = & \frac{|Z|^2_{zz}}{|Z|^2}-
\frac{\left(|Z|^2_z\right)^2}{|Z|^4} \\
& = & \frac{|Z|^2_{zz}}{|Z|^2}-4\langle Y,Z\rangle^2 -8\langle
Y,Z\rangle \psi(Z)|Z|^2-4\psi^2(Z)|Z|^4 \end{eqnarray*}
 Thus, it
follows from Equation~(\ref{Z_t}) that we have
$$\frac{\partial}{\partial
 t}|Z|^2=\frac{(|Z|^2)_{zz}}{|Z|^2}-2|Z|^2|Y|^2-4\langle Z,Y\rangle^2+V,$$
where $|V|\le C'(1+\delta |Y|)$ for a constant $C'$ depending only
on the curvature bound of the ambient flow.

Similar computations show that
 \begin{eqnarray*}\frac{\partial}{\partial
 t}|Y|^2 & = &
\frac{(|Y|^2)_{zz}}{|Z|^2}-\frac{2|\nabla_ZY|^2}{|Z|^2}
 -\frac{4}{|Z|^2}\langle\frac{\nabla_ZZ}{|Z|^2},Y\rangle\langle
 \nabla_ZY,Z\rangle \\
 &  & -2\Ric(Y,Y)+2\frac{\Rm(Y,Z,Y,Z)}{|Z|^2}-2\langle \nabla_Y(\psi(Z)Z),Y\rangle.
 \end{eqnarray*}
Of course, $$\langle \frac{\nabla_ZZ}{|Z|^2},Y\rangle=|Y|^2+\psi(Z)\langle
Z,Y\rangle.$$ Hence, putting all this together and using Claim~\ref{claim:graph-flow-connection-term-bounds}
we have
$$\frac{\partial}{\partial
 t}|Y|^2  =
\frac{(|Y|^2)_{zz}}{|Z|^2}-\frac{2|\nabla_ZY|^2}{|Z|^2}
 -\frac{4|Y|^2}{|Z|^2}\langle
 \nabla_ZY,Z\rangle+W,$$
 where
 $$|W|\le C'|Y|(|Y|+\delta|\nabla_ZY|).$$


 Now let us consider
$$Q=\frac{|Y|^2}{2-|Z|^2}.$$
Notice that since $|f_z|_{h'}<1/10$, it follows that $1-\delta\le
|Z|^2<(1.01)(1+\delta)$ on all of $[t',t_2]$. We now make computation following
the computations on p. 294 of \cite{AG} and adding in the error terms.
\begin{eqnarray*}
Q_t & = & \frac{|Y|^2_t}{(2-|Z|^2)}+\frac{|Y|^2|Z|^2_t}{(2-|Z|^2)^2} \\
& = & \frac{|Y|^2_{zz}}{|Z|^2(2-|Z|^2)}-\frac{2|\nabla_ZY|^2}{|Z|^2(2-|Z|^2)}
 -\frac{4|Y|^2}{|Z|^2(2-|Z|^2)}\langle
 \nabla_ZY,Z\rangle+\frac{W}{(2-|Z|^2)} \\
 & &
 +\frac{|Y|^2|Z|^2_{zz}}{|Z|^2(2-|Z|^2)^2}-\frac{2|Z|^2||Y|^4}{(2-|Z|^2)^2}
 -\frac{4|Y|^2}{(2-|Z|^2)^2}\langle Z,Y\rangle^2+\frac{|Y|^2}{(2-|Z|^2)^2}V.
\end{eqnarray*}

On the other hand,
\begin{eqnarray*}
\frac{Q_{zz}}{|Z|^2}=\frac{|Y|^2_{zz}}{|Z|^2(2-|Z|^2)}+
\frac{|Y|^2|Z|^2_{zz}}{|Z|^2(2-|Z|^2)^2}+\frac{2|Y|^2_z|Z|^2_z}{|Z|^2|(2-|Z|^2)^2}
+\frac{2|Y|^2\left(|Z|^2_z\right)^2}{|Z|^2(2-|Z|^2)^3}.
\end{eqnarray*}


From Claim~\ref{claim:graph-flow-angle-and-norm-bounds} we have
$$|Z|^2_z=2\langle\nabla_ZZ,Z\rangle=2|Z|^2\langle Y,Z\rangle+2\psi(Z)|Z|^4.$$
Plugging in this expansion gives
\begin{eqnarray*}
\frac{Q_{zz}}{|Z|^2} & = & \frac{|Y|^2_{zz}}{|Z|^2(2-|Z|^2)}+
\frac{|Y|^2|Z|^2_{zz}}{|Z|^2(2-|Z|^2)^2} \\
& & +\frac{8\langle \nabla_ZY,Y\rangle\langle
Y,Z\rangle}{(2-|Z|^2)^2}+\frac{8\psi(z)|Z|^2\langle
\nabla_ZY,Y\rangle}{(2-|Z|^2)^2} \\
& & +\frac{8|Z|^2|Y|^2\langle
Y,Z\rangle^2}{(2-|Z|^2)^3}+\frac{16\psi(Z)|Z|^4|Y|^2\langle
Y,Z\rangle}{(2-|Z|^2)^3}+\frac{8\psi^2(Z)|Z|^6|Y|^2}{(2-|Z|^2)^3}.
\end{eqnarray*}
Expanding, we have
\begin{eqnarray*}
\frac{Q_{zz}}{|Z|^2} & = & \frac{|Y|^2_{zz}}{|Z|^2(2-|Z|^2)}+
\frac{|Y|^2|Z|^2_{zz}}{|Z|^2(2-|Z|^2)^2}+\frac{8\langle
\nabla_ZY,Y\rangle\langle Y,Z\rangle}{(2-|Z|^2)^2}  \\
& & +\frac{8|Y|^2|Z|^2\langle Y,Z\rangle^2}{(2-|Z|^2)^3}+U,
\end{eqnarray*}
where
$$|U|\le C'(|Y|^2+\delta|\nabla_ZY||Y|+\delta|Y|^3).$$

Comparing the formulas yields
\begin{eqnarray*}
Q_t & = & \frac{Q_{zz}}{|Z|^2}-\frac{8\langle \nabla_ZY,Y\rangle\langle
Y,Z\rangle}{(2-|Z|^2)^2}-\frac{8|Y|^2|Z|^2\langle Y,Z\rangle^2}{(2-|Z|^2)^3} \\
& & -\frac{2|\nabla_ZY|^2}{|Z|^2(2-|Z|^2)}
 -\frac{4|Y|^2}{|Z|^2(2-|Z|^2)}\langle
 \nabla_ZY,Z\rangle \\
 & & -\frac{2|Z|^2||Y|^4}{(2-|Z|^2)^2}
 -\frac{4|Y|^2}{(2-|Z|^2)^2}\langle Z,Y\rangle^2+A,
\end{eqnarray*}
where
$$|A|\le C'(|Y|^2+\delta|\nabla_ZY||Y|+\delta|Y|^3).$$

Using Item (4) of Claim~\ref{claim:graph-flow-angle-and-norm-bounds} this leads to
\begin{eqnarray*}
Q_t & \le & \frac{Q_{zz}}{|Z|^2}+\frac{8(1+3\delta)|Y|^2|f_z|
|\nabla_ZY|}{(2-|Z|^2)^2} -\frac{4|Y|^2\langle
 \nabla_ZY,Y\rangle }{|Z|^2(2-|Z|^2)} \\
 & &
-\frac{2|\nabla_ZY|^2}{|Z|^2(2-|Z|^2)}
 -\frac{2|Z|^2||Y|^4}{(2-|Z|^2)^2}
 +|A|
 \end{eqnarray*}
Next, we have
\begin{claim}[Mixed derivative bound for the graph flow]
\label{claim:graph-flow-mixed-derivative-bound}
$$|\langle \nabla_ZY,Z\rangle|\le (|f_z|(|\nabla_ZY|+2\delta|Y|)+\delta|Z|^2|Y|)(1+\delta).$$
\end{claim}

\begin{proof}
Since $Y=(0,\phi)$ for some function $\phi$, we have
$\nabla_ZY=(0,\phi_z)+\Gamma(Z,Y)$ and hence
$$|\langle\nabla_ZY,Z\rangle_{h'}| = |\langle \nabla_ZY,(1,f_z)\rangle_{h'}|\le |\langle
f_z ,\phi_z\rangle_{h'}|+|\langle \Gamma(Z,Y),Z\rangle_{h'}|\le
|f_z|_{h'}|\phi_z|_{h'}+\delta|Z|^2|Y|.$$ On the other hand
$\nabla_ZY=(0,\phi_z)+\Gamma(Z,Y)$ so that $|\phi_z|_{h'}\le|\nabla_ZY|+\delta
|Z||Y|$. From this the claim follows.
\end{proof}

Now for $\delta>0$ sufficiently small, using the fact that
$1-\delta<|Z|^2<(1+\delta)(1.01)$ we can rewrite this as
\begin{eqnarray*}
Q_t & \le & \frac{Q_{zz}}{|Z|^2}+\frac{8(1+3\delta)|Y|^2|f_z|
|\nabla_ZY|}{(2-|Z|^2)^2} +\frac{4|Y|^2|(1+\delta)
 |\nabla_ZY||f_z|}{|Z|^2(2-|Z|^2)} \\
 & & -\frac{2|\nabla_ZY|^2}{|Z|^2(2-|Z|^2)}-\frac{(1.95)|Y|^4}{(2-|Z|^2)^2} +\tilde
 A,
 \end{eqnarray*}
where $\tilde A\le C'(|Y|^2+\delta|Y||\nabla_ZY|+\delta|Y|^3)$. Of
course, $|Y||\nabla_ZY|+|Y|^3\le 2|Y|^2+|\nabla_ZY|^2+|Y|^4$. Using
this, provided that $\delta$ is sufficiently small, we can rewrite
this as

\begin{eqnarray*}
Q_t & \le & \frac{Q_{zz}}{|Z|^2} +\frac{1}{(2-|Z|^2)}\cdot
\left[\frac{8(1+3\delta)|\nabla_ZY||f_z||Y|^2-(0.9)|Y|^4}{(2-|Z|^2)} \right.\\&
&
 +\left. \frac{4|Y|^2|(1+\delta)
 |\nabla_ZY||f_z|-(1.9)|\nabla_ZY|^2}{|Z|^2}\right]-Q^2+\tilde A''
\end{eqnarray*}
where $\tilde A''\le C'(|Y|^2)$. We denote the quantity within the
brackets by $B$ and we estimate
\begin{eqnarray*}
B & \le &
8(1+3\delta)\frac{|Y|^2|\nabla_ZY|(1/10)(1+\delta)}{(2-(1.01)(1+\delta))}
+\frac{4(1/10)(1+\delta)|Y|^2|\nabla_ZY|}{(1-\delta)} \\
& & -\frac{(1.9)}{(1.01)(1+\delta)}|\nabla_ZY|^2
-\frac{(0.9)|Y|^4}{2-(1.01)(1+\delta)} \\
& \le  & (1.6)|Y|^2|\nabla_ZY|-(0.8)|\nabla_ZY|^2-(0.8)|Y|^4 \\
& \le  & 0.
\end{eqnarray*}
Therefore,
$$Q_t\le \frac{Q_{zz}}{|Z|^2}-Q^2+|\tilde A|\le \frac{Q_{zz}}{|Z|^2}-(Q-C_1')^2+(C_1')^2,$$
for some constant $C'_1>1$ depending only on the curvature bound for the
ambient Ricci flow.


Denote by $l$ the length of $I'$ under $h'$. As we have already
seen, $(0.8)r\le l\le r$. We translate the $z$-coordinate so that
$z=0$ is one endpoint of $I'$ and $z=l$ is the other endpoint; the
point $x_0$ then corresponds to $z=l/2$. Consider the function
$g=l^2/(z^2(l-z)^2)$ on $I'\times [t',t_2]$. Direct computation
shows that $g_{zz}\le 12g^2$.
 Now set
$$\widetilde Q=Q-C_1'$$
 and
$$h=\frac{1}{t-t'}+\frac{4(1-\delta)^{-1}l^2}{z^2(l-z)^2}+C_1'.$$
 Then
$$ -h_t+(1-\delta)^{-1}h_{zz}+(C_1')^2\le h^2,$$
so that
$$(\widetilde Q-h)_t\le \frac{\widetilde Q_{zz}}{|Z|^2}-\frac{h_{zz}}{1-\delta}-\widetilde
Q^2+h^2.$$ Since both $h$ and $h_{zz}$ are positive, at any point
where $\widetilde Q -h\ge 0$ and $\widetilde Q_{zz}<0$, we have
$(\widetilde Q-h)_t<0$. At any point where $Q_{zz}\ge 0$, using the
fact that $|Z|^2\ge (1-\delta)$ we have
$$(\widetilde Q-h)_t\le (1-\delta)^{-1}(\widetilde Q-h)_{zz}-\widetilde
Q^2+h^2.$$ Thus, for any fixed $t$, at any local maximum for
$(\widetilde Q-h)(\cdot,t)$ at which $(\widetilde Q-h)$ is $\ge 0$
we have $(\widetilde Q-h)_t\le 0$. Since $\widetilde Q-h$ equals
$-\infty$ at the end points of $I'$ for all times, there is a
continuous function $f(t)=\max_{z\in I'}(\widetilde Q-h)(z,t)$,
defined for all $t\in (t',t_2]$ approaching $-\infty$ uniformly as
$t$ approaches $t'$ from above. By the previous discussion, at any
point where $f(t)\ge 0$ we have $f'(t)\le 0$ in the sense of forward
difference quotients. It now follows that $f(t)\le 0$ for all $t\in
(t',t_2]$.
 This means that for all $t\in (t',t_2]$ at the $h'$-midpoint $x_0$ of
$I'$ (the point where $z=l/2$) we have
$$Q(x_0,t)\le \frac{1}{t-t'}+\frac{16\cdot 4(1-\delta)^{-1}}{l^2}+C_1'.$$
Since $l\ge (0.8)r$ and since $t-t'\le \delta r^2$, we see that
provided $\delta$ is sufficiently small (depending on the bound of
the curvature of the ambient flow) we have
$$Q(x_0,t)< \frac{3}{2(t-t')}$$ for all $t\in [t',t_2].$
Of course, since $|Z|^2\ge 1-\delta$ everywhere, this shows that
$$k^2(x_0,t)\le |Y(x_0,t)|^2=
(2-|Z(x_0,t)|^2)Q(x_0,t)< \frac{2}{(t-t')}$$ for all $t\in
[t',t_2]$. Since $x_0$ was an arbitrary point of $c(\cdot,t')$,
 this shows that
 $k(x,t)<\sqrt{\frac{2}{t-t'}}$ for all $x\in c(\cdot,t)$ and all $t\in
 [t',t_2]$. By the definition of $t_2$ this implies that $t_2=t'+\delta r^2$ and completes the
 proof of  Lemma~\ref{lem:curve-shrinking-shi-estimate-first-inequality}.
